%% ===================================================================================
%% VITRINE -- Documentação oficial do projeto (Trabalho de Conclusão de Curso)
%% ===================================================================================
%% Este arquivo é a RAIZ do documento. Ele é a "Fonte da Verdade" textual do projeto:
%% a especificação (requisitos, arquitetura, modelo de dados, API) vive aqui e nos
%% diagramas de doc/diagrams/src/. Quem edita: o agente Arquiteto / Redator Técnico.
%% Fluxo completo: ver CLAUDE.md, ARQUITETO.md e DESENVOLVEDORES.md na raiz do repo.
%%
%% Base: template abntex2 (undergrad_thesis.cls). Compilação:  latexmk -pdf main.tex
%%   (ou `make pdf` dentro de doc/tex). NÃO é necessário -shell-escape.
%% ===================================================================================

\documentclass[
	12pt,        % Tamanho da fonte
	openright,   % Capítulos começam em páginas ímpares
	oneside,     % Impressão em um lado (use 'twoside' para frente e verso)
	a4paper,     % Tamanho do papel
	english,     % Idioma adicional para hifenização
	brazil       % Idioma principal do documento (deve ser o último)
]{undergrad_thesis}

%% ===================================================================================
%% PACOTES
%% ===================================================================================
\usepackage[T1]{fontenc}      % Fontes T1 de 8 bits
\usepackage[utf8]{inputenc}   % UTF-8 (acentos e caracteres especiais)
\usepackage{lmodern}          % Latin Modern
\usepackage{indentfirst}      % Indenta o primeiro parágrafo de cada seção
\usepackage{cmap}             % PDF pesquisável/copiável

\usepackage{graphicx}         % Inclusão de imagens
\usepackage{booktabs}         % Tabelas profissionais
\usepackage{listings}         % Listagens de código (estilo e linguagens em vitrine.sty)
\usepackage[table]{xcolor}    % Cores, inclusive em tabelas
\usepackage{multirow}         % Células mescladas
\usepackage{amsmath, amsfonts, amssymb}
\usepackage{inconsolata}      % Fonte monoespaçada

\usepackage[num]{abntex2cite} % Citações ABNT numéricas (para autor-data: [alf])

% Macros do projeto: figuras com placeholder, pendências por agente, selos de status.
% Troque [rascunho] por [final] na versão de entrega: pendências somem e figura ausente vira erro.
\usepackage[rascunho]{vitrine}

%% ===================================================================================
%% INFORMAÇÕES DO DOCUMENTO (capa, folha de rosto e metadados do PDF)
%% ===================================================================================
% Título, orientador e banca: dados herdados do Relatório de IC; confirmar para o TCC.
\titulo{Vitrine: Um Sistema de Publicação de Conteúdo Científico}
\autor{Caio Augusto Silva dos Reis}
\local{São João del-Rei, Minas Gerais}
\data{2026}
\orientador{Prof. Dr. Matheus Carvalho Viana}
\instituicao{%
  Universidade Federal de São João del-Rei
  \par
  Curso de Graduação em Ciência da Computação}
\tipotrabalho{Trabalho de Conclusão de Curso}
\preambulo{Trabalho de Conclusão de Curso apresentado ao Curso de Graduação em Ciência da Computação da Universidade Federal de São João del-Rei como requisito parcial para a obtenção do grau de Bacharel em Ciência da Computação.}
\logo{imgs/university-logo.jpg.png}

% -----------------------------------------------------------------------------------
% Hiperlinks e metadados do PDF
% -----------------------------------------------------------------------------------
\definecolor{hyperlinkblue}{RGB}{41, 5, 195}

\makeatletter
\hypersetup{
    pdftitle={\@title},
    pdfauthor={\@author},
    pdfsubject={\imprimirpreambulo},
    pdfcreator={LaTeX com abnTeX2},
    pdfkeywords={CMS acadêmico; Next.js; SQLite; Markdown; Documentation-Driven Development},
    colorlinks=true,
    linkcolor=hyperlinkblue,
    citecolor=hyperlinkblue,
    filecolor=magenta,
    urlcolor=hyperlinkblue,
    bookmarksdepth=4,
    bookmarkstype=toc
}
\makeatother

% -----------------------------------------------------------------------------------
% Espaçamento e profundidade de numeração/sumário
% -----------------------------------------------------------------------------------
\setlength{\parindent}{1.5cm}
\setlength{\parskip}{0.2cm}
\setcounter{secnumdepth}{3}
\setcounter{tocdepth}{3}

\makeindex

%% ===================================================================================
%% INÍCIO DO DOCUMENTO
%% ===================================================================================
\begin{document}

\frenchspacing

% ----------------------------------------------------------
% ELEMENTOS PRÉ-TEXTUAIS
% ----------------------------------------------------------
\pretextual

\imprimircapalogo
\imprimirfolhaderosto

% Folha de aprovação, dedicatória, agradecimentos, epígrafe, resumo e abstract
%% ===================================================================================
%% Elementos Pré-Textuais (pretextual_elements.tex)
%% ===================================================================================
%% Contém: folha de aprovação, dedicatória, agradecimentos, epígrafe, resumo e abstract.
%% Campos entre colchetes são de preenchimento pessoal do autor (não são tarefa de agente).
%% O resumo e o abstract são RASCUNHOS e devem ser reescritos pelo Arquiteto/Redator
%% quando os resultados (Cap. 7 e 8) estiverem consolidados.
%% ===================================================================================

% -----------------------------------------------------------------------------------
% FOLHA DE APROVAÇÃO (obrigatória -- ABNT NBR 14724)
% Após a defesa, substitua o conteúdo do ambiente pela página assinada digitalizada:
%   \includegraphics[width=\textwidth]{imgs/folha_aprovacao_assinada.pdf}
% -----------------------------------------------------------------------------------
\begin{folhadeaprovacao}
  \begin{center}
    {\ABNTEXchapterfont\large\imprimirautor}

    \vspace*{\fill}
    {\ABNTEXchapterfont\bfseries\Large\imprimirtitulo}
    \vspace*{\fill}

    \hspace{.45\textwidth}
    \begin{minipage}{.5\textwidth}
      \SingleSpacing
      \imprimirpreambulo
    \end{minipage}%
    \vspace*{\fill}
  \end{center}

  Trabalho aprovado. \imprimirlocal, \underline{\hspace{1cm}} de \underline{\hspace{3cm}} de \underline{\hspace{1.5cm}}:
  \vspace*{2cm}

  \assinatura{\textbf{\imprimirorientador} \\ Orientador(a)}
  \assinatura{\textbf{[Nome completo do membro da banca 1]} \\ {[Instituição]}}
  \assinatura{\textbf{[Nome completo do membro da banca 2]} \\ {[Instituição]}}

\end{folhadeaprovacao}

% -----------------------------------------------------------------------------------
% DEDICATÓRIA (opcional)
% -----------------------------------------------------------------------------------
\begin{dedicatoria}
	\vspace*{\fill}
	\flushright
	\textit{[Dedicatória -- a critério do autor.]}
\end{dedicatoria}

% -----------------------------------------------------------------------------------
% AGRADECIMENTOS (opcional)
% -----------------------------------------------------------------------------------
\begin{agradecimentos}
[Agradecimentos -- a critério do autor. Sugere-se mencionar o orientador,
Prof.\ Dr.\ Matheus Carvalho Viana, pelo acompanhamento da Iniciação Científica que
originou este trabalho, a Universidade Federal de São João del-Rei e os colegas e
familiares que apoiaram o desenvolvimento.]
\end{agradecimentos}

% -----------------------------------------------------------------------------------
% EPÍGRAFE (opcional)
% -----------------------------------------------------------------------------------
\begin{epigrafe}
	\vspace*{\fill}
	\begin{flushright}
		\textit{[Epígrafe -- a critério do autor.]}
	\end{flushright}
\end{epigrafe}

% -----------------------------------------------------------------------------------
% RESUMO EM PORTUGUÊS (obrigatório) -- RASCUNHO
% -----------------------------------------------------------------------------------
\begin{resumo}
Este trabalho apresenta o Vitrine, uma plataforma de gerenciamento de conteúdo (CMS) de código aberto voltada à publicação científica e didática independente por docentes e pesquisadores. O problema que motiva o projeto é a barreira tecnológica enfrentada por autores que, ao buscarem autonomia digital, esbarram na complexidade da infraestrutura dos CMSs tradicionais ou na rigidez e na perda de controle impostas por plataformas proprietárias, que ainda oferecem suporte limitado a elementos da escrita científica, como fórmulas matemáticas e blocos de código. A pesquisa, de natureza aplicada e caráter exploratório-descritivo, adota o desenvolvimento experimental e incremental de um artefato tecnológico. O sistema é construído como um monólito modular sobre o \textit{framework} Next.js, com persistência em SQLite por meio do ORM Prisma, autenticação com a biblioteca better-auth, edição e renderização de conteúdo em Markdown com fórmulas em \LaTeX{} (KaTeX) e personalização visual em tempo de execução, por variáveis CSS persistidas em banco de dados e propagadas por \textit{Incremental Static Regeneration} (ISR). Como contribuição metodológica, o desenvolvimento segue a abordagem \textit{Documentation-Driven Development}, em que esta documentação, escrita em \LaTeX{} e acompanhada de diagramas versionados como código, constitui a fonte da verdade que orienta agentes especializados (dados, servidor e interface) sob a governança de um arquiteto integrador. O documento registra o levantamento de requisitos, o modelo de dados, a arquitetura, as decisões de projeto e as estratégias de implementação e de verificação do sistema. [Os resultados e as conclusões serão consolidados ao final da implementação.]

\vspace{\onelineskip}
\noindent
\textbf{Palavras-chave}: CMS acadêmico; Next.js; SQLite; Markdown; \textit{Documentation-Driven Development}; código aberto.
\end{resumo}

% -----------------------------------------------------------------------------------
% RESUMO EM LÍNGUA ESTRANGEIRA (obrigatório) -- RASCUNHO
% -----------------------------------------------------------------------------------
\begin{resumo}[Abstract]
\begin{otherlanguage*}{english}
This work presents Vitrine, an open-source content management system (CMS) aimed at independent scientific and educational publishing by lecturers and researchers. The project is motivated by the technological barrier faced by authors who, when seeking digital autonomy, run into the infrastructure complexity of traditional CMSs or into the rigidity and loss of control imposed by proprietary platforms, which still offer limited support for elements of scientific writing such as mathematical formulas and code blocks. The research is applied and exploratory-descriptive in nature and adopts experimental, incremental development of a technological artifact. The system is built as a modular monolith on the Next.js framework, with SQLite persistence through the Prisma ORM, authentication with the better-auth library, Markdown content editing and rendering with \LaTeX{} formulas (KaTeX), and runtime visual customization through CSS variables persisted in the database and propagated by Incremental Static Regeneration (ISR). As a methodological contribution, development follows the Documentation-Driven Development approach, in which this documentation, written in \LaTeX{} and accompanied by diagrams versioned as code, is the source of truth guiding specialized agents (data, server and interface) under the governance of an integrating architect. The document records requirements elicitation, the data model, the architecture, design decisions, and the implementation and verification strategies of the system. [Results and conclusions will be consolidated at the end of the implementation.]

\vspace{\onelineskip}
\noindent
\textbf{Keywords}: academic CMS; Next.js; SQLite; Markdown; Documentation-Driven Development; open source.
\end{otherlanguage*}
\end{resumo}


\pdfbookmark[0]{\listfigurename}{lof}
\listoffigures*
\cleardoublepage

\pdfbookmark[0]{\listtablename}{lot}
\listoftables*
\cleardoublepage

%% ===================================================================================
%% Siglas e Abreviaturas (acronyms.tex)
%% ===================================================================================
%% Mantida pelo Arquiteto/Redator. Ordem alfabética. Ao introduzir uma sigla nova em
%% qualquer capítulo, acrescente-a aqui no mesmo commit.
%% ===================================================================================

\begin{siglas}
	\item[ADR] \textit{Architecture Decision Record} (Registro de Decisão Arquitetural)
	\item[API] \textit{Application Programming Interface} (Interface de Programação de Aplicações)
	\item[BD] Banco de Dados
	\item[CMS] \textit{Content Management System} (Sistema de Gerenciamento de Conteúdo)
	\item[CRUD] \textit{Create, Read, Update, Delete} (Criar, Ler, Atualizar, Excluir)
	\item[CSR] \textit{Client-Side Rendering} (Renderização no Cliente)
	\item[CSRF] \textit{Cross-Site Request Forgery}
	\item[CSS] \textit{Cascading Style Sheets}
	\item[DBML] \textit{Database Markup Language}
	\item[DDD] \textit{Documentation-Driven Development} (Desenvolvimento Orientado por Documentação)
	\item[DER] Diagrama Entidade-Relacionamento
	\item[GFM] \textit{GitHub Flavored Markdown}
	\item[HTML] \textit{HyperText Markup Language}
	\item[HTTP] \textit{HyperText Transfer Protocol}
	\item[IC] Iniciação Científica
	\item[ISR] \textit{Incremental Static Regeneration} (Regeneração Estática Incremental)
	\item[JS] JavaScript
	\item[JSON] \textit{JavaScript Object Notation}
	\item[JSX] \textit{JavaScript XML}
	\item[LCP] \textit{Largest Contentful Paint}
	\item[ORM] \textit{Object-Relational Mapping} (Mapeamento Objeto-Relacional)
	\item[OWASP] \textit{Open Worldwide Application Security Project}
	\item[RF] Requisito Funcional
	\item[RN] Regra de Negócio
	\item[RNF] Requisito Não Funcional
	\item[SEO] \textit{Search Engine Optimization} (Otimização para Mecanismos de Busca)
	\item[SGBD] Sistema de Gerenciamento de Banco de Dados
	\item[SQL] \textit{Structured Query Language} (Linguagem de Consulta Estruturada)
	\item[SSG] \textit{Static Site Generation} (Geração Estática de Páginas)
	\item[SSR] \textit{Server-Side Rendering} (Renderização no Servidor)
	\item[SUS] \textit{System Usability Scale}
	\item[UC] Caso de Uso (\textit{Use Case})
	\item[UI] \textit{User Interface} (Interface de Usuário)
	\item[UML] \textit{Unified Modeling Language}
	\item[URL] \textit{Uniform Resource Locator}
	\item[UX] \textit{User Experience} (Experiência de Usuário)
	\item[VPS] \textit{Virtual Private Server} (Servidor Virtual Privado)
	\item[WAL] \textit{Write-Ahead Logging}
	\item[WCAG] \textit{Web Content Accessibility Guidelines}
\end{siglas}


\pdfbookmark[0]{\contentsname}{toc}
\tableofcontents*
\cleardoublepage

% ----------------------------------------------------------
% ELEMENTOS TEXTUAIS
% Dica: para compilar só um capítulo, use \includeonly{capitulos/cap5_arquitetura} aqui.
% ----------------------------------------------------------
\textual

%% ===================================================================================
%% Capítulo 1 -- Introdução
%% Documento de Arquitetura e Requisitos de Software do Vitrine. Origem do produto:
%% Relatório de IC (doc/ref). Dono: Arquiteto/Redator.
%% ===================================================================================
\chapter{Introdução}\label{cap:intro}

\section{Propósito e Público do Documento}\label{sec:proposito}

Este documento é a especificação oficial do Vitrine: reúne os requisitos, a arquitetura, o modelo de dados, os contratos de API, o plano de implementação e o plano de verificação do sistema. É a \textbf{Fonte da Verdade} do projeto: o código a implementa e o Arquiteto confere a aderência a cada integração (Seção~\ref{sec:processo-ddd}). Ele é mantido junto ao código, em \LaTeX{}, e evolui com o produto.

O público do documento é o Engenheiro Chefe, que decide o rumo do produto; os agentes e desenvolvedores dos domínios Dados, Core e UI, que implementam; o Implantador, que instala e opera o sistema (Apêndice~\ref{ap:implantacao}); e quem quiser estender ou reimplementar o Vitrine.

\section{Visão Geral do Produto}\label{sec:visao-geral-produto}

A disseminação do conhecimento acadêmico encontrou na web o seu principal canal de expansão \cite{ginige2001web}. Docentes e pesquisadores, especialmente os de áreas não tecnológicas, dependem cada vez mais de um espaço próprio para publicar artigos, notas de aula e materiais didáticos, com identidade reconhecível e endereço estável. Entretanto, a criação desse espaço ainda esbarra em um dilema: de um lado, a dependência de plataformas proprietárias e rígidas, que limitam o controle sobre a identidade visual, a formatação de elementos técnicos --- fórmulas matemáticas e blocos de código --- e a perenidade do conteúdo; de outro, a complexidade de instalar, configurar e manter sistemas de gerenciamento de conteúdo (\en{Content Management Systems}, CMS) tradicionais \cite{mikkonen2010mashware}.

O Vitrine é a resposta a esse problema: uma plataforma de gerenciamento de conteúdo de código aberto, sob licença MIT (\adr{ADR-015}), concebida para a criação de blogs e portfólios acadêmicos. Combina um editor Markdown com suporte a fórmulas em \LaTeX{}, uma estrutura hierárquica de documentos e coleções, um estado de publicação por conteúdo e um sistema de personalização visual aplicado em tempo de execução. A primeira versão foi concebida e prototipada em um projeto de Iniciação Científica (IC) entre setembro de 2025 e fevereiro de 2026 \cite{reis2026ic}; esta documentação consolida essa base em uma arquitetura monolítica modular sobre Next.js, SQLite e better-auth.

A premissa de produto é que cada educador deve poder possuir a sua própria ``vitrine'' de conhecimento, de forma independente e personalizada: a implantação deve ser um ato acessível, e a barreira técnica não deve impedir a presença digital qualificada do docente. Essa premissa dialoga com a visão de que sistemas de informação eficazes devem ser facilitadores dos processos humanos e não obstáculos a eles \cite{rainer2019intro}.

\section{Objetivos do Produto}\label{sec:objetivos}

\subsection{Objetivo Geral}

Oferecer uma plataforma CMS de código aberto, modular e de fácil implantação, baseada no \en{framework} Next.js, voltada à criação de blogs e portfólios acadêmicos por professores e pesquisadores, descrita por uma documentação que seja a fonte da verdade que orienta e verifica a sua implementação.

\subsection{Objetivos Específicos}\label{sec:objetivos-especificos}

Os objetivos OE1 a OE5 são herdados do relatório de IC \cite{reis2026ic}; o único ajuste é a adoção da biblioteca better-auth no lugar do Auth.js inicialmente previsto (decisão \adr{ADR-004}).

\begin{enumerate}[label=\textbf{OE\arabic*.}, leftmargin=1.6cm]
  \item Viabilizar a publicação independente por meio de um processo de clonagem e implantação simplificado;
  \item Implementar um editor de texto baseado em Markdown que atenda às necessidades da escrita científica, como realce de sintaxe e fórmulas em \LaTeX{};
  \item Prover um sistema de personalização visual dinâmica que permita alterar a paleta de cores e a identidade do site por meio de interface administrativa, sem novos ciclos de compilação (\en{build});
  \item Adotar uma arquitetura de dados portátil, com o ORM Prisma e o banco de dados SQLite;
  \item Integrar um sistema de autenticação simplificado, por meio do better-auth, para garantir a segurança no gerenciamento de conteúdos.
\end{enumerate}

\section{Justificativa Técnica}\label{sec:justificativa}

O equilíbrio entre modernidade tecnológica e portabilidade orienta as escolhas. O Next.js permite uma renderização híbrida de alto desempenho, essencial para a indexação por motores de busca (\en{Search Engine Optimization}, SEO) \cite{nextjsdocs}. O SQLite, por ser um sistema baseado em arquivo, simplifica a infraestrutura e reduz custos para o usuário final \cite{sqliteabout,sqlitewhentouse}. A camada de acesso a dados com Prisma ORM isola o restante do sistema do dialeto SQL e oferece um modelo declarativo versionável \cite{prismadocs}. As customizações visuais são feitas por variáveis CSS em tempo de execução, com persistência em banco de dados e reflexo imediato por meio de \en{Incremental Static Regeneration} (ISR) \cite{nextjsdocs}.

Embora o Vitrine seja intuitivo para usuários leigos, ele também atende a desenvolvedores, ao oferecer uma base sólida que elimina o trabalho repetitivo de configurar estruturas do zero. O valor do produto pode ser resumido em três dimensões:

\begin{itemize}
  \item \textbf{Social:} contribui para a democratização da publicação acadêmica ao reduzir a dependência de plataformas proprietárias;
  \item \textbf{Prática:} entrega ao docente uma ferramenta pronta para uso, com suporte a fórmulas, código e identidade visual própria;
  \item \textbf{Tecnológica:} demonstra uma arquitetura monolítica enxuta, sem serviços externos obrigatórios, que combina desempenho estático e flexibilidade dinâmica.
\end{itemize}

\section{Escopo e Delimitações}\label{sec:escopo}

O escopo desta versão do Vitrine compreende: autenticação de um administrador; gestão hierárquica de documentos e coleções, com estado de publicação (rascunho ou publicado); edição e renderização de Markdown com fórmulas \LaTeX{}, tabelas e código; navegação por barra superior e \en{sidebar} hierárquica; personalização de identidade, paleta de cores e links sociais do rodapé; e entrega das páginas públicas por ISR com invalidação sob demanda, implantadas por imagem Docker em uma VPS.

Ficam fora do escopo desta versão, embora documentados como possibilidades de evolução (Seção~\ref{sec:trabalhos-futuros}): múltiplos usuários com papéis e permissões distintos; comentários e interação de leitores; busca textual; histórico de versões de conteúdo; envio e gerenciamento de mídia; internacionalização; e edição colaborativa em tempo real.

\section{Histórico do Projeto: da Iniciação Científica à Versão Atual}\label{sec:historico}

O relatório de IC \cite{reis2026ic} registra a concepção do sistema, os requisitos iniciais, os diagramas de casos de uso e de entidade-relacionamento, os protótipos de baixa fidelidade e a decisão de usar Next.js, Prisma e SQLite. Esta documentação reaproveita integralmente esse acervo e o evolui. A Tabela~\ref{tab:ic-vs-atual} resume o que permanece e o que muda.

\begin{table}[htb]
\caption{Comparação entre o relatório de IC e a versão atual}\label{tab:ic-vs-atual}
\small
\begin{tabularx}{\linewidth}{@{}L{2.9cm}YY@{}}
\toprule
\textbf{Aspecto} & \textbf{Relatório de IC (set/2025 a fev/2026)} & \textbf{Versão atual}\\
\midrule
Objetivo & Conceber e prototipar o CMS acadêmico & Consolidar, documentar e verificar o monólito\\
\addlinespace
Autenticação & Auth.js (previsto) & better-auth (\adr{ADR-004})\\
\addlinespace
Modelo de dados & Três entidades: usuários, conteúdos e configurações & Seis modelos: quatro do better-auth (\cod{user}, \cod{session}, \cod{account}, \cod{verification}) mais conteúdos e configurações\\
\addlinespace
Especificação & Relatório em PDF, diagramas exportados & Documentação \LaTeX{} viva e diagramas como código (PlantUML e DBML)\\
\addlinespace
Processo & Incremental, individual & \en{Documentation-Driven Development} com agentes especializados e mural\\
\addlinespace
Interface & Protótipos de baixa fidelidade e uma instância de demonstração & Implementação completa conforme a especificação (Capítulo~\ref{cap:implementacao})\\
\addlinespace
Cache & ISR citado como requisito & ISR com revalidação sob demanda por rota e por \en{layout} (Seção~\ref{sec:arq-isr})\\
\addlinespace
Implantação & Clonagem do repositório e implantação em plataforma como a Vercel (tutoriais previstos) & Imagem Docker em VPS tradicional, com volume persistente para o SQLite (\adr{ADR-011})\\
\addlinespace
Conteúdo & Sem estado de publicação no modelo de dados (o protótipo oferece ``Salvar localmente'' e ``Enviar nova versão'') & Estado de publicação \texttt{DRAFT} ou \texttt{PUBLISHED}, com padrão \texttt{PUBLISHED} (\adr{ADR-012})\\
\bottomrule
\end{tabularx}
\fonte{Elaborado pelos autores, com base em \citeonline{reis2026ic} e no estado do repositório.}
\end{table}

Uma consequência dessa evolução é que alguns termos do relatório de IC --- por exemplo, ``Auth.js'' e a tabela ``\cod{users}'' --- foram substituídos nesta documentação. Sempre que o texto fizer referência ao relatório de IC, a origem será explicitada.

\section{Convenções do Documento}\label{sec:convencoes}

Os elementos do projeto têm identificadores estáveis, citados nos \en{commits}, nos pedidos de integração e nos testes: \texttt{UC} (casos de uso), \texttt{RF} e \texttt{RNF} (requisitos), \texttt{RN} (regras de negócio), \texttt{ADR} (decisões de projeto), \texttt{M} (módulos) e \texttt{CT} (casos de teste); as referências cruzadas são hiperligações. Trechos marcados como \emph{normativos} são contratos a cumprir; os marcados como \emph{esboço} apenas ilustram e não vinculam a implementação. Blocos tracejados amarelos indicam figuras ou trechos de código ainda não produzidos e são listados no Apêndice~\ref{ap:pendencias}.

\section{Organização do Documento}\label{sec:organizacao}

Além desta introdução, o documento está organizado da seguinte forma:

\begin{itemize}
  \item \textbf{Capítulo~\ref{cap:fundamentacao}:} apresenta o contexto tecnológico --- CMS, arquiteturas e estratégias de renderização, tecnologias adotadas, engenharia de requisitos e \en{Documentation-Driven Development};
  \item \textbf{Capítulo~\ref{cap:metodologia}:} descreve a abordagem de desenvolvimento, o planejamento por fases, o processo com agentes, o ambiente e os riscos;
  \item \textbf{Capítulo~\ref{cap:requisitos}:} apresenta o levantamento e a análise de requisitos: atores, casos de uso, requisitos funcionais e não funcionais, regras de negócio e prototipagem;
  \item \textbf{Capítulo~\ref{cap:arquitetura}:} detalha o projeto da arquitetura: visões, decisões, modelo de dados, autenticação, renderização e cache, contrato de API, segurança e implantação;
  \item \textbf{Capítulo~\ref{cap:implementacao}:} descreve a implementação por módulo, com o estado de cada um;
  \item \textbf{Capítulo~\ref{cap:testes}:} define a estratégia e o plano de verificação e validação;
  \item \textbf{Capítulo~\ref{cap:resultados}:} registra a situação da verificação, as limitações conhecidas e o roteiro de evolução.
\end{itemize}

Os apêndices reúnem o dicionário de dados, o esquema Prisma, os fontes dos diagramas, o guia de implantação, os registros de decisão de projeto (ADR), o mapa de pendências da documentação e o glossário.

%% ===================================================================================
%% Capítulo 2 -- Contexto Tecnológico
%% Fontes: Relatório de IC (Cap. 2), documentação oficial das tecnologias adotadas
%% (conferida nas cópias locais em node_modules e nos sítios citados em references.bib).
%% Dono: Arquiteto/Redator.
%% ===================================================================================
\chapter{Contexto Tecnológico}\label{cap:fundamentacao}

Este capítulo reúne os conceitos necessários para compreender as decisões do projeto. Inicia pelo domínio do problema, a publicação científica na web (Seção~\ref{sec:fund-publicacao}), situa os sistemas de gerenciamento de conteúdo (Seção~\ref{sec:fund-cms}) e as estratégias de renderização (Seção~\ref{sec:fund-renderizacao}), detalha as tecnologias adotadas (Seção~\ref{sec:fund-tecnologias}), apresenta os fundamentos de engenharia de requisitos e modelagem (Seção~\ref{sec:fund-modelagem}) e de \en{Documentation-Driven Development} (Seção~\ref{sec:fund-ddd}) e encerra com as soluções existentes e o posicionamento do produto (Seção~\ref{sec:trabalhos-relacionados}).

\section{Publicação Científica e Didática na Web}\label{sec:fund-publicacao}

\subsection{A Web como Canal de Disseminação}

A engenharia web consolidou-se como a aplicação de princípios de engenharia ao desenvolvimento, à implantação e à manutenção de sistemas baseados na web \cite{ginige2001web}. Ao longo do tempo, o que começou como um conjunto de documentos interligados evoluiu para aplicações completas, e a web tornou-se o meio natural de circulação de artigos, materiais didáticos e dados de pesquisa. Essa evolução, porém, nem sempre foi acompanhada de práticas de engenharia: a facilidade de combinar tecnologias heterogêneas --- o fenômeno que \citeonline{mikkonen2010mashware} descrevem como \en{mashware} --- produz sistemas que funcionam, mas cuja manutenção e cujo entendimento dependem de conhecimento tácito. Esse diagnóstico reforça, para o Vitrine, a necessidade de especificação explícita e de rastreabilidade entre requisitos, projeto e código (Capítulo~\ref{cap:metodologia}).

\subsection{Autonomia Digital de Docentes e Pesquisadores}

Para o autor acadêmico, publicar na web envolve três necessidades: \textbf{controle} (sobre o conteúdo, a identidade visual e o endereço), \textbf{perenidade} (o material continua acessível independentemente de mudanças de serviço) e \textbf{adequação técnica} (capacidade de expressar a sua área de conhecimento). As alternativas disponíveis atendem a essas necessidades de forma desigual. Plataformas de terceiros oferecem facilidade, mas impõem limites de formatação e de identidade; CMS tradicionais oferecem liberdade, mas exigem competências de infraestrutura que muitos docentes não têm nem desejam adquirir \cite{mikkonen2010mashware,reis2026ic}.

\subsection{Elementos da Escrita Científica}

A escrita científica e didática usa recursos que raramente são prioridade em sistemas genéricos de publicação: fórmulas matemáticas (inline e em bloco), trechos de código com realce de sintaxe, tabelas, figuras numeradas e referências. No ecossistema acadêmico, a notação \LaTeX{} é amplamente adotada para fórmulas; por isso, um sistema voltado a docentes e pesquisadores deve aceitá-la diretamente no texto, sem exigir imagens ou ferramentas externas. Esse requisito orienta a escolha de Markdown com suporte a matemática (Seção~\ref{sec:fund-markdown}).

\section{Sistemas de Gerenciamento de Conteúdo}\label{sec:fund-cms}

\subsection{Conceito e Componentes}

Um sistema de gerenciamento de conteúdo (CMS) é um \en{software} que permite criar, organizar, editar e publicar conteúdo digital separando o conteúdo de sua apresentação, em geral por meio de uma interface de edição acessível a pessoas sem formação em programação. Os componentes típicos são: um repositório de conteúdo (banco de dados ou arquivos); uma interface de edição; um mecanismo de apresentação, que transforma o conteúdo em páginas; mecanismos de administração, como usuários, permissões e configurações; e uma política de publicação.

\subsection{Taxonomia}\label{sec:fund-taxonomia}

Na prática, as soluções disponíveis agrupam-se em quatro categorias (Tabela~\ref{tab:taxonomia-cms}). A classificação é didática: produtos reais podem combinar características de mais de uma categoria.

\begin{table}[htb]
\caption{Categorias de soluções para publicação de conteúdo}\label{tab:taxonomia-cms}
\small
\begin{tabularx}{\linewidth}{@{}L{3.1cm}YL{4.6cm}@{}}
\toprule
\textbf{Categoria} & \textbf{Características} & \textbf{Exemplos usuais}\\
\midrule
Plataformas hospedadas & Hospedagem e banco de dados gerenciados por terceiros; edição pela web; pouca liberdade de identidade visual e de formatação técnica & Redes de publicação e blogs hospedados\\
\addlinespace
CMS tradicionais (acoplados) & Edição, regras e apresentação no mesmo \en{software}; exigem servidor, linguagem de \en{script} e banco relacional; grande ecossistema de temas e \en{plugins} & WordPress, Drupal, Joomla\\
\addlinespace
CMS desacoplados (\en{headless}) & O CMS expõe o conteúdo por API; a apresentação é construída à parte, em qualquer tecnologia & Strapi, Contentful, Sanity\\
\addlinespace
Geradores de sites estáticos e documentação como código & Conteúdo em arquivos versionados no Git; um processo de \en{build} gera HTML estático; normalmente sem interface web de edição & Jekyll, Hugo, MkDocs, Docusaurus\\
\bottomrule
\end{tabularx}
\fonte{Elaborado pelos autores.}
\end{table}

\subsection{Posicionamento do Vitrine}

O Vitrine reúne características de duas categorias. Do ponto de vista da implantação, é um CMS acoplado --- uma única aplicação com interface de edição, regras e apresentação, sem serviços externos --- o que simplifica a instalação. Do ponto de vista da entrega, aproxima-se dos geradores estáticos: as páginas públicas são pré-renderizadas e servidas de cache, e regeneradas de forma incremental quando o conteúdo muda (Seção~\ref{sec:fund-isr}). O resultado pretendido é o equilíbrio entre autonomia e simplicidade que motivou o projeto.

\section{Arquiteturas e Estratégias de Renderização}\label{sec:fund-renderizacao}

\subsection{CSR, SSR, SSG e ISR}\label{sec:fund-isr}

A forma como o HTML de uma página é produzido determina desempenho, custo e facilidade de indexação. Em aplicações de página única (\en{Single Page Applications}), o navegador recebe um documento quase vazio e monta a interface com JavaScript, estratégia chamada \en{Client-Side Rendering} (CSR) \cite{jadhav2015single}. Em \en{Server-Side Rendering} (SSR), o servidor gera o HTML a cada requisição. Em \en{Static Site Generation} (SSG), o HTML é gerado uma vez, no \en{build}. A \en{Incremental Static Regeneration} (ISR) combina as duas últimas: as páginas são servidas como estáticas e regeneradas depois de um intervalo ou sob demanda, sem reconstruir o site inteiro \cite{nextjsdocs}. A Tabela~\ref{tab:renderizacao} compara as estratégias.

\begin{table}[htb]
\caption{Estratégias de renderização e adequação ao Vitrine}\label{tab:renderizacao}
\small
\begin{tabularx}{\linewidth}{@{}L{1.6cm}L{3.3cm}L{3.0cm}YL{2.4cm}@{}}
\toprule
\textbf{Sigla} & \textbf{Onde o HTML é gerado} & \textbf{Frescor do conteúdo} & \textbf{Observações} & \textbf{Uso no Vitrine}\\
\midrule
CSR & No navegador & Sempre atual & Primeira exibição mais lenta; indexação dependente de execução de JS & Somente em componentes interativos\\
\addlinespace
SSR & No servidor, a cada requisição & Sempre atual & Custo de CPU por requisição; latência maior & Painel administrativo\\
\addlinespace
SSG & No \en{build} & Atualizado só em novo \en{build} & Máximo desempenho; conteúdo novo exige reconstruir & Não isoladamente\\
\addlinespace
ISR & No \en{build} e regenerado depois & Por tempo ou sob demanda & Desempenho de estático com atualização incremental & Páginas públicas\\
\bottomrule
\end{tabularx}
\fonte{Elaborado pelos autores, com base em \citeonline{nextjsdocs}.}
\end{table}

\subsection{Jamstack}

O termo Jamstack (JavaScript, APIs e \en{Markup}) designa uma abordagem em que páginas são pré-renderizadas e servidas por rede de distribuição, enquanto funcionalidades dinâmicas são obtidas por APIs \cite{biilmann2019jamstack}. O relatório de IC alinhou o Vitrine a esse paradigma pela dissociação entre camada de apresentação e dados, com ganhos de segurança e escalabilidade \cite{reis2026ic}. No monólito, a camada de dados reside no mesmo processo, de modo que o Vitrine adota da Jamstack a \textbf{pré-renderização} e a \textbf{invalidação seletiva}, mas não a separação física entre os serviços.

\subsection{Monólito Modular versus Microsserviços}\label{sec:fund-monolito}

Um monólito é uma unidade de implantação única. Um monólito \emph{modular} mantém fronteiras internas claras --- por domínio, com dependências controladas --- de modo que a estrutura possa evoluir sem que cada módulo precise ser um serviço independente. Microsserviços trazem autonomia de equipes e escalabilidade seletiva, mas impõem custos de operação, consistência distribuída e observabilidade \cite{newman2019monolith}. O argumento conhecido como \en{MonolithFirst} recomenda começar por um monólito bem estruturado e extrair serviços apenas quando as fronteiras forem bem compreendidas e houver necessidade real \cite{fowler2015monolithfirst}. Para um sistema de um único \en{tenant}, de tráfego moderado e implantação pelo próprio docente, o monólito modular é a escolha que minimiza a complexidade acidental (decisão \adr{ADR-001}).

\subsection{Estilo REST}

O estilo arquitetural REST organiza a comunicação em torno de recursos identificados por URLs, manipulados por um conjunto uniforme de operações (os métodos HTTP), com mensagens autodescritivas e interação sem estado entre requisições \cite{fielding2000rest}. As rotas JSON do Vitrine (Seção~\ref{sec:arq-api}) seguem um REST pragmático: recursos \cod{/api/contents} e \cod{/api/configuration}, verbos HTTP com semântica usual, códigos de status e formato de erro padronizados. O estado de autenticação trafega em \en{cookie} de sessão, o que é um compromisso consciente em relação ao ideal de ausência total de estado.

\section{Tecnologias do Vitrine}\label{sec:fund-tecnologias}

\subsection{React e Componentização}

React é uma biblioteca para construção de interfaces a partir de componentes: funções que recebem dados (\en{props}) e descrevem a interface correspondente, com estado local e ciclo de vida expressos por \en{hooks} \cite{reactdocs,fedosejev2015react}. A versão 19 consolidou os \emph{Server Components}, componentes executados apenas no servidor, que podem acessar dados diretamente e não enviam o seu código ao navegador. O projeto também habilita o React Compiler, que aplica memoização automática em tempo de compilação.

\subsection{Next.js}\label{sec:fund-nextjs}

Next.js é o \en{framework} sobre React que fornece roteamento, renderização híbrida e otimizações de produção \cite{nextjsdocs,konshin2018next}. Os elementos relevantes para o Vitrine são:

\begin{itemize}
  \item \textbf{App Router:} o roteamento é definido pela estrutura de diretórios em \cod{src/app}; arquivos \cod{page.js}, \cod{layout.js} e \cod{route.js} declaram páginas, \en{layouts} compartilhados e \en{endpoints} HTTP (\en{Route Handlers});
  \item \textbf{Server e Client Components:} por padrão os componentes do App Router são de servidor; a diretiva \texttt{'use client'} marca os que dependem de estado ou eventos do navegador;
  \item \textbf{Proxy:} a partir da versão 16, o antigo \en{Middleware} passou a se chamar \en{Proxy} (arquivo \cod{proxy.js}); serve para redirecionamentos e verificações otimistas, mas a documentação oficial adverte que não deve ser a única barreira de autorização;
  \item \textbf{ISR e revalidação:} o segmento de rota pode declarar \texttt{export const revalidate = N} (regeneração por tempo) e usar \arq{generateStaticParams} para pré-renderizar parâmetros dinâmicos; as funções \arq{revalidatePath} e \texttt{revalidateTag}, chamadas em Route Handlers ou Server Functions, invalidam o cache sob demanda.
\end{itemize}

A documentação da versão instalada registra ressalvas importantes para este projeto: a ISR só é suportada no \en{runtime} Node.js e não em exportação estática; com múltiplas instâncias, o cache padrão em sistema de arquivos é por instância, de modo que a invalidação só atinge a instância que recebe a chamada; e o cabeçalho \arq{x-nextjs-cache} informa se a resposta foi \texttt{HIT}, \texttt{STALE}, \texttt{MISS} ou \texttt{REVALIDATED} \cite{nextjsdocs}. Essas propriedades fundamentam a estratégia de cache (Seção~\ref{sec:arq-isr}) e os testes de ISR (Capítulo~\ref{cap:testes}).

\subsection{Persistência: SQLite e Prisma ORM}\label{sec:fund-sqlite}

O SQLite é um mecanismo de banco de dados SQL embutido, sem servidor e sem configuração, que armazena o banco inteiro em um único arquivo e oferece transações com propriedades ACID \cite{sqliteabout}. Sua documentação indica que ele é adequado à maioria dos sites de tráfego baixo a moderado \cite{sqlitewhentouse}, perfil do Vitrine. As limitações relevantes são: a escrita é serializada --- há um escritor por vez --- e a durabilidade depende de um sistema de arquivos persistente. O modo \en{Write-Ahead Logging} (WAL) permite que leitores e escritor atuem simultaneamente, o que é conveniente para um site majoritariamente de leitura \cite{sqlitewal}.

O Prisma ORM descreve o modelo de dados em um arquivo declarativo (\cod{schema.prisma}), gera um cliente tipado e gerencia migrações \cite{prismadocs}. Na versão 7, instalada no projeto, a configuração da CLI passa a residir em um arquivo \cod{prisma.config.ts}, o gerador recomendado é o \arq{prisma-client} e o acesso ao banco exige um \emph{\en{driver adapter}} (por exemplo, \arq{better-sqlite3} ou \texttt{libSQL} para SQLite). Um detalhe a considerar: o Prisma aceita campos do tipo \en{enum}, mas o SQLite não impõe os valores no nível do banco, o que justifica validar o tipo de conteúdo na aplicação (regra \texttt{RN02}, Seção~\ref{sec:regras-negocio}).

\subsection{Autenticação e Sessões}\label{sec:fund-auth}

\emph{Autenticação} é o processo de verificar quem é o usuário; \emph{autorização} é decidir o que ele pode fazer. Em aplicações web com senha, as boas práticas recomendadas pelo OWASP incluem armazenar apenas um \en{hash} da senha, gerado por função adaptativa e com sal, limitar tentativas de acesso e transportar identificadores de sessão em \en{cookies} com os atributos \texttt{HttpOnly}, \texttt{Secure} e \texttt{SameSite} \cite{owaspcheatsheets,owaspasvs}. Sessões persistidas no servidor --- e não apenas \en{tokens} autossuficientes --- permitem revogação imediata.

O better-auth é uma biblioteca de autenticação para aplicações JavaScript/TypeScript, independente de \en{framework}, que fornece autenticação por e-mail e senha, gerenciamento de sessões e um sistema de \emph{plugins} \cite{betterauthdocs}. Na versão instalada (1.7.6), o pacote expõe adaptador para Prisma (\cod{better-auth/adapters/prisma}), integração com Next.js (\cod{better-auth/next-js}, com o tratador de rotas e o \en{plugin} de \en{cookies}), cliente React (\cod{better-auth/react}), utilitário de leitura do \en{cookie} de sessão (\cod{better-auth/cookies}) e \en{plugins} como \texttt{admin}, \texttt{username} e \arq{two-factor}, que sinalizam caminhos de evolução para papéis e autenticação reforçada. O esquema de dados que ele requer --- tabelas \texttt{user}, \texttt{session}, \texttt{account} e \texttt{verification} --- é o que compõe a primeira metade do modelo de dados do Vitrine (Seção~\ref{sec:arq-dados}).

\subsection{Markdown, Matemática e o Ecossistema unified}\label{sec:fund-markdown}

Markdown é uma linguagem de marcação leve, legível em texto puro; a sua sintaxe básica é especificada pelo CommonMark \cite{commonmark}, e o GitHub Flavored Markdown (GFM) acrescenta tabelas, listas de tarefas, texto tachado e links automáticos \cite{gfmspec}. O ecossistema \texttt{unified} processa conteúdo como árvores sintáticas: \emph{remark} trata Markdown e \emph{rehype} trata HTML \cite{unifiedjs}. O \arq{react-markdown} usa esse ecossistema para converter Markdown em elementos React, sem empregar \arq{dangerouslySetInnerHTML} e sem interpretar HTML bruto por padrão --- uma propriedade importante de segurança \cite{reactmarkdown}.

O suporte a matemática é obtido pela combinação de \arq{remark-math}, que reconhece as delimitações \texttt{\$...\$} (inline) e \texttt{\$\$...\$\$} (bloco), e \arq{rehype-katex}, que renderiza as expressões com o KaTeX \cite{katex}. O KaTeX exige a folha de estilos correspondente. Com essa cadeia, o autor escreve fórmulas em \LaTeX{} diretamente no texto, cumprindo o objetivo OE2.

\subsection{Tailwind CSS e Variáveis CSS}\label{sec:fund-css}

O Tailwind CSS é um \en{framework} de utilitários; a versão 4 adota configuração orientada a CSS (\texttt{@import "tailwindcss"} e o bloco \texttt{@theme}) \cite{tailwinddocs}. As propriedades personalizadas de CSS (\emph{custom properties}) permitem definir valores reutilizáveis que participam da cascata e da herança e podem ser alterados em tempo de execução \cite{mdncssvars}. A combinação dos dois é o mecanismo de personalização do Vitrine: os tokens de cor do tema referenciam variáveis CSS, e um atributo \texttt{style} no elemento raiz define os valores vindos do banco, sem novo \en{build} (Seção~\ref{sec:arq-tema}).

\section{Engenharia de Requisitos e Modelagem}\label{sec:fund-modelagem}

\subsection{Requisitos}

Requisitos funcionais descrevem o que o sistema deve fazer; requisitos não funcionais descrevem restrições e qualidades sob as quais ele deve operar \cite{sommerville2011,pressman2021}. A norma ISO/IEC/IEEE 29148 estabelece processos e atributos de boa especificação --- identificação única, verificabilidade, consistência e rastreabilidade \cite{iso29148}. As qualidades do produto são classificadas pelo modelo da ISO/IEC 25010 em adequação funcional, eficiência de desempenho, compatibilidade, usabilidade, confiabilidade, segurança, manutenibilidade e portabilidade \cite{iso25010}; os requisitos não funcionais do Vitrine (Seção~\ref{sec:rnf}) usam essa taxonomia.

\subsection{UML, Modelo Entidade-Relacionamento e Diagramas como Código}

Os diagramas de casos de uso da UML descrevem as interações entre atores e sistema, com as relações \texttt{include} (obrigatória) e \texttt{extend} (opcional) \cite{umlspec}. O modelo entidade-relacionamento descreve dados e associações. Ambos são mantidos neste projeto como texto: PlantUML para diagramas UML e de arquitetura \cite{plantumldocs} e DBML para o modelo de dados \cite{dbmldocs}. Essa escolha permite versionar, comparar e revisar diagramas como se fossem código. A arquitetura é descrita pelas visões do modelo C4 --- contexto, contêineres e componentes \cite{c4model} --- complementadas por vistas de processo e de implantação no espírito do modelo 4+1 \cite{kruchten1995}. As decisões relevantes são registradas como \en{Architecture Decision Records} (ADR), documentos curtos que capturam contexto, decisão e consequências \cite{nygard2011adr}.

\section{Documentation-Driven Development}\label{sec:fund-ddd}

\subsection{Conceito}

Neste projeto, \emph{Documentation-Driven Development} (DDD) designa o processo em que a documentação é escrita \emph{antes} do código, passa a ser o contrato que o descreve e é atualizada em sincronia com ele; o código é verificado contra a documentação, e não o contrário. Não se deve confundir a sigla com \emph{Domain-Driven Design}, também abreviado como DDD, que trata da modelagem de domínios de negócio \cite{evans2003ddd} e cujos conceitos não são empregados aqui.

Essa prática tem antecedentes. \citeonline{parnas1986rational} argumentam que, mesmo que o projeto real nunca siga um processo racional ideal, produzir a documentação como se ele tivesse sido seguido traz benefícios à compreensão e à manutenção. O \en{Readme-Driven Development} propõe escrever primeiro o documento que explica o que o \en{software} fará \cite{rdd}. A abordagem \en{docs as code} trata a documentação com as mesmas ferramentas do código --- texto puro, controle de versão, revisão e automação \cite{writethedocs,gentle2017docs}.

\subsection{Desenvolvimento Conduzido por Agentes Especializados}

Quando parte do desenvolvimento é delegada a agentes de inteligência artificial, surgem riscos específicos: cada agente enxerga apenas uma parte do sistema, pode assumir premissas incorretas e produz alterações em volume que dificulta a revisão humana. O DDD mitiga esses riscos de três maneiras: (i) a especificação compartilhada reduz a necessidade de premissas; (ii) a separação de domínios e o \emph{escritor único} da documentação evitam conflitos; e (iii) um agente integrador (Arquiteto) confronta cada entrega com a especificação antes de aceitá-la. O Capítulo~\ref{cap:metodologia} detalha o processo adotado; o fluxo de ramificações segue práticas consagradas de controle de versão \cite{chacon2014progit} e a convenção de mensagens \en{Conventional Commits} \cite{conventionalcommits}.

\section{Soluções Existentes e Posicionamento}\label{sec:trabalhos-relacionados}

\subsection{Critérios de Comparação}

A comparação é feita por categoria de solução, segundo critérios derivados dos objetivos do produto: (C1) possibilidade de hospedagem própria; (C2) necessidade de administrar um servidor de banco de dados; (C3) edição de conteúdo por interface web; (C4) suporte nativo a fórmulas \LaTeX{} e a código com realce; e (C5) personalização visual sem programação.

\subsection{Limitações das Plataformas Existentes}

Redes sociais e plataformas de terceiros impõem três classes de limitação, que o relatório de IC \cite{reis2026ic} já identificava: a \textbf{perda de controle}, pela dificuldade de manter a identidade visual e a formatação técnica sob regras e limites de caracteres da plataforma, que ainda submetem a produção acadêmica à competição visual com conteúdos de outras naturezas; as \textbf{barreiras técnicas}, já que os CMS tradicionais exigem conhecimentos de infraestrutura para instalação e manutenção; e a \textbf{formatação técnica}, pela falta de suporte nativo a elementos essenciais, como fórmulas em \LaTeX{} e blocos de código.

Os geradores de sites estáticos e as ferramentas de documentação como código resolvem bem a formatação técnica --- tipicamente por extensões de matemática e de realce de sintaxe --- e a perenidade, já que o conteúdo são arquivos. O custo é o fluxo de trabalho: editar arquivos localmente, executar um \en{build} e publicar, o que pressupõe familiaridade com linha de comando e Git. Os CMS \en{headless} entregam um painel de edição, mas transferem ao autor a tarefa de construir e manter a camada de apresentação.

\subsection{Síntese e Posicionamento}

A Tabela~\ref{tab:trabalhos-relacionados} sintetiza a comparação. Os valores são qualitativos e descrevem tendências das categorias, não produtos específicos.

\begin{table}[htb]
\caption{Comparação qualitativa entre categorias de solução e o Vitrine}\label{tab:trabalhos-relacionados}
\scriptsize
\begin{tabularx}{\linewidth}{@{}L{2.9cm}YYYYY@{}}
\toprule
\textbf{Categoria} & \textbf{C1 Hospedagem própria} & \textbf{C2 Sem servidor de BD} & \textbf{C3 Edição pela web} & \textbf{C4 Fórmulas e código nativos} & \textbf{C5 Personalização sem código}\\
\midrule
Plataformas hospedadas & Não & Sim & Sim & Variável, em geral limitado & Limitada\\
\addlinespace
CMS tradicionais & Sim & Não & Sim & Via \en{plugins} ou temas & Sim, com curva de aprendizado\\
\addlinespace
CMS \en{headless} & Sim ou serviço & Varia & Sim (painel) & Depende do \en{front-end} & Exige construir o \en{front-end}\\
\addlinespace
Geradores estáticos e docs como código & Sim (estático) & Sim & Não, por padrão & Sim, por extensões & Via temas e CSS\\
\addlinespace
\textbf{Vitrine} & \textbf{Sim} & \textbf{Sim (SQLite)} & \textbf{Sim} & \textbf{Sim (nativo)} & \textbf{Sim (painel)}\\
\bottomrule
\end{tabularx}
\fonte{Elaborado pelos autores. Síntese qualitativa das tendências de cada categoria; não substitui a avaliação de um produto específico.}
\end{table}

O diferencial do Vitrine não está em uma funcionalidade isolada, mas na combinação: hospedagem própria sem servidor de banco de dados, edição pela web, matemática e código como parte nativa do texto, personalização visual sem programação e entrega com desempenho de página estática. Em contraste com ecossistemas fechados que dificultam a exportação de dados, a arquitetura privilegia a soberania do autor: o conteúdo reside em um arquivo de banco e em texto Markdown, o que permite ao pesquisador usar a ferramenta como a base central do seu conhecimento \cite{reis2026ic}.

\section{Síntese do Capítulo}

O capítulo situou o Vitrine entre as categorias de solução para publicação de conteúdo e justificou, com base na documentação das tecnologias, as escolhas de monólito modular, renderização incremental, SQLite com Prisma, better-auth, Markdown com \texttt{unified} e KaTeX, e personalização por variáveis CSS. Apresentou também os fundamentos de especificação e de \en{Documentation-Driven Development} que sustentam o processo descrito a seguir.

%% ===================================================================================
%% Capítulo 3 -- Processo de Desenvolvimento e Ambiente
%% Fontes: CLAUDE.md, ARQUITETO.md, DESENVOLVEDORES.md. Dono: Arquiteto/Redator.
%% Este capítulo descreve o processo; as regras operacionais vivem nos arquivos de
%% governança na raiz do repositório e devem permanecer coerentes.
%% ===================================================================================
\chapter{Processo de Desenvolvimento e Ambiente}\label{cap:metodologia}

Este capítulo descreve a abordagem e as etapas do desenvolvimento (Seção~\ref{sec:abordagem}), o planejamento por fases (Seção~\ref{sec:planejamento}) e, principalmente, o processo de desenvolvimento orientado por documentação e executado por agentes especializados (Seção~\ref{sec:processo-ddd}). Encerra com o ambiente técnico (Seção~\ref{sec:ambiente}), a estratégia de verificação (Seção~\ref{sec:estrategia-vv}) e os riscos do projeto (Seção~\ref{sec:riscos}).

\section{Abordagem e Etapas do Desenvolvimento}\label{sec:abordagem}

O Vitrine é construído de forma iterativa e incremental \cite{sommerville2011}: cada incremento entrega uma fatia funcional de ponta a ponta --- da persistência no SQLite à renderização no \en{front-end} ---, verificada antes do incremento seguinte. A partir da consolidação do monólito, a especificação passa a ser escrita antes da implementação e atualizada após cada integração (Seção~\ref{sec:processo-ddd}).

O processo percorre quatro macroetapas, fundamentadas nos processos clássicos de engenharia de \en{software} \cite{sommerville2011,pressman2021}:

\begin{enumerate}
  \item \textbf{Levantamento de requisitos:} identificação das necessidades dos usuários --- professores e desenvolvedores --- e definição dos requisitos funcionais (RF) e não funcionais (RNF);
  \item \textbf{Modelagem de sistemas:} uso de UML e DBML para visualizar a lógica do sistema, incluindo diagramas de casos de uso e modelagem de dados (ER);
  \item \textbf{Prototipagem:} criação de protótipos de baixa fidelidade para validar a usabilidade do editor de Markdown e do painel administrativo antes da implementação final;
  \item \textbf{Desenvolvimento incremental (\en{sprints}):} implementação do código-fonte com Next.js e Prisma ORM, da persistência no SQLite até a renderização no \en{front-end} e a configuração de cache (ISR).
\end{enumerate}

\section{Planejamento por Fases}\label{sec:planejamento}

A consolidação do sistema em arquitetura monolítica documentada é organizada em fases sucessivas. Cada fase termina com a integração das ramificações correspondentes na \texttt{main}, a atualização da documentação e a revisão do estado dos módulos (Seção~\ref{sec:modulos}). As fases não têm datas fixas neste documento: o avanço é registrado no mural de agentes e no Capítulo~\ref{cap:implementacao}.

\begin{table}[htb]
\caption{Fases planejadas da consolidação}\label{tab:fases}
\small
\begin{tabularx}{\linewidth}{@{}L{1.1cm}L{4.2cm}YL{2.8cm}@{}}
\toprule
\textbf{Fase} & \textbf{Nome} & \textbf{Entregas principais} & \textbf{Agentes}\\
\midrule
F1 & Especificação & Documentação \LaTeX{} inicial, diagramas como código, governança e mural & \agente{Arquiteto}\\
\addlinespace
F2 & Fundação & Esquema, migração e \en{seed}; acesso a dados; autenticação; proteção de rotas & \agente{Dados}, \agente{Core}\\
\addlinespace
F3 & Conteúdo e leitura & API de conteúdo; páginas públicas com ISR; \en{sidebar} e barra superior; renderização Markdown com \LaTeX{} & \agente{Core}, \agente{UI}\\
\addlinespace
F4 & Edição e painel & Editor com pré-visualização; gestão da árvore (criar, mover, ordenar, excluir) & \agente{UI}, \agente{Core}\\
\addlinespace
F5 & Personalização & Identidade e paleta por variáveis CSS; revalidação de \en{layout} & \agente{UI}, \agente{Core}, \agente{Dados}\\
\addlinespace
F6 & Verificação e validação & Testes automatizados, desempenho, segurança, acessibilidade e usabilidade & Todos\\
\addlinespace
F7 & Consolidação da documentação & Relato de implementação por módulo e situação da verificação & \agente{Arquiteto}\\
\bottomrule
\end{tabularx}
\fonte{Elaborado pelos autores.}
\end{table}

\section{Processo de Desenvolvimento: Documentation-Driven Development com Agentes}\label{sec:processo-ddd}

\subsection{Princípios}\label{sec:principios-ddd}

O processo segue sete princípios, que também constam dos arquivos de governança do repositório.

\begin{enumerate}[label=\textbf{PR\arabic*.}, leftmargin=1.6cm]
  \item \textbf{A documentação é a fonte da verdade.} Em caso de divergência entre código e documentação, a divergência é tratada como defeito de um dos dois e resolvida explicitamente; nunca é ignorada.
  \item \textbf{A especificação precede a implementação.} Nenhuma funcionalidade nova é implementada sem que exista, na documentação, o requisito, o contrato e o critério de aceite correspondentes.
  \item \textbf{Escritor único da documentação.} Somente o agente Arquiteto edita a documentação (\arq{doc/}) e os arquivos de governança; os demais agentes propõem mudanças por meio do mural.
  \item \textbf{Domínios com dono.} Cada diretório do código pertence a um agente; alterações fora do domínio são solicitadas ao dono.
  \item \textbf{Código apenas em ramificações.} É vedado alterar código ou configuração diretamente na \texttt{main}; a integração ocorre por \en{merge} autorizado pelo Arquiteto. Cada agente trabalha em um \en{worktree} próprio, e a cópia principal permanece na \texttt{main}.
  \item \textbf{O mural é o único canal de coordenação e só é editado na \texttt{main}} (na cópia principal). Isso serializa a comunicação e evita conflitos de fusão.
  \item \textbf{Rastreabilidade por identificadores.} Requisitos, casos de uso, regras, decisões e módulos possuem identificadores estáveis, citados nos \en{commits}, nos pedidos de integração e nos testes.
\end{enumerate}

\subsection{Papéis e Domínios}\label{sec:papeis}

O time é composto por quatro agentes. A Tabela~\ref{tab:papeis} resume a responsabilidade e o domínio de cada um; a distribuição de diretórios é normativa e está detalhada em \arq{DESENVOLVEDORES.md}.

\begin{table}[htb]
\caption{Papéis, domínios e diretórios}\label{tab:papeis}
\small
\begin{tabularx}{\linewidth}{@{}L{2.6cm}YL{5.3cm}@{}}
\toprule
\textbf{Agente} & \textbf{Responsabilidade} & \textbf{Diretórios e arquivos}\\
\midrule
\agente{Arquiteto} (Redator Técnico e \en{Merge Master}) & Governa a fonte da verdade; revisa pedidos de integração; autoriza e executa \en{merges}; atualiza a documentação; mantém o mural limpo & \arq{doc/}, \arq{CLAUDE.md}, \arq{ARQUITETO.md}, \arq{DESENVOLVEDORES.md}, \arq{MURAL_DE_AGENTES.md} (limpeza)\\
\addlinespace
\agente{Dados} & Modela o esquema e as rotinas de acesso a dados; migrações e \en{seed} & \arq{prisma/}, \arq{src/lib/db.js}, \arq{src/lib/content-tree.js}\\
\addlinespace
\agente{Core} & Lógica de servidor: API, autenticação, proteção de rotas, revalidação de cache; roteamento do App Router & \arq{src/app/api/}, \arq{src/proxy.js}, \arq{src/lib/auth*.js}, \arq{src/lib/validators.js}, \arq{src/lib/revalidate.js}, arquivos de rota (\arq{page.js}, \arq{layout.js}, \arq{route.js})\\
\addlinespace
\agente{UI} & Interface: componentes, estilo global, editor e pipeline Markdown & \arq{src/components/}, \arq{src/app/globals.css}, \arq{src/lib/theme.js}\\
\bottomrule
\end{tabularx}
\fonte{Elaborado pelos autores, com base em \arq{DESENVOLVEDORES.md} e \arq{ARQUITETO.md}.}
\end{table}

\subsection{Fontes da Verdade}\label{sec:fontes-verdade}

A Tabela~\ref{tab:fontes-verdade} relaciona cada tipo de informação ao artefato que o define. Quem implementa deve consultar esses arquivos antes de escrever código e citar os identificadores (RF, RNF, RN, UC, ADR, D) no pedido de integração.

\begin{table}[htb]
\caption{Fontes da verdade do projeto}\label{tab:fontes-verdade}
\small
\begin{tabularx}{\linewidth}{@{}L{3.3cm}YL{3.5cm}@{}}
\toprule
\textbf{Informação} & \textbf{Artefato (caminho)} & \textbf{Responsável}\\
\midrule
Requisitos, regras e casos de uso & \arq{doc/tex/capitulos/cap4_requisitos.tex}; \arq{doc/diagrams/src/uc-visao-geral.puml} & \agente{Arquiteto}\\
\addlinespace
Arquitetura, API, rotas e cache & \arq{doc/tex/capitulos/cap5_arquitetura.tex}; diagramas \arq{arq-*} e \arq{seq-*} & \agente{Arquiteto}\\
\addlinespace
Modelo de dados & Implementação: \arq{prisma/schema.prisma}; documentação: \arq{doc/diagrams/src/er-vitrine.dbml} e \arq{doc/tex/apendices/apendice_a_dicionario_dados.tex} & \agente{Dados} (código); \agente{Arquiteto} (documentação)\\
\addlinespace
Decisões de projeto & \arq{doc/tex/apendices/apendice_e_adr.tex} & \agente{Arquiteto}\\
\addlinespace
Estado de cada módulo & \arq{doc/tex/capitulos/cap6_implementacao.tex} e mapa de pendências (Apêndice~\ref{ap:pendencias}) & \agente{Arquiteto}\\
\addlinespace
Processo e regras de trabalho & \arq{CLAUDE.md}, \arq{ARQUITETO.md}, \arq{DESENVOLVEDORES.md} & \agente{Arquiteto}\\
\addlinespace
Comunicação e pedidos de integração & \arq{MURAL_DE_AGENTES.md} (somente na \texttt{main}) & Todos\\
\bottomrule
\end{tabularx}
\fonte{Elaborado pelos autores.}
\end{table}

\subsection{Fluxo de Git e Uso do Mural}\label{sec:fluxo-git}

O ciclo completo está na Figura~\ref{fig:fluxo-ddd-agentes}. Em resumo: o Arquiteto especifica e publica avisos no mural; cada desenvolvedor cria, no primeiro passo, um \en{worktree} próprio com uma ramificação a partir da \texttt{main} e implementa nele, conforme a especificação; ao terminar, vai à cópia principal, que permanece na \texttt{main}, apenas para registrar o pedido de integração no mural; o Arquiteto compara a ramificação com a \texttt{main}, autoriza ou veta, e após o \en{merge} atualiza a documentação e limpa o mural.

\vitrinefigura[tipo=diagrama, arquivo=fluxo-ddd-agentes, agente=Arquiteto, largura=0.95,
  origem={doc/diagrams/src/fluxo-ddd-agentes.puml},
  descricao={Diagrama de atividades com raias (Arquiteto, Desenvolvedor) mostrando especificação, implementação, pedido no mural, revisão, merge ou veto e atualização da documentação.}]
  {Ciclo de \en{Documentation-Driven Development} com agentes}{fig:fluxo-ddd-agentes}

O mecanismo de fusão assíncrona por ramificações e a regra de que o mural só é editado na \texttt{main} formam um protocolo simples: como cada pedido é um \en{commit} na \texttt{main}, o histórico do Git registra quem pediu o quê e quando (Figura~\ref{fig:fluxo-git-mural}).

\textbf{Cópia de trabalho própria.} Os agentes executam ao mesmo tempo e não se veem. Se dividissem a mesma pasta, a troca de ramificação de um deslocaria a ramificação ativa do outro, e um \en{commit} cairia na ramificação errada --- o que ocorreu em 02/10/2026, quando um \en{commit} de código foi parar na \texttt{main} e outro, na ramificação de um colega (risco R15). Por isso cada agente cria, como primeiro passo, um \en{worktree} (\texttt{git worktree add ../vitrine-<agente> -b <ramificação> main}); a pasta principal fica sempre na \texttt{main}, serve só ao mural e ao Arquiteto, e ninguém troca a ramificação dela; confere-se \texttt{git branch -{}-show-current} antes de cada \en{commit}; e o Arquiteto remove o \en{worktree} depois do \en{merge}. Um \en{commit} de código que chegue à \texttt{main} por engano não é apagado do histórico: é revisado, mantido e registrado como exceção no \en{merge} correspondente.

\vitrinefigura[tipo=diagrama, arquivo=fluxo-git-mural, agente=Arquiteto, largura=0.95,
  origem={doc/diagrams/src/fluxo-git-mural.puml},
  descricao={Diagrama de sequência entre Desenvolvedor, ramificação de feature, main, MURAL\_DE\_AGENTES.md e Arquiteto, incluindo os caminhos de aprovação e de veto.}]
  {Fluxo Git e solicitação de integração pelo mural}{fig:fluxo-git-mural}

\subsection{Critérios de Aceite de uma Integração}\label{sec:criterios-aceite}

Ao revisar um pedido, o Arquiteto verifica, no mínimo:

\begin{enumerate}
  \item a ramificação parte da \texttt{main} atualizada e altera apenas arquivos do domínio do autor (ou declara as exceções no pedido);
  \item o comportamento implementado corresponde aos requisitos, às regras e aos contratos citados (identificadores RF, RNF, RN, UC);
  \item alterações de esquema, de API e de dependências estão declaradas no pedido e são compatíveis com a documentação --- ou acompanham um pedido de atualização da documentação;
  \item \texttt{npm run lint} e \texttt{npm run build} passam, e os testes aplicáveis foram executados;
  \item nenhum segredo, arquivo gerado ou banco de dados é versionado;
  \item há evidências visuais (capturas de tela nos caminhos previstos) quando a mudança é visível ao usuário;
  \item a ramificação não altera o mural nem a documentação (\arq{doc/});
  \item os critérios de segurança aplicáveis foram atendidos: validação de entrada no servidor, sessão verificada no servidor e ausência de HTML bruto no conteúdo.
\end{enumerate}

Se algum critério falha, o Arquiteto não realiza o \en{merge} e responde no mural, apontando o que deve ser corrigido.

\subsection{Rastreabilidade e Identificadores}\label{sec:rastreabilidade}

Todos os artefatos usam identificadores estáveis: \texttt{UC} (casos de uso), \texttt{RF} e \texttt{RNF} (requisitos), \texttt{RN} (regras de negócio), \texttt{ADR} (decisões aceitas) e \texttt{D} (identificadores das decisões iniciais, hoje todas registradas nas ADR), \texttt{M} (módulos) e \texttt{CT} (casos de teste). A Seção~\ref{sec:matriz-rastreabilidade} relaciona requisitos, casos de uso e módulos, e o Capítulo~\ref{cap:testes} relaciona requisitos e casos de teste. As referências cruzadas no documento são hiperligações: clicar em \req{RF01} leva à definição do requisito.

\subsection{Convenções de Ramificações, \emph{Commits} e Idioma}

As ramificações seguem o padrão \arq{tipo/dominio-tema} --- por exemplo, \arq{feat/dados-schema-inicial} ou \arq{fix/core-sessao-expirada} --- e as mensagens de \en{commit}, o padrão \en{Conventional Commits} \cite{conventionalcommits} em português, com os identificadores afetados em rodapé (por exemplo, \texttt{Refs: RF05, RN03}). Identificadores de código e nomes de arquivos são em inglês; comentários, textos de interface e documentação são em português do Brasil.

\subsection{Uso de Agentes de Inteligência Artificial e Responsabilidade}\label{sec:ia-autoria}

Os agentes de inteligência artificial são ferramentas de apoio sob supervisão humana: não assumem responsabilidade pelo produto. O Engenheiro Chefe (autor do projeto) é o responsável pela revisão do código e da documentação que integram o produto, pelas decisões de projeto e pela veracidade das referências. O registro das interações relevantes --- pedidos, vetos e decisões --- é preservado no histórico do Git, o que permite auditar como cada parte do sistema foi produzida.

\section{Ambiente e Ferramentas}\label{sec:ambiente}

A Tabela~\ref{tab:stack} lista as tecnologias e as versões efetivamente declaradas no repositório, conferidas em \arq{package.json} e em \arq{node_modules}; os valores foram atualizados após o merge M-001 (30/09/2026) e revistos até o merge M-011 (03/10/2026).

\begin{table}[htb]
\caption{Tecnologias e versões do projeto}\label{tab:stack}
\small
\begin{tabularx}{\linewidth}{@{}L{4.0cm}L{3.0cm}Y@{}}
\toprule
\textbf{Tecnologia} & \textbf{Versão} & \textbf{Papel}\\
\midrule
Node.js & $\geq$ 20.9 (24 no desenvolvimento e na imagem) & \en{Runtime} exigido pelo Next.js 16. O \arq{jsdom} dos testes de interação (merge M-011) exige 22.22.2, 24.15 ou superior: com o 22.22.0 o npm só avisa (\texttt{EBADENGINE}) e os testes passam\\
Next.js & 16.3.7 & \en{Framework} web (App Router, ISR, Route Handlers, Proxy)\\
React e React DOM & 19.2.8 & Biblioteca de interface\\
React Compiler (plugin Babel) & 1.0.0 & Memoização automática (\texttt{reactCompiler: true})\\
Tailwind CSS & 4.3.3 (instalada; \texttt{\^{}4}) & Estilização por utilitários\\
Prisma ORM e \arq{@prisma/client} & 7.10.0 & Acesso a dados e migrações; configuração em \arq{prisma.config.ts}; cliente gerado em \arq{src/generated/prisma} (gerador \texttt{prisma-client})\\
\arq{@prisma/adapter-better-sqlite3} e \arq{better-sqlite3} & 7.10.0 e 12.11.1 & \en{Driver adapter} do Prisma para SQLite (módulo nativo) e o próprio SQLite em arquivo (\adr{ADR-014})\\
\arq{prisma} (CLI) e \arq{dotenv} & 7.10.0 e 18.0.5 (\texttt{dependencies}) & Migrações e carregamento do \arq{.env} pela CLI; necessários na imagem de produção (\texttt{prisma migrate deploy})\\
Vitest & 5.0.3 (desenvolvimento) & Testes de unidade e de integração com banco SQLite isolado (\texttt{npm test}; Seção~\ref{sec:testes-banco})\\
\arq{jsdom}, \arq{@testing-library/react} e \arq{@testing-library/dom} & 30.1.1, 16.3.3 e 10.4.2 (desenvolvimento; MIT) & Testes de interação dos Componentes Cliente (cliques, estado, \texttt{localStorage}): cada arquivo de teste escolhe o ambiente com \texttt{// @vitest-environment jsdom}, sem mudar a configuração do Vitest; 0~byte no \en{bundle} (merge M-011)\\
\arq{@playwright/test} & 1.63.0 (desenvolvimento; Apache-2.0) & Testes de ponta a ponta no navegador, sobre o \en{build} de produção (\texttt{npm run test:e2e}; \decisao{D-12}; merge M-010). Os navegadores ficam fora do repositório e da imagem; os três pacotes (\texttt{@playwright/test}, \texttt{playwright} e \texttt{playwright-core}) entram na imagem como qualquer dependência de desenvolvimento, cerca de 18~MiB em \arq{node_modules}\\
\arq{tsx} & 4.23.15 (desenvolvimento) & Execução do \en{seed} (\arq{prisma/seed.js}) com o cliente gerado em TypeScript\\
better-auth & 1.7.6 & Autenticação e sessões (e-mail e senha, cadastro público desabilitado, limitação de tentativas)\\
react-markdown & 10.1.0 & Renderização de Markdown em React\\
remark-math e rehype-katex & 6.0.0 e 7.0.1 & Delimitadores e renderização de fórmulas\\
remark-gfm & 4.0.1 & Tabelas, tachado, listas de tarefas, links automáticos e notas de rodapé (RF11; \decisao{D-08}; merge M-009)\\
rehype-highlight, \arq{lowlight} e \arq{highlight.js} & 7.0.2, 3.3.0 e 11.11.2 & Realce de sintaxe por classes, no servidor e no cliente (RF10; \decisao{D-07}; merge M-009; o \arq{highlight.js} é BSD-3-Clause, os demais são MIT). No merge M-011, o \arq{highlight.js} (fixado em \texttt{\~{}11.11.2}) e o \arq{lowlight} passaram a dependências diretas, com LaTeX, MATLAB, Julia, Haskell e Fortran (42 linguagens; mais 6,3~kB comprimidos)\\
react-hook-form & 7.89.0 (MIT) & Formulário do editor, com validação por campo e foco no primeiro erro, sobre os esquemas do Zod por um resolvedor próprio (M09; merge M-011). Só entra no código do editor, carregado sob demanda (12,4~kB comprimidos)\\
KaTeX & 0.16.47 & Composição de fórmulas; declarado diretamente no merge M-009, pois o \arq{MarkdownRenderer} importa a folha de estilos do pacote\\
Zod & 4.6.5 (\texttt{dependencies}) & Validação de entrada da API e do comando do administrador (\decisao{D-11}); declarado diretamente no merge M-004 (já era dependência do better-auth)\\
ESLint e \arq{eslint-config-next} & 9.x e 16.3.7 & Análise estática\\
\LaTeX{} (MiKTeX) e \texttt{latexmk} & pdfTeX 4.19 & Composição desta documentação\\
PlantUML e Graphviz & a instalar & Renderização dos diagramas (\arq{doc/diagrams/})\\
Git & --- & Controle de versão e fluxo por ramificações\\
Docker Engine e Docker Compose & versão do servidor & Implantação por imagem em VPS tradicional (\adr{ADR-011})\\
\bottomrule
\end{tabularx}
\fonte{Elaborado pelos autores, a partir de \arq{package.json} e das instalações locais.}
\end{table}

\section{Estratégia de Verificação e Validação}\label{sec:estrategia-vv}

A verificação (o sistema foi construído corretamente?) e a validação (foi construído o sistema correto?) serão realizadas em quatro frentes: (i) testes automatizados de unidade, integração e ponta a ponta sobre os casos de uso; (ii) medições de desempenho e de comportamento do cache ISR; (iii) verificações de segurança e acessibilidade orientadas pelo OWASP e pelo WCAG; e (iv) avaliação de usabilidade com usuários, com o questionário SUS. O plano detalhado, com critérios e casos de teste, está no Capítulo~\ref{cap:testes}.

\section{Riscos do Projeto}\label{sec:riscos}

A Tabela~\ref{tab:riscos} registra os principais riscos identificados, com probabilidade (P) e impacto (I) qualitativos --- baixo (B), médio (M) ou alto (A) --- e as mitigações adotadas.

\begin{longtable}{@{}L{1.0cm}L{5.0cm}L{0.7cm}L{0.7cm}L{6.9cm}@{}}
\caption{Registro de riscos}\label{tab:riscos}\\
\toprule
\textbf{ID} & \textbf{Risco} & \textbf{P} & \textbf{I} & \textbf{Mitigação}\\
\midrule
\endfirsthead
\toprule
\textbf{ID} & \textbf{Risco} & \textbf{P} & \textbf{I} & \textbf{Mitigação}\\
\midrule
\endhead
\bottomrule
\endfoot
R01 & O sistema de arquivos efêmero de plataformas \en{serverless} impede a persistência do arquivo SQLite & B & A & Mitigado pela decisão \decisao{D-04}: implantação por imagem Docker em VPS com volume persistente (\adr{ADR-011}); plataformas \en{serverless} não são alvo desta versão\\
R02 & Divergência entre código e documentação (deriva) & M & A & Princípios PR1 e PR2; critérios de aceite; Arquiteto como revisor único\\
R03 & Conflitos de fusão em \arq{package-lock.json} quando várias ramificações adicionam dependências & A & B & Novas dependências declaradas no pedido; o Arquiteto regenera o \en{lockfile} ao integrar\\
R04 & Mudanças incompatíveis nas dependências (Next.js 16, Prisma 7, better-auth 1.x) em relação ao conhecimento dos agentes & M & M & Versões fixadas no \en{lockfile}; consulta obrigatória à documentação local e às \en{skills} do Prisma antes de codificar\\
R05 & Injeção de HTML/\en{script} (XSS) por meio do Markdown & M & A & Não habilitar HTML bruto; testes de segurança; revisão de dependências de Markdown\\
R06 & Cores escolhidas com contraste insuficiente & M & M & Aviso de contraste no painel (WCAG 2.1 AA); valores padrão acessíveis\\
R07 & Limite de concorrência de escrita do SQLite & B & M & Modo WAL; poucas escritas; transações curtas\\
R08 & Agentes introduzem código fora do escopo ou sem testes & M & M & Domínios com dono; critérios de aceite; veto fundamentado no mural\\
R09 & Perda de dados por exclusão em cascata acidental & M & A & Confirmação explícita; rotina de \en{backup} documentada (Apêndice~\ref{ap:implantacao})\\
R10 & Escopo maior do que a capacidade de entrega do time & M & A & Priorização Essencial, Importante e Desejável; fases com entregas incrementais\\
R11 & Vazamento de rascunhos (\texttt{DRAFT}) por consulta pública sem filtro de estado, na página, na \en{sidebar}, na geração estática ou no mapa do site & M & A & Regra RN15; filtro centralizado em \arq{db.js}; caso de teste CT31; revisão de todo \en{endpoint} ou página nova\\
R12 & Perda do volume do SQLite na VPS (falha de disco ou do provedor) & M & A & Rotina de cópia de segurança com \texttt{VACUUM INTO} e cópias fora da VPS (Apêndice~\ref{ap:implantacao}); teste de restauração (CT30)\\
R13 & Alertas de vulnerabilidade do \texttt{npm audit} em dependências transitivas do Prisma 7 (\arq{mysql2} e \arq{deepmerge-ts}), cuja correção automática rebaixaria o Prisma para a versão 6 & B & M & Alertas avaliados e aceitos (Seção~\ref{sec:arq-npm-audit}, \adr{ADR-019}): os pacotes não são exercidos pelo Vitrine; \texttt{npm audit} a cada entrega e nova avaliação quando o Prisma publicar versão corrigida\\
R14 & Perda do acesso ao painel por senha esquecida do administrador único, sem envio de e-mail nem redefinição & B & A & Guardar a senha com segurança; evolução prevista de um comando de redefinição (Seção~\ref{sec:trabalhos-futuros}); \en{backup} do banco (Apêndice~\ref{ap:implantacao})\\
R15 & Agentes simultâneos na mesma cópia de trabalho: a troca de ramificação de um desloca a do outro, e \en{commits} caem na ramificação errada (ocorrido em 02/10/2026, com um \en{commit} de código na \texttt{main}) & A & M & \en{Worktree} próprio por agente desde o primeiro passo (PR5; Seção~\ref{sec:fluxo-git}); cópia principal sempre na \texttt{main}; conferir a ramificação antes de cada \en{commit}; o Arquiteto remove o \en{worktree} após o \en{merge}\\
\end{longtable}
\fonte{Elaborado pelos autores.}

%% ===================================================================================
%% Capítulo 4 -- Levantamento e Análise de Requisitos
%% FONTE DA VERDADE dos requisitos, das regras de negócio e dos casos de uso.
%% Origem: Relatório de IC (RF01 a RF04, diagrama de casos de uso, protótipos) + derivação
%% feita na consolidação do monólito (RF05 em diante, RNF e RN).
%% Dono: Arquiteto/Redator. Desenvolvedores: citem os identificadores (RF, RNF, RN, UC)
%% nos commits e nos pedidos do mural; peçam mudanças pelo mural.
%% ===================================================================================
\chapter{Levantamento e Análise de Requisitos}\label{cap:requisitos}

Este capítulo especifica \emph{o que} o Vitrine deve fazer e sob que restrições. Os requisitos RF01 a RF04 foram definidos no relatório de IC \cite{reis2026ic}; os demais requisitos funcionais, os requisitos não funcionais e as regras de negócio foram derivados dos casos de uso, dos protótipos e dos objetivos (Seção~\ref{sec:objetivos-especificos}) durante a consolidação do monólito, e estão marcados com a sua origem. Todos os requisitos e regras deste capítulo, inclusive RF26 e RN15 a RN17, foram aprovados pelo Engenheiro Chefe e são oficiais. A especificação segue as recomendações de identificação única, verificabilidade e rastreabilidade da ISO/IEC/IEEE 29148 \cite{iso29148}.

\section{Visão Geral do Produto}\label{sec:visao-produto}

\textbf{Propósito.} Permitir que docentes e pesquisadores publiquem e organizem conteúdo científico e didático --- artigos, notas de aula, tutoriais --- em um site próprio, com identidade visual personalizada, edição em Markdown com fórmulas \LaTeX{} e entrega com desempenho de página estática.

\textbf{Perspectiva.} O Vitrine é uma aplicação web autônoma, implantada como um contêiner Docker com um único processo Node.js, acompanhado de um arquivo de banco de dados SQLite em volume persistente. Não depende de serviços externos obrigatórios (Figura~\ref{fig:arq-contexto}, Capítulo~\ref{cap:arquitetura}).

\textbf{Ambiente operacional.} Navegadores modernos no cliente; no servidor, uma VPS tradicional com Docker, volume persistente em disco local e, em produção, HTTPS (decisão \decisao{D-04}, \adr{ADR-011}).

\textbf{Restrições.} (i) A pilha tecnológica é Next.js, Prisma, SQLite e better-auth; (ii) a interface é em português do Brasil; (iii) o código é aberto, sob licença MIT (\adr{ADR-015}); (iv) a capacidade de entrega do time é limitada, o que exige priorização (\emph{risco} R10, Tabela~\ref{tab:riscos}).

\textbf{Premissas e dependências.} Assume-se a disponibilidade de uma VPS com Docker e armazenamento persistente local (\decisao{D-04}); a estabilidade das interfaces de programação das versões fixadas (Tabela~\ref{tab:stack}); e a existência de um único administrador na versão inicial (\decisao{D-13}).

\section{Atores e Partes Interessadas}\label{sec:atores}

\begin{table}[htb]
\caption{Atores do sistema e partes interessadas}\label{tab:atores}
\small
\begin{tabularx}{\linewidth}{@{}L{2.8cm}L{2.4cm}Y@{}}
\toprule
\textbf{Ator ou parte} & \textbf{Tipo} & \textbf{Descrição e interesse}\\
\midrule
Leitor & Ator primário & Visitante --- estudante, colega, pesquisador --- que navega e lê o conteúdo publicado; quer encontrar e ler rápido, em qualquer dispositivo. Não se autentica.\\
\addlinespace
Administrador & Ator primário & Docente-autor que cria, organiza e publica conteúdo e define a identidade visual do site. Tem conhecimentos básicos de informática e não deseja administrar infraestrutura.\\
\addlinespace
Sistema & Ator secundário & Automação interna que dispara a invalidação do cache ISR após mutações (caso de uso UC12).\\
\addlinespace
Implantador & Parte interessada & Quem instala e opera o sistema --- em geral o próprio docente com apoio técnico ou um desenvolvedor. Interessa-se por simplicidade de instalação, configuração e \en{backup}.\\
\addlinespace
Engenheiro Chefe & Parte interessada & Autor e responsável pelo produto; decide o escopo e as questões de projeto e aprova as integrações e as regras de negócio.\\
\addlinespace
Agentes de desenvolvimento & Parte interessada & Dados, Core e UI; consomem esta especificação como contrato de implementação.\\
\bottomrule
\end{tabularx}
\fonte{Elaborado pelos autores, com base em \citeonline{reis2026ic}.}
\end{table}

\section{Casos de Uso}\label{sec:casos-uso}

\subsection{Diagrama de Casos de Uso}

A Figura~\ref{fig:uc-visao-geral} apresenta o diagrama de casos de uso do Vitrine. Ele parte do diagrama do relatório de IC (Figura~\ref{fig:ic-uc}), com uma alteração: \emph{Gerenciar Conteúdo} passa a incluir \emph{Invalidar Cache (ISR)}, porque toda mutação de conteúdo --- e não apenas a personalização --- deve refletir-se nas páginas públicas (\req{RF17}).

\vitrinefigura[tipo=diagrama, arquivo=uc-visao-geral, agente=Arquiteto, largura=0.95,
  origem={doc/diagrams/src/uc-visao-geral.puml},
  descricao={Diagrama de casos de uso com os atores Leitor, Administrador e Sistema, os pacotes Navegação e Leitura, Gestão de Estrutura e Configurações e Performance, e as relações include e extend.}]
  {Diagrama de casos de uso do Vitrine}{fig:uc-visao-geral}

\vitrinefigura[tipo=ic, arquivo=ic-fig1-casos-de-uso, largura=0.4,
  fonte={Relatório de IC \cite{reis2026ic}, Figura 1.}]
  {Diagrama de casos de uso original, do relatório de IC}{fig:ic-uc}

\subsection{Catálogo de Casos de Uso}

\begin{longtable}{@{}L{1.2cm}L{4.4cm}L{2.5cm}L{5.0cm}@{}}
\caption{Catálogo de casos de uso}\label{tab:catalogo-uc}\\
\toprule
\textbf{ID} & \textbf{Caso de uso} & \textbf{Ator} & \textbf{Descrição resumida}\\
\midrule
\endfirsthead
\toprule
\textbf{ID} & \textbf{Caso de uso} & \textbf{Ator} & \textbf{Descrição resumida}\\
\midrule
\endhead
\bottomrule
\endfoot
UC01 & Explorar \en{sidebar} hierárquica & Leitor & Navegar pela árvore de coleções e documentos e pelas coleções raiz na barra superior.\\
UC02 & Visualizar página única & Leitor & Abrir um documento por sua URL amigável.\\
UC03 & Ler conteúdo (Markdown) & Leitor & Ler o documento renderizado, com fórmulas, tabelas e código. Incluído por UC01 e UC02.\\
UC04 & Gerenciar conteúdo & Administrador & Criar, editar, definir o estado de publicação (rascunho ou publicado), mover, ordenar e excluir conteúdos no painel. Inclui UC07, UC08 e UC12.\\
UC05 & Criar documento & Administrador & Criar conteúdo do tipo DOCUMENT. Estende UC04.\\
UC06 & Criar coleção/pasta & Administrador & Criar conteúdo do tipo COLLECTION. Estende UC04.\\
UC07 & Aninhar e ordenar conteúdos & Administrador & Definir pai e ordem de um conteúdo. Incluído por UC04.\\
UC08 & Autenticar-se & Administrador & Entrar e sair da área administrativa. Incluído por UC04 e UC09.\\
UC09 & Personalizar UI & Administrador & Alterar a aparência do site. Inclui UC08 e UC12.\\
UC10 & Definir identidade (nome e \en{branding}) & Administrador & Definir nome, descrição, logotipo, \en{favicon} e links sociais do rodapé. Estende UC09.\\
UC11 & Configurar paleta de cores & Administrador & Escolher as cores do tema. Estende UC09.\\
UC12 & Invalidar cache (ISR) & Sistema & Invalidar as páginas afetadas após uma mutação.\\
\end{longtable}
\fonte{Elaborado pelos autores, a partir da Figura~\ref{fig:uc-visao-geral}.}

\subsection{Especificações Detalhadas}\label{sec:uc-detalhados}

Os casos de uso com maior impacto sobre a arquitetura são especificados a seguir. Os demais seguem os mesmos padrões e são detalhados conforme a necessidade de cada módulo.

\begin{table}[htb]
\caption{Especificação do caso de uso UC08 --- Autenticar-se}\label{tab:uc08}
\small
\begin{tabularx}{\linewidth}{@{}L{3.0cm}Y@{}}
\toprule
\textbf{Ator} & Administrador\\
\textbf{Pré-condição} & Existe ao menos um usuário cadastrado (bootstrap, \req{RF22}).\\
\textbf{Gatilho} & Acesso a uma rota protegida sem sessão válida, ou acesso direto a \arq{/login}.\\
\textbf{Fluxo principal} & (1) O sistema exibe o formulário de e-mail e senha. (2) O administrador informa as credenciais. (3) O sistema verifica as credenciais e cria uma sessão persistida no banco. (4) O sistema grava o \en{cookie} de sessão (\en{HttpOnly}) e redireciona ao destino original ou a \arq{/admin}.\\
\textbf{Fluxos alternativos} & A1: credenciais inválidas --- mensagem genérica, que não revela qual campo falhou; limitação de tentativas (\req{RNF03}). A2: sessão expirada --- retorno ao formulário preservando o destino (parâmetro \texttt{redirectTo}, que aceita apenas caminhos internos do site; destinos externos levam ao painel). A3: encerrar sessão (\req{RF20}) --- a sessão é invalidada e o \en{cookie} removido.\\
\textbf{Pós-condição} & Sessão válida associada ao usuário; operações de escrita autorizadas (\req{RF21}).\\
\textbf{Requisitos e regras} & \req{RF01}, \req{RF20}, \req{RF21}, \req{RF22}; RN01, RN12.\\
\bottomrule
\end{tabularx}
\fonte{Elaborado pelos autores.}
\end{table}

\begin{table}[htb]
\caption{Especificação do caso de uso UC05 --- Criar documento}\label{tab:uc05}
\small
\begin{tabularx}{\linewidth}{@{}L{3.0cm}Y@{}}
\toprule
\textbf{Ator} & Administrador (autenticado, UC08)\\
\textbf{Pré-condição} & Sessão válida. O documento será criado na raiz ou dentro de uma coleção existente.\\
\textbf{Fluxo principal} & (1) O administrador aciona ``Nova publicação'' (na raiz ou na coleção desejada). (2) Informa o título; o \en{slug} é sugerido a partir do título e pode ser editado; escolhe a coleção pai e o estado de publicação (padrão: Publicado). (3) Escreve o texto em Markdown, acompanhando a pré-visualização. (4) Aciona ``Enviar nova versão''. (5) O sistema valida os dados (RN02, RN03, RN06), persiste o conteúdo com o usuário como autor e (6) dispara a invalidação do cache (UC12). (7) O painel confirma a publicação.\\
\textbf{Fluxos alternativos} & A1: \en{slug} já existente entre os irmãos --- o sistema responde com conflito e sugere um \en{slug} alternativo. A2: o pai escolhido não é uma coleção --- erro de validação. A3: o administrador opta por ``Salvar localmente'' e retoma depois (\req{RF15}).\\
\textbf{Pós-condição} & Documento persistido. Se ele e todos os seus ancestrais estiverem publicados, a \en{sidebar} e a página refletem a mudança a partir da próxima requisição; um rascunho só é visível no painel (RN15).\\
\textbf{Requisitos e regras} & \req{RF05}, \req{RF08}, \req{RF09}, \req{RF15}, \req{RF16}, \req{RF17}, \req{RF26}; RN01 a RN06, RN10, RN11, RN13, RN15, RN16.\\
\bottomrule
\end{tabularx}
\fonte{Elaborado pelos autores.}
\end{table}

\begin{table}[htb]
\caption{Especificação do caso de uso UC07 --- Aninhar e ordenar conteúdos}\label{tab:uc07}
\small
\begin{tabularx}{\linewidth}{@{}L{3.0cm}Y@{}}
\toprule
\textbf{Ator} & Administrador (autenticado, UC08)\\
\textbf{Pré-condição} & Sessão válida; existem conteúdos a reorganizar.\\
\textbf{Fluxo principal} & (1) O administrador visualiza a árvore de conteúdos no painel. (2) Arrasta um item para outra posição ou outro pai, ou usa as ações ``subir'' e ``descer''. (3) A interface atualiza a árvore de forma otimista e envia o lote de alterações. (4) O sistema valida o lote (pai do tipo coleção, ausência de ciclos, \en{slugs} únicos no destino) e o persiste em uma transação. (5) O sistema invalida o \en{layout} raiz (UC12) e devolve a árvore atualizada.\\
\textbf{Fluxos alternativos} & A1: movimento inválido --- o sistema rejeita o lote inteiro, a interface desfaz a atualização otimista e exibe o motivo.\\
\textbf{Pós-condição} & Nova estrutura e nova ordem persistidas; a \en{sidebar} pública reflete a mudança.\\
\textbf{Requisitos e regras} & \req{RF02}, \req{RF03}, \req{RF17}; RN02, RN03, RN05, RN06.\\
\bottomrule
\end{tabularx}
\fonte{Elaborado pelos autores.}
\end{table}

\begin{table}[htb]
\caption{Especificação do caso de uso UC11 --- Configurar paleta de cores}\label{tab:uc11}
\small
\begin{tabularx}{\linewidth}{@{}L{3.0cm}Y@{}}
\toprule
\textbf{Ator} & Administrador (autenticado, UC08)\\
\textbf{Pré-condição} & Sessão válida.\\
\textbf{Fluxo principal} & (1) O administrador abre Configurações. (2) Escolhe as cores primária, secundária e de destaque por seletor ou por código hexadecimal. (3) Pré-visualiza o resultado no próprio painel, com variáveis CSS locais. (4) Salva. (5) O sistema valida o formato e calcula o contraste, devolvendo avisos se for o caso. (6) Persiste a configuração e invalida o \en{layout} raiz. (7) Os leitores veem a nova paleta na próxima requisição, sem novo \en{build}.\\
\textbf{Fluxos alternativos} & A1: valor fora do formato \texttt{\#RRGGBB} --- erro de validação no campo. A2: contraste abaixo de 4,5:1 entre texto e fundo --- aviso, sem bloquear o salvamento (RN09).\\
\textbf{Pós-condição} & Configuração persistida; páginas regeneradas com as novas variáveis CSS.\\
\textbf{Requisitos e regras} & \req{RF04}, \req{RF17}, \req{RF19}; RN08, RN09, RN10.\\
\bottomrule
\end{tabularx}
\fonte{Elaborado pelos autores.}
\end{table}

\section{Requisitos Funcionais}\label{sec:rf}

A Tabela~\ref{tab:rf} lista os requisitos funcionais. A coluna \emph{Origem} indica se o requisito foi herdado do relatório de IC (``IC'') ou derivado nesta especificação a partir de casos de uso (UC), protótipos, objetivos (OE), regras (RN) ou decisões (D). A prioridade segue a escala \emph{Essencial}, \emph{Importante} e \emph{Desejável}. A coluna \emph{Agente} indica o responsável principal pela implementação; requisitos que atravessam domínios são coordenados pelo mural.

\begin{longtable}{@{}L{1.3cm}L{7.2cm}L{1.6cm}L{2.2cm}L{1.8cm}@{}}
\caption{Requisitos funcionais}\label{tab:rf}\\
\toprule
\textbf{ID} & \textbf{Requisito} & \textbf{Prior.} & \textbf{Origem} & \textbf{Agente}\\
\midrule
\endfirsthead
\toprule
\textbf{ID} & \textbf{Requisito} & \textbf{Prior.} & \textbf{Origem} & \textbf{Agente}\\
\midrule
\endhead
\bottomrule
\endfoot
\reqrow{RF01}{Autenticação}{O acesso à área administrativa exige e-mail e senha, com sessão persistida em banco (better-auth).}{Essencial}{IC (RF01); UC08}{Core}
\reqrow{RF02}{Estrutura hierárquica}{Os conteúdos formam uma árvore de coleções e documentos aninhados, por auto-relacionamento (pai e filhos).}{Essencial}{IC (RF02); UC04, UC07}{Dados}
\reqrow{RF03}{Ordenação}{Cada conteúdo possui uma ordem relativa entre irmãos, editável pelo administrador.}{Essencial}{IC (RF03); UC07}{Dados}
\reqrow{RF04}{Personalização}{Interface para alterar identidade visual e cores por variáveis CSS. Requisito agregador, detalhado em RF18 e RF19.}{Essencial}{IC (RF04); UC09}{UI}
\reqrow{RF05}{Criar documento}{Criar conteúdo do tipo DOCUMENT com título, \en{slug}, pai, estado de publicação (padrão PUBLISHED) e corpo em Markdown.}{Essencial}{UC05}{Core}
\reqrow{RF06}{Criar coleção}{Criar conteúdo do tipo COLLECTION (pasta) para agrupar documentos e subcoleções, com estado de publicação (padrão PUBLISHED).}{Essencial}{UC06}{Core}
\reqrow{RF07}{Editar e excluir}{Atualizar título, \en{slug}, corpo, estado, pai e ordem; excluir com confirmação e remoção em cascata dos descendentes.}{Essencial}{UC04; RN07}{Core}
\reqrow{RF08}{Editor Markdown}{Editor com barra de ferramentas e pré-visualização em tempo real (painel dividido).}{Essencial}{IC (OE2); Fig. 6}{UI}
\reqrow{RF09}{Fórmulas matemáticas}{Renderizar fórmulas \LaTeX{} em linha e em bloco no conteúdo e na pré-visualização (KaTeX).}{Essencial}{IC (OE2); UC03}{UI}
\reqrow{RF10}{Realce de sintaxe}{Realçar blocos de código por linguagem, no conteúdo e na pré-visualização.}{Importante}{IC (OE2); UC03}{UI}
\reqrow{RF11}{Markdown estendido}{Suportar tabelas, listas de tarefas, texto tachado e links automáticos (GFM).}{Importante}{IC (Fig. 3); UC03}{UI}
\reqrow{RF12}{\en{Sidebar} hierárquica}{Exibir a árvore da coleção corrente, com o item ativo destacado, níveis recolhíveis e navegação por teclado.}{Essencial}{UC01; Fig. 3}{UI}
\reqrow{RF13}{Barra superior}{Exibir logotipo, nome do site e as coleções raiz como abas ou links de navegação.}{Essencial}{UC01; Fig. 3}{UI}
\reqrow{RF14}{Página de documento}{Exibir o documento em URL derivada da cadeia de \en{slugs}; responder 404 a caminhos inexistentes ou não publicados (RN15).}{Essencial}{UC02; RN04}{Core}
\reqrow{RF15}{Rascunho local}{Guardar rascunho no navegador (``Salvar localmente'') sem publicá-lo.}{Desejável}{IC (Fig. 6)}{UI}
\reqrow{RF16}{Publicar versão}{Enviar ao servidor a versão editada (``Enviar nova versão''), substituindo o conteúdo publicado.}{Essencial}{IC (Fig. 6); UC05}{Core}
\reqrow{RF17}{Invalidação de cache}{Revalidar as páginas afetadas após qualquer mutação de conteúdo ou de configuração.}{Essencial}{IC (UC12); RN10}{Core}
\reqrow{RF18}{Identidade do site}{Configurar nome, descrição, logotipo e \en{favicon}.}{Essencial}{UC10; Fig. 5}{Core}
\reqrow{RF19}{Paleta de cores}{Configurar as cores primária, secundária e de destaque, aplicadas por variáveis CSS sem novo \en{build}.}{Essencial}{IC (OE3); UC11}{UI}
\reqrow{RF20}{Encerrar sessão}{Permitir \en{logout}, invalidando a sessão no servidor.}{Essencial}{UC08}{Core}
\reqrow{RF21}{Autorização de escrita}{Bloquear, no servidor, toda operação de escrita sem sessão válida.}{Essencial}{UC08; RN01}{Core}
\reqrow{RF22}{Bootstrap do administrador}{Criar o primeiro administrador na instalação, pelo comando \texttt{npm run admin:create}, a partir das variáveis de ambiente \texttt{ADMIN\_EMAIL}, \texttt{ADMIN\_PASSWORD} e, opcionalmente, \texttt{ADMIN\_NAME}; o comando é idempotente, recusa um segundo usuário com outro e-mail e não imprime a senha. O cadastro público é desabilitado e não há página de primeiro acesso.}{Essencial}{D-03; ADR-018}{Core}
\reqrow{RF23}{Configuração inicial}{Inserir a configuração padrão (registro único) na instalação.}{Essencial}{RN08}{Dados}
\reqrow{RF24}{Rodapé configurável}{Exibir rodapé com os links sociais configurados: seis campos opcionais e independentes (Discord, GitHub, YouTube, Instagram, LinkedIn e Twitch), apenas URLs http(s); campo vazio não exibe o ícone (RN17). O texto de direitos é gerado a partir do nome do site e do ano.}{Desejável}{IC (Fig. 3 a 5); D-19; ADR-013}{UI}
\reqrow{RF25}{Metadados de SEO}{Gerar título, descrição e metadados por página; expor \arq{robots.txt} e mapa do site, este somente com conteúdo publicado (RN15).}{Importante}{IC (OE1); RNF07}{Core}
\reqrow{RF26}{Estado de publicação}{Cada conteúdo tem estado DRAFT (rascunho) ou PUBLISHED (publicado), com padrão PUBLISHED. Somente conteúdo publicado, com todos os ancestrais publicados, é exibido ao público; o painel mostra todos os conteúdos e permite alternar o estado.}{Essencial}{D-15; ADR-012; RN15, RN16}{Core}
\end{longtable}
\fonte{Elaborado pelos autores. RF01 a RF04: \citeonline{reis2026ic}.}

\section{Requisitos Não Funcionais}\label{sec:rnf}

Os requisitos não funcionais (Tabela~\ref{tab:rnf}) estão associados às características de qualidade da ISO/IEC 25010 \cite{iso25010}. As metas numéricas são as metas oficiais do produto e são verificadas no Capítulo~\ref{cap:testes}.

\begin{longtable}{@{}L{1.4cm}L{7.6cm}L{2.7cm}L{1.9cm}@{}}
\caption{Requisitos não funcionais}\label{tab:rnf}\\
\toprule
\textbf{ID} & \textbf{Requisito} & \textbf{Característica (ISO 25010)} & \textbf{Prior.}\\
\midrule
\endfirsthead
\toprule
\textbf{ID} & \textbf{Requisito} & \textbf{Característica (ISO 25010)} & \textbf{Prior.}\\
\midrule
\endhead
\bottomrule
\endfoot
\reqtarget{RNF01} & \textbf{Portabilidade e implantação simples.} Implantável a partir do repositório com poucos comandos (instalar, configurar variáveis de ambiente, migrar, compilar e iniciar), com um único arquivo de banco e nenhum serviço externo obrigatório. & Portabilidade & Essencial\\[1pt]
\reqtarget{RNF02} & \textbf{Desempenho de leitura.} Páginas públicas servidas do cache ISR. Meta: LCP $\leq$ 2,5 s (percentil 75) e pontuação de desempenho $\geq$ 90 no Lighthouse, em perfil móvel simulado. & Eficiência de desempenho & Importante\\[1pt]
\reqtarget{RNF03} & \textbf{Segurança.} Senhas somente como \en{hash}; \en{cookies} de sessão \texttt{HttpOnly}, \texttt{SameSite} e \texttt{Secure} em produção; validação de toda entrada no servidor; proibição de HTML bruto no conteúdo; limitação de tentativas de login (no máximo cinco por minuto por endereço IP); recusa de mutações vindas de outra origem; segredos apenas em variáveis de ambiente; cabeçalhos de segurança HTTP; alinhamento ao OWASP Top 10 \cite{owasptop10}. & Segurança & Essencial\\[1pt]
\reqtarget{RNF04} & \textbf{Usabilidade.} Interface em português do Brasil; fluxo de criar e publicar um documento concluível por usuário sem formação técnica e sem consultar documentação. Meta: pontuação SUS $\geq$ 68. & Usabilidade & Importante\\[1pt]
\reqtarget{RNF05} & \textbf{Acessibilidade.} Conformidade com o WCAG 2.1 nível AA \cite{wcag21} nas páginas públicas e no painel: contraste mínimo de 4,5:1, navegação por teclado, foco visível, rótulos e textos alternativos. & Usabilidade & Importante\\[1pt]
\reqtarget{RNF06} & \textbf{Responsividade.} Layout adaptável de 320 px a telas largas; \en{sidebar} recolhível em telas pequenas. & Usabilidade; Compatibilidade & Importante\\[1pt]
\reqtarget{RNF07} & \textbf{SEO.} HTML pré-renderizado com título, descrição e metadados por página; URLs estáveis e legíveis. & Adequação funcional & Importante\\[1pt]
\reqtarget{RNF08} & \textbf{Manutenibilidade e modularidade.} Fronteiras entre Dados, Core e UI respeitadas; acesso ao banco somente por \arq{src/lib/db.js}; \en{lint} sem erros; documentação sincronizada com o código. & Manutenibilidade & Essencial\\[1pt]
\reqtarget{RNF09} & \textbf{Extensibilidade.} Novos \en{plugins} \texttt{remark}/\texttt{rehype}, tipos de conteúdo e campos de configuração adicionáveis sem alterar a estrutura central. & Manutenibilidade & Desejável\\[1pt]
\reqtarget{RNF10} & \textbf{Integridade e confiabilidade.} Operações compostas (mover, ordenar) em transação; rotina documentada de \en{backup} e restauração; falha na revalidação não pode causar perda de dados. & Confiabilidade & Importante\\[1pt]
\reqtarget{RNF11} & \textbf{Compatibilidade de navegadores.} Duas últimas versões estáveis de Chrome, Firefox, Safari e Edge. & Compatibilidade & Importante\\[1pt]
\reqtarget{RNF12} & \textbf{Internacionalização.} Nesta versão apenas pt-BR; textos de interface centralizados para viabilizar tradução futura. & Adequação funcional & Desejável\\[1pt]
\reqtarget{RNF13} & \textbf{Licenciamento.} Código publicado sob licença de código aberto (D-02), com dependências de licenças compatíveis. & Restrição & Essencial\\[1pt]
\reqtarget{RNF14} & \textbf{Observabilidade mínima.} Erros de servidor registrados em \en{log} com contexto suficiente ao diagnóstico, sem expor dados sensíveis ao usuário. & Manutenibilidade & Desejável\\[1pt]
\reqtarget{RNF15} & \textbf{Privacidade.} Não coletar dados pessoais de leitores além dos registros de acesso do servidor; sem rastreadores de terceiros por padrão; dados do administrador limitados a nome e e-mail. & Segurança & Importante\\[1pt]
\end{longtable}
\fonte{Elaborado pelos autores, com base em \citeonline{iso25010}.}

\section{Regras de Negócio}\label{sec:regras-negocio}

As regras de negócio (RN) constituem invariantes do domínio; devem ser garantidas no servidor, independentemente da interface.

\begin{longtable}{@{}L{1.3cm}L{9.6cm}L{2.8cm}@{}}
\caption{Regras de negócio}\label{tab:rn}\\
\toprule
\textbf{ID} & \textbf{Regra} & \textbf{Relacionados}\\
\midrule
\endfirsthead
\toprule
\textbf{ID} & \textbf{Regra} & \textbf{Relacionados}\\
\midrule
\endhead
\bottomrule
\endfoot
\reqtarget{RN01} & Somente usuários autenticados podem criar, alterar, mover ou excluir conteúdo e alterar a configuração. & RF21; UC08\\[1pt]
\reqtarget{RN02} & Todo conteúdo tem tipo DOCUMENT ou COLLECTION. Apenas COLLECTION pode ter filhos; DOCUMENT é folha. O tipo é validado na aplicação, pois o SQLite não impõe valores de enumeração. & RF02, RF05, RF06\\[1pt]
\reqtarget{RN03} & O \en{slug} contém apenas letras minúsculas, dígitos e hífens (\texttt{\^{}[a-z0-9]+(-[a-z0-9]+)*\$}), com 1 a 80 caracteres; é único entre irmãos; e os \en{slugs} \texttt{admin}, \texttt{api}, \texttt{login} e \arq{_next} são reservados no nível raiz (D-05). & RF05, RF14\\[1pt]
\reqtarget{RN04} & A URL pública de um conteúdo é a concatenação dos \en{slugs} da raiz até o nó, separados por barras. & RF14\\[1pt]
\reqtarget{RN05} & A ordem é um inteiro não negativo; entre irmãos, a listagem ordena por ordem crescente e, em empate, por data de criação. & RF03\\[1pt]
\reqtarget{RN06} & Mover um conteúdo não pode criar ciclo: o destino não pode ser o próprio nó nem um de seus descendentes. & RF02, RF03\\[1pt]
\reqtarget{RN07} & Excluir um conteúdo exclui todos os seus descendentes e exige confirmação explícita no painel. & RF07\\[1pt]
\reqtarget{RN08} & Existe exatamente um registro de configuração; se ausente, é criado com os valores padrão (nome ``Vitrine'', cores \texttt{\#3b82f6}, \texttt{\#1e293b} e \texttt{\#f59e0b}). & RF18, RF19, RF23\\[1pt]
\reqtarget{RN09} & Cores são aceitas no formato \texttt{\#RRGGBB}; o sistema avisa, sem bloquear, quando o contraste de uma cor do tema com o branco é inferior a 4,5:1 (primária e secundária, usadas com texto) ou a 3:1 (destaque, que não é texto); os pares estão na Tabela~\ref{tab:pares-contraste}. & RF19; RNF05\\[1pt]
\reqtarget{RN10} & Toda mutação de conteúdo ou de configuração dispara a revalidação das páginas afetadas: somente a página, se apenas o corpo mudou; o \en{layout} raiz, se título, \en{slug}, pai, ordem, estado de publicação, tipo (documento ou coleção) ou configuração mudaram. & RF17\\[1pt]
\reqtarget{RN11} & O corpo em Markdown nunca é interpretado como HTML bruto. & RNF03\\[1pt]
\reqtarget{RN12} & Sessões têm prazo de expiração: sete dias, renovados no máximo uma vez por dia enquanto a sessão é usada; uma sessão expirada equivale a nenhuma sessão. & RF01\\[1pt]
\reqtarget{RN13} & O autor de um conteúdo é o usuário autenticado que o criou. & RF05\\[1pt]
\reqtarget{RN14} & A abertura da URL de uma COLLECTION publicada redireciona ao primeiro documento descendente publicado, na ordem da árvore; se não houver, exibe um estado vazio (D-16). & RF14\\[1pt]
\reqtarget{RN15} & Um conteúdo é visível ao público --- por URL, na \en{sidebar} e na barra superior, na geração estática e no mapa do site --- somente se ele e todos os seus ancestrais estiverem PUBLISHED. Conteúdo não visível responde 404 por URL direta. O painel administrativo exibe todos os conteúdos, com o seu estado. & RF14, RF25, RF26\\[1pt]
\reqtarget{RN16} & O estado padrão na criação é PUBLISHED. O administrador pode alternar entre DRAFT e PUBLISHED a qualquer momento. Alterar o estado de uma COLLECTION afeta a visibilidade de toda a sua subárvore, sem modificar o estado dos descendentes. Mudar o estado é uma mudança estrutural para a revalidação (RN10). & RF26, RF17\\[1pt]
\reqtarget{RN17} & Os links sociais do rodapé (Discord, GitHub, YouTube, Instagram, LinkedIn e Twitch) são opcionais e independentes. Só são aceitas URLs absolutas com protocolo \texttt{http} ou \texttt{https}; campo vazio equivale a link ausente e o ícone não é exibido; os links são renderizados com os atributos \texttt{noopener} e \texttt{noreferrer}. & RF24, RNF03\\[1pt]
\end{longtable}
\fonte{Elaborado pelos autores.}

\vitrinependencia{Dados}{A regra RN02 (apenas COLLECTION tem filhos) e a unicidade de slug entre irmãos (RN03) ainda não estão garantidas pelo esquema atual; decidir entre restrição no banco e validação na aplicação (decisões D-05 e D-14) e registrar no mural.}

\section{Matriz de Rastreabilidade}\label{sec:matriz-rastreabilidade}

A Tabela~\ref{tab:matriz-rastreabilidade} relaciona cada requisito funcional ao caso de uso e aos módulos de implementação (definidos na Seção~\ref{sec:modulos}). A relação com os casos de teste está no Capítulo~\ref{cap:testes}.

\begin{longtable}{@{}L{1.4cm}L{2.5cm}L{3.8cm}L{3.0cm}@{}}
\caption{Rastreabilidade entre requisitos funcionais, casos de uso e módulos}\label{tab:matriz-rastreabilidade}\\
\toprule
\textbf{RF} & \textbf{Caso(s) de uso} & \textbf{Módulo(s)} & \textbf{Agente(s)}\\
\midrule
\endfirsthead
\toprule
\textbf{RF} & \textbf{Caso(s) de uso} & \textbf{Módulo(s)} & \textbf{Agente(s)}\\
\midrule
\endhead
\bottomrule
\endfoot
RF01 & UC08 & M03 & Core\\
RF02 & UC04, UC07 & M01, M02 & Dados\\
RF03 & UC07 & M01, M02, M10 & Dados, UI\\
RF04 & UC09 & M11 & UI, Core\\
RF05 & UC05 & M04, M09 & Core, UI\\
RF06 & UC06 & M04, M10 & Core, UI\\
RF07 & UC04 & M04, M10 & Core, UI\\
RF08 & UC04 & M09 & UI\\
RF09 & UC03 & M08 & UI\\
RF10 & UC03 & M08 & UI\\
RF11 & UC03 & M08 & UI\\
RF12 & UC01 & M07 & UI\\
RF13 & UC01 & M07 & UI\\
RF14 & UC02 & M02, M07 & Core, Dados\\
RF15 & UC05 & M09 & UI\\
RF16 & UC05 & M04, M09 & Core, UI\\
RF17 & UC12 & M06 & Core\\
RF18 & UC10 & M05, M11 & Core, UI\\
RF19 & UC11 & M05, M11 & UI, Core\\
RF20 & UC08 & M03 & Core\\
RF21 & UC08 & M03, M04, M05 & Core\\
RF22 & UC08 & M01, M03, M13 & Core, Dados\\
RF23 & UC09 & M01 & Dados\\
RF24 & UC01, UC10 & M01, M07, M11 & UI, Dados\\
RF25 & UC02 & M12 & Core\\
RF26 & UC02, UC04, UC05 & M01, M04, M07, M10 & Core, Dados, UI\\
\end{longtable}
\fonte{Elaborado pelos autores.}

\section{Prototipagem de Baixa Fidelidade}\label{sec:prototipagem}

A prototipagem de baixa fidelidade permitiu validar a jornada do usuário e a disposição dos elementos antes da codificação final \cite{reis2026ic}. As quatro interfaces do relatório de IC, reproduzidas a seguir, são a referência visual para a implementação da interface (agente UI) e originaram requisitos funcionais.

\subsection{Visualização do Sistema com Conteúdo}

A Figura~\ref{fig:ic-instancia} mostra uma instância do sistema em operação, preenchida com conteúdos de programação. Ela evidencia a barra superior com as coleções raiz (RF13), a \en{sidebar} hierárquica com níveis recolhíveis (RF12), a renderização de títulos, listas, tabelas e texto tachado (RF11) e o rodapé com ícones de redes sociais (RF24).

\vitrinefigura[tipo=ic, arquivo=ic-fig3-instancia-fromcaio, largura=0.7,
  fonte={Relatório de IC \cite{reis2026ic}, Figura 3.}]
  {Instância do sistema populada com cursos de programação (protótipo da IC)}{fig:ic-instancia}

\subsection{Gestão e Organização de Conteúdos}

O sistema permite uma estrutura hierárquica em que coleções e arquivos são aninhados. A Figura~\ref{fig:ic-gestao} apresenta a interface administrativa de organização, em que o usuário gerencia a árvore e define a precedência de exibição de cada item, com ações de criar, editar e excluir por nó (RF03, RF06 e RF07).

\vitrinefigura[tipo=ic, arquivo=ic-fig4-gestao-conteudo, largura=0.6,
  fonte={Relatório de IC \cite{reis2026ic}, Figura 4.}]
  {Interface de organização de documentos e coleções aninhadas (protótipo da IC)}{fig:ic-gestao}

\subsection{Customização Dinâmica}

A tela de customização (Figura~\ref{fig:ic-personalizacao}) permite ao administrador alterar a identidade visual, incluindo o logotipo e as cores. No protótipo, as cores foram rotuladas por função na barra de navegação (``principal \en{background}'' e ``principal \en{hover}''); o modelo de dados adota três cores semânticas --- primária, secundária e de destaque --- e o mapeamento entre ambos é definido na Seção~\ref{sec:arq-tema}.

\vitrinefigura[tipo=ic, arquivo=ic-fig5-personalizacao, largura=0.55,
  fonte={Relatório de IC \cite{reis2026ic}, Figura 5.}]
  {Interface de personalização de cores e identidade visual (protótipo da IC)}{fig:ic-personalizacao}

\subsection{Editor de Escrita Científica}

A interface do editor (Figura~\ref{fig:ic-editor}) foi priorizada para suportar a escrita acadêmica: Markdown integrado a fórmulas \LaTeX{} e blocos de código com realce, com barra de ferramentas, painel dividido com pré-visualização e dois comandos de publicação --- ``Salvar localmente'' e ``Enviar nova versão'' (RF08, RF15 e RF16).

\vitrinefigura[tipo=ic, arquivo=ic-fig6-editor, largura=0.65,
  fonte={Relatório de IC \cite{reis2026ic}, Figura 6.}]
  {Interface do editor de texto com suporte a elementos técnicos (protótipo da IC)}{fig:ic-editor}

\subsection{Lacunas Identificadas nos Protótipos}

A análise dos protótipos em confronto com o modelo de dados revelou duas lacunas, ambas resolvidas por decisão do Engenheiro Chefe em 30/09/2026: (i) o rodapé apresenta links de redes sociais, mas o esquema original não previa onde armazená-los, o que foi resolvido com seis campos opcionais e independentes em \texttt{Configuration} (\decisao{D-19}, \adr{ADR-013}); e (ii) o botão ``Enviar nova versão'' sugeria um histórico de versões, que o modelo de dados não contempla; adotou-se o estado de publicação \texttt{DRAFT}/\texttt{PUBLISHED}, com padrão \texttt{PUBLISHED} (\decisao{D-15}, \adr{ADR-012}), e o histórico de versões permanece fora do escopo desta versão.

\section{Síntese do Capítulo}

Foram especificados doze casos de uso, vinte e seis requisitos funcionais, quinze requisitos não funcionais e dezessete regras de negócio, com rastreabilidade para módulos e agentes. Os requisitos herdados da IC foram preservados e refinados; os derivados foram marcados. Este capítulo é o contrato de comportamento: o Capítulo~\ref{cap:arquitetura} define \emph{como} o sistema atende a ele.

%% ===================================================================================
%% Capítulo 5 -- Projeto da Arquitetura e da Solução
%% FONTE DA VERDADE da arquitetura, do contrato de API, do roteamento e da estratégia de
%% cache. Os diagramas citados estão em doc/diagrams/src/ (ver doc/diagrams/README.md).
%% Dono: Arquiteto/Redator. Mudanças propostas pelos desenvolvedores entram pelo mural.
%% Convenção: "normativo" = contrato a ser cumprido; "esboço" = ilustração não vinculante.
%% ===================================================================================
\chapter{Projeto da Arquitetura e da Solução}\label{cap:arquitetura}

Este capítulo define \emph{como} o Vitrine atende aos requisitos do Capítulo~\ref{cap:requisitos}. Suas seções são normativas para a implementação, exceto os trechos explicitamente chamados de \emph{esboço}. Quando a implementação precisar divergir, o desenvolvedor propõe a mudança no mural e o Arquiteto atualiza este capítulo antes da integração.

\section{Visão Arquitetural e Princípios}\label{sec:arq-visao}

A arquitetura é um \textbf{monólito modular}: um único repositório, um único processo Node.js e um único arquivo de banco de dados. Os principais direcionadores são a portabilidade e a simplicidade de implantação (\req{RNF01}), o desempenho de leitura (\req{RNF02}), a segurança (\req{RNF03}) e a manutenibilidade por domínios (\req{RNF08}). Esses direcionadores se traduzem nos princípios a seguir.

\begin{enumerate}[label=\textbf{PA\arabic*.}, leftmargin=1.6cm]
  \item \textbf{Monólito modular.} Um processo, um banco, um repositório; módulos separados por domínio (Dados, Core, UI), com dependências unidirecionais (Figura~\ref{fig:arq-camadas}).
  \item \textbf{Pré-renderizar por padrão.} Páginas públicas são estáticas e regeneradas de forma incremental; o dinamismo fica restrito ao painel administrativo e à API.
  \item \textbf{Acesso a dados centralizado.} Somente \arq{src/lib/db.js} instancia o cliente do Prisma; rotas e componentes não escrevem consultas.
  \item \textbf{Segurança em profundidade.} Verificação otimista no \en{Proxy}, verificação de sessão em cada Route Handler e Server Component protegido, e validação de entrada no servidor.
  \item \textbf{Configuração por ambiente.} Segredos e parâmetros de implantação em variáveis de ambiente, nunca no repositório \cite{twelvefactor}.
  \item \textbf{Conteúdo é dado, não código.} O corpo é Markdown; HTML bruto é proibido (RN11).
  \item \textbf{Personalização por dados.} Identidade e cores vêm do banco e chegam ao CSS por variáveis, sem novo \en{build}.
  \item \textbf{Simplicidade operacional.} Sem fila, cache externo ou serviço adicional obrigatório; a implantação é um contêiner Docker com um volume para o banco (\adr{ADR-011}).
\end{enumerate}

\section{Decisões Arquiteturais}\label{sec:arq-decisoes}

As decisões estruturantes estão registradas como ADR no Apêndice~\ref{ap:adr}. A Tabela~\ref{tab:resumo-adr} as resume. As ADR-011 a ADR-014 foram acrescentadas em 30/09/2026, após o merge M-001 e a atualização P-001, as ADR-015 a ADR-017 em 01/10/2026, quando o Engenheiro Chefe encerrou as decisões que restavam, as ADR-018 e ADR-019 em 01/10/2026, após o merge M-004, e a ADR-020 em 01/10/2026, após o merge M-005; os identificadores D-xx das decisões iniciais continuam citados no texto e levam à ADR que as registra.

\begingroup\small\renewcommand{\arraystretch}{1.25}
\begin{longtable}{@{}L{2.0cm}L{5.6cm}L{6.4cm}@{}}
\caption{Resumo das decisões arquiteturais aceitas}\label{tab:resumo-adr}\\
\toprule
\textbf{ADR} & \textbf{Decisão} & \textbf{Motivo principal}\\
\midrule
\endfirsthead
\toprule
\textbf{ADR} & \textbf{Decisão} & \textbf{Motivo principal}\\
\midrule
\endhead
\bottomrule
\endfoot
\adr{ADR-001} & Monólito modular com Next.js (App Router) & Simplicidade operacional; uma unidade de implantação\\
\adr{ADR-002} & SQLite como banco embutido & Portabilidade; arquivo único; custo zero\\
\adr{ADR-003} & Prisma ORM 7 como camada de acesso & Modelo declarativo, migrações e cliente tipado\\
\adr{ADR-004} & better-auth no lugar do Auth.js previsto na IC & E-mail e senha com sessão em banco e adaptador Prisma\\
\adr{ADR-005} & ISR com revalidação sob demanda & Desempenho estático com atualização imediata\\
\adr{ADR-006} & Markdown com remark/rehype e KaTeX & Escrita científica nativa, sem HTML bruto\\
\adr{ADR-007} & Tematização por variáveis CSS e Tailwind 4 & Personalização sem \en{build}\\
\adr{ADR-008} & JavaScript com JSDoc, sem TypeScript & Herdado do projeto base; menor fricção\\
\adr{ADR-009} & Documentation-Driven Development com \LaTeX{} e diagramas como código & Coerência entre agentes\\
\adr{ADR-010} & Fluxo Git por domínios com mural na \texttt{main} & Coordenação assíncrona rastreável\\
\adr{ADR-011} & Implantação por imagem Docker em VPS tradicional, com volume para o SQLite & Persistência do arquivo de banco e ambiente reprodutível (resolve D-04)\\
\adr{ADR-012} & Estado de publicação \texttt{DRAFT}/\texttt{PUBLISHED}, padrão \texttt{PUBLISHED} & Rascunhos sem histórico de versões (resolve D-15)\\
\adr{ADR-013} & Seis campos opcionais de redes sociais em \texttt{Configuration} & Rodapé configurável sem tabela adicional (resolve D-19)\\
\adr{ADR-014} & Adaptador \texttt{better-sqlite3} e conformidade com o Prisma 7 & O Prisma 7 exige \en{driver adapter} para o SQLite (resolve D-09)\\
\adr{ADR-015} & Nome Vitrine e licença MIT & Identidade e licença permissiva do código aberto (resolve D-01 e D-02)\\
\adr{ADR-016} & Documento de Arquitetura e Requisitos de Software, sem formato acadêmico & Documentação a serviço do produto (resolve D-18)\\
\adr{ADR-017} & Decisões técnicas com valor padrão adotado & Encerra as questões técnicas iniciais, refináveis por RFC (D-03, D-05 a D-08, D-10 a D-14, D-16, D-17)\\
\adr{ADR-018} & Administrador por \texttt{npm run admin:create} e contrato da API (envelope \texttt{data}, erros de entrada e de origem) & Aprova as nove propostas do Core (RFC P-002; fecha o D-03)\\
\adr{ADR-019} & Manter o Prisma 7 apesar de quatro alertas do \texttt{npm audit} & A correção automática rebaixaria o Prisma para a versão 6\\
\adr{ADR-020} & Texto sobre a cor primária escolhido por \texttt{readableTextOn} e indicador de foco neutro & O branco fixo reprova o nível AA com a paleta padrão, e o foco não pode depender de uma cor configurável\\
\end{longtable}
\fonte{Elaborado pelos autores.}
\endgroup

\section{Visões da Arquitetura}\label{sec:arq-visoes}

\subsection{Contexto}

O sistema atende dois atores humanos --- Leitor e Administrador --- e depende apenas de uma VPS tradicional com Docker (\decisao{D-04}). O Implantador obtém o código de um repositório Git, constrói a imagem Docker e a executa na VPS (Figura~\ref{fig:arq-contexto}).

\vitrinefigura[tipo=diagrama, arquivo=arq-contexto, agente=Arquiteto, largura=0.9,
  origem={doc/diagrams/src/arq-contexto.puml},
  descricao={Diagrama de contexto (C4 nível 1): Leitor, Administrador, Implantador, o sistema Vitrine, a VPS tradicional com Docker e o repositório Git.}]
  {Diagrama de contexto do Vitrine}{fig:arq-contexto}

\subsection{Contêineres}

Na terminologia do modelo C4, há um único \emph{contêiner} de aplicação --- o servidor Next.js, que na implantação é empacotado em uma imagem Docker (\adr{ADR-011}) --- mais o navegador do usuário e o arquivo SQLite. Dentro do servidor convivem as páginas públicas (Server Components com ISR), o painel administrativo, os Route Handlers, o better-auth, o Proxy, a camada de dados e o cache de ISR (Figura~\ref{fig:arq-conteineres}).

\vitrinefigura[tipo=diagrama, arquivo=arq-conteineres, agente=Arquiteto, largura=0.95,
  origem={doc/diagrams/src/arq-conteineres.puml},
  descricao={Diagrama de contêineres (C4 nível 2): navegador, servidor Node.js em contêiner Docker com páginas públicas, painel, Route Handlers, better-auth, Proxy, camada de dados (Prisma e adaptador better-sqlite3) e cache ISR, e o banco SQLite em volume.}]
  {Diagrama de contêineres do Vitrine}{fig:arq-conteineres}

\subsection{Componentes e Regras de Dependência}

A Figura~\ref{fig:arq-componentes} mostra os componentes por diretório e o agente responsável por cada um, e a Figura~\ref{fig:arq-camadas} resume as camadas. As regras de dependência são normativas:

\begin{itemize}
  \item componentes em \arq{src/components} \textbf{não} importam \arq{src/lib/db.js} nem o cliente gerado do Prisma (\arq{@/generated/prisma}); recebem dados por \en{props} dos Server Components ou os buscam nas rotas \arq{/api};
  \item rotas em \arq{src/app/api} \textbf{não} escrevem consultas Prisma; chamam funções de \arq{db.js};
  \item apenas \arq{db.js} instancia o cliente do Prisma na aplicação, com o adaptador \arq{better-sqlite3} (uma instância por processo, \adr{ADR-014}); a única exceção é o \en{script} \arq{prisma/seed.js}, executado pela CLI fora do Next.js;
  \item toda rota de escrita valida a sessão e a entrada antes de chamar a camada de dados.
\end{itemize}

\vitrinefigura[tipo=diagrama, arquivo=arq-componentes, agente=Arquiteto, largura=0.95,
  origem={doc/diagrams/src/arq-componentes.puml},
  descricao={Diagrama de componentes por diretório (src/app, src/components, src/lib, prisma), coloridos por agente responsável, com as regras de dependência.}]
  {Componentes do monólito e regras de dependência}{fig:arq-componentes}

\vitrinefigura[tipo=diagrama, arquivo=arq-camadas, agente=Arquiteto, largura=0.85,
  origem={doc/diagrams/src/arq-camadas.puml},
  descricao={Camadas Apresentação, Aplicação, Dados e Infraestrutura, com o sentido permitido das dependências.}]
  {Camadas e sentido das dependências}{fig:arq-camadas}

\section{Estrutura do Repositório e Convenções}\label{sec:arq-estrutura}

O Código~\ref{lst:arvore-atual} mostra a estrutura verificada no repositório na data desta redação, após os merges M-001 a M-005: o \en{scaffold} do Next.js, o esquema do Prisma 7 com a migração inicial e o \en{seed}, as funções de dados, a autenticação, as rotas de API de conteúdo e de configuração, o comando do administrador, o tema e os primeiros componentes da UI (\arq{src/components/admin/} e \arq{src/components/ui/}), os testes e a governança. O Código~\ref{lst:arvore-alvo} mostra a estrutura-alvo, que os agentes constroem por módulo, já com as pastas criadas pelos agentes fora do desenho inicial: \arq{scripts/}, \arq{tests/} (\arq{setup/}, \arq{dados/}, \arq{core/} e \arq{ui/}), \arq{src/app/api/_lib/} e \arq{src/app/admin/_lib/}. Os diretórios \arq{_lib} começam por sublinhado e, no App Router do Next.js, são \emph{pastas privadas}: não geram rota e guardam o código compartilhado pelos Route Handlers e pelas páginas da mesma área. O alias \arq{@/} aponta para \arq{src/} (definido em \arq{jsconfig.json}).

\begin{lstlisting}[caption={Estrutura atual do repositório (verificada)}, label={lst:arvore-atual}]
.
|-- CLAUDE.md  ARQUITETO.md  DESENVOLVEDORES.md  MURAL_DE_AGENTES.md  LICENSE
|-- doc/                      # documentação (ref/, tex/, diagrams/)
|-- prisma/                   # schema.prisma, migrations/ (init) e seed.js (Dados)
|-- prisma.config.ts          # configuração da CLI (datasource.url vem de DATABASE_URL)
|-- scripts/create-admin.js   # npm run admin:create (Core)
|-- tests/                    # setup/ (comum), dados/, core/ e ui/ (Vitest)
|-- vitest.config.mjs         # configuração do Vitest (alias @, banco temporário)
|-- .env.example              # nomes das variáveis de ambiente, sem valores
|-- public/                   # ativos estáticos do scaffold
|-- src/proxy.js              # verificação otimista de sessão em /admin (Core)
|-- src/app/                  # layout.js e page.js (scaffold); globals.css (tema, UI)
|   |-- login/                # page.js (cartão de login; usa o LoginForm)
|   |-- admin/                # layout.js, page.js e _lib/ (session, actions)
|   `-- api/                  # auth/[...all], contents/, configuration/ e _lib/
|-- src/components/           # admin/ (painel, LoginForm) e ui/ (Button, icons)
|-- src/lib/                  # db.js, content-tree.js (Dados); auth.js, auth-client.js,
|                             #   validators.js, revalidate.js (Core); theme.js (UI)
|-- src/generated/prisma/     # cliente gerado; ignorado pelo Git (ADR-014)
|-- eslint.config.mjs  jsconfig.json  next.config.mjs  postcss.config.mjs
|-- package.json  package-lock.json  skills-lock.json
`-- .agents/ .claude/ .windsurf/   # skills do Prisma para os agentes
\end{lstlisting}

\begin{lstlisting}[caption={Estrutura-alvo do código (normativa)}, label={lst:arvore-alvo}]
src/
|-- proxy.js                      # verificação otimista de sessão (Core)
|-- app/
|   |-- layout.js                 # RootLayout: fontes, lang="pt-BR", tema (Core+UI)
|   |-- globals.css               # Tailwind 4 + tokens de tema (UI)
|   |-- page.js                   # página inicial
|   |-- [...slug]/page.js         # leitura de documento/coleção (ISR)
|   |-- not-found.js
|   |-- login/page.js             # cartão de login; usa o LoginForm (UI)
|   |-- admin/                    # painel (dinâmico): layout, page, conteudo/, configuracoes/
|   |   `-- _lib/                 # pasta privada (Core)
|   |       |-- session.js        # requireAdminSession
|   |       `-- actions.js        # Server Function signOut ("use server")
|   `-- api/
|       |-- _lib/                 # http.js (envelope, erros, sessão, origem) e contents.js
|       |-- auth/[...all]/route.js
|       |-- contents/route.js
|       |-- contents/[id]/route.js
|       |-- contents/reorder/route.js
|       `-- configuration/route.js
|-- components/                   # layout/, markdown/, editor/, admin/, ui/ (UI)
`-- lib/
    |-- db.js                     # único ponto de acesso ao Prisma (Dados)
    |-- content-tree.js           # montagem da árvore e resolução de caminhos (Dados)
    |-- auth.js  auth-client.js   # better-auth servidor e cliente (Core)
    |-- validators.js             # esquemas de entrada (Core)
    |-- revalidate.js             # regras de revalidação (Core)
    `-- theme.js                  # variáveis CSS e contraste (UI)
src/generated/prisma/             # cliente gerado por "prisma generate" (não versionado)
prisma/                           # schema.prisma, migrations/, seed.js (Dados)
prisma.config.ts                  # CLI do Prisma: schema, migrations, seed, DATABASE_URL
scripts/create-admin.js           # npm run admin:create: administrador inicial (Core)
tests/                            # setup/ (comum), dados/, core/ e ui/ (Vitest)
vitest.config.mjs                 # configuração dos testes
Dockerfile  docker-compose.yml  .dockerignore   # implantação (M13)
.env.example                      # nomes das variáveis de ambiente, sem valores
\end{lstlisting}

Convenções: identificadores em inglês; comentários e textos de interface em português; módulos ES (\texttt{import}/\texttt{export}); funções públicas de \arq{src/lib} documentadas com JSDoc (\adr{ADR-008}); o arquivo \arq{.env.example} lista os nomes das variáveis de ambiente, sem valores (o \arq{.gitignore} ignora \arq{.env*} e traz a exceção \arq{!.env.example}). Código compartilhado por Route Handlers ou páginas de uma mesma área fica em uma pasta privada \arq{_lib} dessa área, e não em \arq{src/lib}, que é reservado aos módulos de uso geral. Scripts de linha de comando ficam em \arq{scripts/} e são executados com o \texttt{tsx}, pois os módulos de \arq{src/} usam o alias \arq{@/} e o cliente gerado do Prisma é TypeScript. O cliente do Prisma é gerado em \arq{src/generated/prisma}, diretório ignorado pelo Git, e importado de \arq{@/generated/prisma/client} somente por \arq{db.js}; por isso \texttt{prisma generate} deve rodar depois de \texttt{npm install} e antes de \texttt{next build} (\adr{ADR-014}).

\section{Modelo de Dados}\label{sec:arq-dados}

\subsection{Modelo de Domínio}

O domínio possui seis entidades (Figura~\ref{fig:classes-dominio}): quatro de autenticação, gerenciadas pelo better-auth (\texttt{User}, \texttt{Session}, \texttt{Account}, \texttt{Verification}), e duas de negócio do Vitrine (\texttt{Content} e \texttt{Configuration}), além do enumerado \texttt{ContentStatus}. \texttt{Content} possui um auto-relacionamento que forma a árvore de documentos e coleções, um estado de publicação (\texttt{DRAFT} ou \texttt{PUBLISHED}, padrão \texttt{PUBLISHED}; \adr{ADR-012}) e se relaciona com o autor (\texttt{User}). \texttt{Configuration} é um registro único e guarda também os seis links sociais do rodapé, todos opcionais (\adr{ADR-013}).

\vitrinefigura[tipo=diagrama, arquivo=classes-dominio, agente=Arquiteto, largura=0.95,
  origem={doc/diagrams/src/classes-dominio.puml},
  descricao={Diagrama de classes do domínio espelhando prisma/schema.prisma: User, Session, Account, Verification, Content (auto-relacionamento pai e filhos, com status), Configuration (com os seis links sociais), o enumerado ContentStatus e o valor lógico ContentType.}]
  {Modelo de domínio do Vitrine}{fig:classes-dominio}

\subsection{Modelo Físico}

O modelo físico (Figura~\ref{fig:er-vitrine}) é implementado por \arq{prisma/schema.prisma}. O dicionário de dados completo está no Apêndice~\ref{ap:dicionario} e o esquema no Apêndice~\ref{ap:schema}. As convenções adotadas são: identificadores \texttt{cuid}; nomes de tabela pela diretiva \texttt{@@map} (\texttt{user}, \texttt{session}, \texttt{account}, \texttt{verification}, \texttt{contents}, \arq{configurations}); nomes de coluna iguais aos campos do modelo (\emph{camelCase}); carimbos \texttt{createdAt} e \texttt{updatedAt}; e exclusão em cascata nas relações, exceto o autor do conteúdo (\texttt{Content.author}), que usa \texttt{Restrict}: excluir um usuário que tenha conteúdo é bloqueado (RNF10). O enumerado \texttt{ContentStatus} é gravado como texto; o SQLite não impõe o conjunto de valores, de modo que a verificação cabe ao Prisma e à aplicação. Os seis links sociais são colunas de texto opcionais da tabela \arq{configurations}: \arq{discordUrl}, \arq{githubUrl}, \arq{youtubeUrl}, \arq{instagramUrl}, \arq{linkedinUrl} e \arq{twitchUrl}.

\vitrinefigura[tipo=diagrama, arquivo=er-vitrine, agente=Arquiteto, largura=0.95,
  origem={doc/diagrams/src/er-vitrine.dbml},
  descricao={Diagrama entidade-relacionamento gerado no dbdiagram.io a partir de er-vitrine.dbml: tabelas user, session, account, verification, contents (auto-relacionamento e coluna status do enumerado ContentStatus) e configurations (com as seis colunas de links sociais).}]
  {Diagrama entidade-relacionamento do Vitrine}{fig:er-vitrine}

A Figura~\ref{fig:ic-er} reproduz o modelo do relatório de IC, com três tabelas. A evolução consiste em substituir a tabela genérica \texttt{users} pelas tabelas do better-auth e em explicitar as relações de cascata.

\vitrinefigura[tipo=ic, arquivo=ic-fig2-er, largura=0.35,
  fonte={Relatório de IC \cite{reis2026ic}, Figura 2.}]
  {Diagrama entidade-relacionamento original, do relatório de IC}{fig:ic-er}

\subsection{Hierarquia de Conteúdo}\label{sec:arq-hierarquia}

A árvore de conteúdos é representada pelo campo \texttt{parentId} (nulo para nós raiz). A URL pública é a cadeia de \en{slugs} da raiz ao nó (RN04). Para a navegação, a árvore é carregada por \textbf{uma única consulta}, sem o corpo dos documentos, e montada em memória em tempo linear; a resolução de um caminho percorre os níveis comparando \en{slugs}. O Código~\ref{lst:algoritmo-arvore} é um \emph{esboço} do algoritmo a ser implementado em \arq{src/lib/content-tree.js}. O parâmetro \texttt{onlyPublished} implementa a RN15: um nó \texttt{DRAFT} é removido junto com toda a sua subárvore, de modo que um descendente publicado de um ancestral em rascunho também não aparece. A área pública usa a árvore podada; o painel usa a árvore completa.

\begin{lstlisting}[language=VJS, caption={Esboço: montagem da árvore e resolução de caminho}, label={lst:algoritmo-arvore}]
// rows: [{ id, title, slug, type, status, order, parentId, createdAt }] (sem o corpo)
export function buildTree(rows, { onlyPublished = false } = {}) {
  const byId = new Map(rows.map((r) => [r.id, { ...r, children: [] }]));
  const roots = [];
  for (const node of byId.values()) {
    const parent = node.parentId ? byId.get(node.parentId) : null;
    (parent ? parent.children : roots).push(node);
  }
  // RN15: nó não publicado oculta toda a sua subárvore
  const prune = (list) =>
    list
      .filter((n) => !onlyPublished || n.status === "PUBLISHED")
      .map((n) => ({ ...n, children: prune(n.children) }));
  const sortLevel = (list) => {
    list.sort((a, b) => a.order - b.order || a.createdAt - b.createdAt); // RN05
    list.forEach((n) => sortLevel(n.children));
  };
  const visible = prune(roots);
  sortLevel(visible);
  return visible;
}

export function resolvePath(roots, slugs) {
  let level = roots;
  let node = null;
  for (const slug of slugs) {
    node = level.find((n) => n.slug === slug) ?? null;
    if (!node) return null; // 404 (RF14, RN15)
    level = node.children;
  }
  return node;
}
\end{lstlisting}

A alternativa de gravar o caminho materializado em cada nó foi descartada nesta versão, pois exigiria atualizar todos os descendentes a cada renomeação ou movimentação; fica registrada como evolução possível, para o caso de o volume de conteúdo o justificar (\decisao{D-06}, \adr{ADR-017}).

\subsection{Observações sobre o Esquema Atual}\label{sec:arq-observacoes-schema}

A revisão do esquema (\arq{prisma/schema.prisma}) contra esta especificação resultou nas observações da Tabela~\ref{tab:obs-schema}, que foram entradas para o agente Dados. O merge M-001, aprovado em 30/09/2026, integrou a conformidade com o Prisma 7 (\adr{ADR-014}), o enumerado \texttt{ContentStatus} com a coluna \texttt{status} (\adr{ADR-012}) e os seis links sociais de \texttt{Configuration} (\adr{ADR-013}). O merge M-002, aprovado em 01/10/2026, resolveu as observações 1, 3, 4 e 5 na migração inicial e manteve a 2 como validação na aplicação; o merge M-003 aplicou a observação 8 na camada de dados (filtro \texttt{PUBLISHED} nas leituras públicas, Seção~\ref{sec:arq-visibilidade}).

\begin{longtable}{@{}L{0.8cm}L{5.6cm}L{2.2cm}L{4.8cm}@{}}
\caption{Observações sobre o esquema atual do Prisma}\label{tab:obs-schema}\\
\toprule
\textbf{\#} & \textbf{Observação} & \textbf{Regra ou RNF} & \textbf{Recomendação}\\
\midrule
\endfirsthead
\toprule
\textbf{\#} & \textbf{Observação} & \textbf{Regra ou RNF} & \textbf{Recomendação}\\
\midrule
\endhead
\bottomrule
\endfoot
1 & \texttt{slug} sem restrição de unicidade & RN03 & \emph{Resolvido (M-002)}: índice único (\texttt{parentId}, \texttt{slug}) mais índice parcial para a raiz, em SQL (Seção~\ref{sec:arq-indice-parcial}; D-05)\\[1pt]
2 & \texttt{type} é texto livre, ao passo que \texttt{status} já é um enumerado & RN02 & \emph{Mantido como texto}: DOCUMENT e COLLECTION são validados na aplicação (Zod, Core), pois o SQLite não impõe enumerados e um \texttt{enum} do Prisma só duplicaria a regra\\[1pt]
3 & Sem índice em \texttt{parentId} & RNF02 & \emph{Resolvido (M-002)}: índice (\texttt{parentId}, \texttt{order}) para a consulta da árvore\\[1pt]
4 & \arq{Content.author} com \texttt{onDelete: Cascade}: excluir o usuário apaga todo o seu conteúdo & RNF10 & \emph{Resolvido (M-002)}: \texttt{onDelete: Restrict}; excluir um usuário que tenha conteúdo é bloqueado pelo banco\\[1pt]
5 & \texttt{Configuration} sem garantia de registro único & RN08 & \emph{Resolvido (M-002)}: identificador fixo \texttt{default} (padrão do esquema) e acesso por \emph{upsert}\\[1pt]
6 & Datas e booleanos sem valor padrão em \arq{User.emailVerified} & --- & Manter: o valor é fornecido pelo better-auth\\[1pt]
7 & Tabela \texttt{verification} sem relação com \texttt{user} & --- & Manter: desenho do better-auth\\[1pt]
8 & \texttt{status} tem padrão \texttt{PUBLISHED}: um conteúdo recém-criado torna-se público de imediato (decisão D-15) & RN15, RN16, risco R11 & Manter o padrão; a API deve aceitar a criação já como \texttt{DRAFT}, a interface deve exibir o estado com destaque e a camada de dados deve filtrar por \texttt{PUBLISHED} nas leituras públicas\\[1pt]
\end{longtable}
\fonte{Elaborado pelos autores.}

\subsection{Índice Parcial da Raiz}\label{sec:arq-indice-parcial}

A regra RN03 exige \en{slug} único entre irmãos. O índice único composto (\texttt{parentId}, \texttt{slug}) resolve a regra para todos os nós com pai, mas não para a raiz: no SQLite, \texttt{NULL} é distinto de \texttt{NULL}, de modo que dois conteúdos de raiz (\texttt{parentId} nulo) com o mesmo \en{slug} passariam pelo índice composto. O agente Dados resolveu a lacuna com um \textbf{índice único parcial} escrito à mão no SQL da migração inicial (Código~\ref{lst:indice-parcial}), que cobre apenas as linhas da raiz. O Prisma não modela índices parciais, mas também não os trata como divergência (\en{drift}): \texttt{migrate diff} contra as migrações e contra o banco não acusa diferenças. O resultado é a garantia de unicidade no próprio banco, sem depender só da validação na aplicação, que permanece como primeira barreira (\decisao{D-05}, \adr{ADR-017}).

\begin{lstlisting}[caption={Índice parcial da raiz (trecho da migração inicial)}, label={lst:indice-parcial}]
CREATE UNIQUE INDEX "contents_parentId_slug_key"
  ON "contents"("parentId", "slug");

-- manual: cobre somente a raiz
CREATE UNIQUE INDEX "contents_root_slug_key"
  ON "contents"("slug") WHERE "parentId" IS NULL;
\end{lstlisting}

Como a migração é SQL versionado, qualquer nova migração gerada pelo Prisma deve preservar esse índice; o agente Dados deve conferir isso ao alterar o esquema. Em banco descartável, o agente Dados verificou que um \en{slug} repetido na raiz e entre irmãos é recusado e que o mesmo \en{slug} em pais diferentes é aceito; os testes de integração correspondentes (casos CT08 e CT28) ainda devem ser automatizados.

\subsection{Migrações e Dados Iniciais}

As alterações de esquema são versionadas por migrações do Prisma (\arq{prisma/migrations}), geradas em desenvolvimento e aplicadas em produção por um comando não interativo (Apêndice~\ref{ap:implantacao}). O \en{seed} insere a configuração padrão (RF23) e cria, sob controle, o primeiro administrador (RF22). A primeira migração (\arq{20261001034824_init}), o \en{seed} idempotente da configuração padrão (\arq{prisma/seed.js}, executado por \texttt{tsx} porque o cliente gerado é TypeScript) e a instância do cliente em \arq{src/lib/db.js} foram integrados pelo merge M-002; as funções de acesso a dados foram integradas pelo merge M-003; o administrador inicial não passa pelo \en{seed}: é criado pelo comando \texttt{npm run admin:create} do Core, integrado no merge M-004 (\decisao{D-03}, \adr{ADR-018}). Como \arq{prisma} e \arq{dotenv} estão em \texttt{dependencies}, a imagem de produção aplica as migrações com \texttt{prisma migrate deploy}. No Prisma 7, o comando \texttt{migrate dev} não gera o cliente nem executa o \en{seed} automaticamente, e o \en{seed} é declarado em \arq{prisma.config.ts} (chave \texttt{migrations.seed}), não mais no \arq{package.json}; por isso \texttt{prisma generate} e \texttt{prisma db seed} são passos explícitos. Um cuidado adicional: um caminho relativo em \arq{DATABASE_URL} pode ser resolvido de forma diferente pela CLI e pelo adaptador em tempo de execução, criando silenciosamente um segundo banco; o agente Dados verificou que, com \texttt{file:./dev.db} e execução a partir da raiz, a CLI e o adaptador abrem o mesmo arquivo; o adaptador resolve o caminho a partir do diretório de trabalho do processo, e em produção o caminho deve ser absoluto (por exemplo, \texttt{file:/data/vitrine.db}). As instruções de CLI da versão 7 estão nas \en{skills} do Prisma instaladas em \arq{.claude/skills/}.

\section{Autenticação e Autorização}\label{sec:arq-auth}

\subsection{Configuração}

O better-auth 1.7.6 é configurado em \arq{src/lib/auth.js} (merge M-004) com o adaptador do Prisma para SQLite, autenticação por e-mail e senha, sessão persistida no banco, cadastro público desabilitado e o \en{plugin} de integração de \en{cookies} para o Next.js. As rotas são expostas por um único Route Handler em \arq{src/app/api/auth/[...all]/route.js}. A sessão no servidor é obtida por \texttt{getSession(headers)}, usada pelos Route Handlers e pelas páginas do painel. O cliente React (\arq{src/lib/auth-client.js}) usa a mesma origem da página e dispensa variáveis \texttt{NEXT\_PUBLIC\_*}. Variáveis de ambiente: \arq{DATABASE_URL}, \arq{BETTER_AUTH_SECRET} e \arq{BETTER_AUTH_URL} (nomes confirmados na versão instalada, que as lê sozinha) e, lidas somente pelo comando do administrador inicial, \arq{ADMIN_EMAIL}, \arq{ADMIN_PASSWORD} e \arq{ADMIN_NAME}. A Tabela~\ref{tab:auth-config} resume a configuração efetiva, aprovada na \adr{ADR-018}.

\begingroup\small
\begin{longtable}{@{}L{3.0cm}L{9.6cm}L{1.9cm}@{}}
\caption{Configuração efetiva do better-auth}\label{tab:auth-config}\\
\toprule
\textbf{Aspecto} & \textbf{Configuração} & \textbf{Origem}\\
\midrule
\endfirsthead
\toprule
\textbf{Aspecto} & \textbf{Configuração} & \textbf{Origem}\\
\midrule
\endhead
\bottomrule
\endfoot
Cadastro & Desabilitado (\texttt{disableSignUp}); \texttt{POST /api/auth/sign-up/email} responde 400 & RF22\\
\addlinespace
Senha & De 8 a 128 caracteres; o \en{hash} é do próprio better-auth, o mesmo no login e no comando do administrador & RNF03\\
\addlinespace
Sessão & Persistida no banco; expira em 7 dias e, em uso, é renovada no máximo uma vez por dia & RN12\\
\addlinespace
\en{Cookie} & \arq{better-auth.session_token}, \texttt{HttpOnly} e \texttt{SameSite=Lax}; com \arq{BETTER_AUTH_URL} em \texttt{https}, ganha \texttt{Secure} e o prefixo \texttt{\_\_Secure-} & RNF03\\
\addlinespace
Limitação de tentativas & Ligada em todos os ambientes, com contadores em memória por endereço IP (cabeçalho \texttt{x-forwarded-for}) e por rota: 100 requisições por 60 s no geral e \textbf{5 por 60 s em \texttt{/sign-in/email}} (resposta 429) & RNF03\\
\addlinespace
Telemetria & Desligada & RNF15\\
\addlinespace
Segredo & \arq{BETTER_AUTH_SECRET} obrigatório em produção; sem ele, as rotas de autenticação e de API respondem 500 (falha fechada: nenhuma rota protegida é servida) e o \en{build} registra o erro do segredo padrão sem falhar & RNF03\\
\end{longtable}
\fonte{Elaborado pelos autores, a partir de \arq{src/lib/auth.js} e da verificação do Arquiteto no merge M-004.}
\endgroup

\begin{lstlisting}[language=VJS, caption={Configuração efetiva do servidor de autenticação (resumo de src/lib/auth.js)}, label={lst:auth-esboco}]
export const auth = betterAuth({
  appName: "Vitrine",
  database: prismaAdapter(prisma, { provider: "sqlite" }),
  emailAndPassword: {
    enabled: true,
    disableSignUp: true,            // RF22
    minPasswordLength: 8,
    maxPasswordLength: 128,
  },
  session: { expiresIn: 7 * 86400, updateAge: 86400 },   // RN12
  rateLimit: {
    enabled: true, window: 60, max: 100,
    customRules: { "/sign-in/email": { window: 60, max: 5 } },
  },
  telemetry: { enabled: false },    // RNF15
  plugins: [nextCookies()],         // deve ser o último
});
// src/app/api/auth/[...all]/route.js
export const { GET, POST } = toNextJsHandler(auth);
\end{lstlisting}

A limitação de tentativas identifica o cliente pelo cabeçalho \texttt{x-forwarded-for}, que o better-auth 1.7.6 só aceita quando contém um \emph{único} endereço (um valor com vários endereços pode ter sido forjado pelo cliente). Quando o endereço não pode ser determinado, todos os clientes passam a compartilhar um único contador por rota --- por exemplo, apenas cinco tentativas de login por minuto no total ---, o que mantém a proteção, mas deixa um atacante bloquear momentaneamente o login do administrador. Por isso o \en{proxy} reverso deve \emph{definir} esse cabeçalho com o endereço real do cliente, substituindo o valor recebido, e não acrescentar-lhe um segundo endereço (Apêndice~\ref{ap:implantacao}).

\subsection{Fluxo e Proteção em Camadas}

O fluxo de login e de acesso à área administrativa está na Figura~\ref{fig:seq-autenticacao}. A proteção é feita em camadas, porque o Proxy é, por documentação do Next.js, uma verificação otimista e não uma solução completa de autorização \cite{nextjsdocs}:

\begin{table}[htb]
\caption{Camadas de proteção das rotas administrativas}\label{tab:camadas-protecao}
\small
\begin{tabularx}{\linewidth}{@{}L{3.3cm}L{5.2cm}Y@{}}
\toprule
\textbf{Camada} & \textbf{Mecanismo} & \textbf{Finalidade}\\
\midrule
1. Proxy & \arq{src/proxy.js} (casador \arq{/admin/:path*}) lê o \en{cookie} de sessão, sem consultar o banco, e redireciona para \arq{/login?redirectTo=...} se ausente & Melhor experiência: evita renderizar a página para quem não tem sessão\\
\addlinespace
2. Layout/página do painel & Server Component valida a sessão no servidor e redireciona se inválida; a página repete a validação, pois layouts não são renderizados de novo a cada navegação & Barreira real para páginas\\
\addlinespace
3. Route Handlers & Cada handler valida a sessão no banco antes de qualquer ação (RF21) e, nas mutações, antes disso, a origem da requisição; só então lê o corpo & Barreira real para a API\\
\addlinespace
4. Validação de entrada & Esquemas de entrada rejeitam dados inválidos antes da camada de dados & Integridade e segurança\\
\bottomrule
\end{tabularx}
\fonte{Elaborado pelos autores.}
\end{table}

\textbf{Interface do login e do encerramento da sessão.} O formulário de login é o componente \texttt{LoginForm} (\arq{src/components/admin/LoginForm.js}, Componente Cliente), que chama \texttt{authClient.signIn.email}: em qualquer erro de credencial mostra uma única mensagem genérica (UC08, fluxo A1); no status \texttt{429} avisa que as tentativas foram excedidas; sem resposta do servidor, informa a falha de conexão; e, em caso de sucesso, faz a navegação completa para o destino que o servidor já validou (\texttt{safeRedirectPath}). O encerramento da sessão (RF20) \emph{não} passa pela rota \arq{/api/auth/sign-out}: o botão ``Sair'' do cabeçalho do painel envia um formulário à \en{Server Function} \texttt{signOut}, definida em \arq{src/app/admin/_lib/actions.js} (diretiva \texttt{"use server"}), que chama \texttt{auth.api.signOut} com os cabeçalhos da requisição --- apaga a sessão do banco e remove o \en{cookie} --- e redireciona para \arq{/login}. O layout do painel, que é um arquivo de rota do Core, entrega a função ao cabeçalho por \en{prop}; assim, os componentes da UI não importam o módulo de autenticação do servidor (regras de dependência, Seção~\ref{sec:arq-visoes}). Toda função exportada de um arquivo \texttt{"use server"} é um ponto de entrada HTTP, alcançável por uma requisição direta; por isso esse arquivo só pode conter funções que validem no servidor tudo o que exigem (a \texttt{signOut} não recebe argumentos).

\vitrinefigura[tipo=diagrama, arquivo=seq-autenticacao, agente=Arquiteto, largura=0.95,
  origem={doc/diagrams/src/seq-autenticacao.puml},
  descricao={Sequência de autenticação: acesso a /admin sem cookie, redirecionamento a /login, POST de credenciais ao better-auth, criação de sessão, Set-Cookie e acesso autorizado.}]
  {Sequência de autenticação do administrador}{fig:seq-autenticacao}

\subsection{Política de Cadastro e Papéis}

Na versão inicial há um único perfil, o administrador. O primeiro usuário é criado pelo comando \texttt{npm run admin:create} (\decisao{D-03}, \adr{ADR-018}), executado por quem instala o sistema: ele lê \arq{ADMIN_EMAIL}, \arq{ADMIN_PASSWORD} e, opcionalmente, \arq{ADMIN_NAME}, valida o formato e chama \texttt{ensureInitialAdmin} (\arq{src/lib/auth.js}), que usa o adaptador interno do better-auth --- com o mesmo \en{hash} de senha do login --- porque a rota de cadastro está desabilitada. O comando é \textbf{idempotente}: com o mesmo e-mail, nada muda (nem a senha) e a saída é 0; recusa um segundo usuário com outro e-mail, pois há um único administrador (\decisao{D-13}; saída 1); completa a credencial de um usuário que ficou sem ela por uma execução interrompida; e nunca imprime a senha. Não há página aberta de primeiro acesso, para evitar o sequestro da instalação. As credenciais não passam pelo \arq{prisma/seed.js}, que cria apenas a configuração padrão (RF23); o comando roda com o \texttt{tsx}, como o \en{seed}. Na imagem Docker: \texttt{docker compose run -{}-rm vitrine npm run admin:create} (Apêndice~\ref{ap:implantacao}). O \en{plugin} \texttt{admin} do better-auth, presente no pacote instalado, é o caminho previsto para papéis e permissões futuros (\decisao{D-13}).

\section{Renderização, Cache e ISR}\label{sec:arq-isr}

\subsection{Estratégia por Rota}

A Tabela~\ref{tab:rotas-render} define o modo de renderização de cada conjunto de rotas. O projeto usa a ISR \emph{sem} Cache Components (o recurso \arq{cacheComponents} não está habilitado em \arq{next.config.mjs}), conforme o guia correspondente da documentação do Next.js 16 \cite{nextjsdocs}.

\begin{table}[htb]
\caption{Modo de renderização por conjunto de rotas}\label{tab:rotas-render}
\small
\begin{tabularx}{\linewidth}{@{}L{3.3cm}L{3.0cm}Y@{}}
\toprule
\textbf{Rota} & \textbf{Modo} & \textbf{Regras}\\
\midrule
\arq{/} e \arq{/[...slug]} & Estático com ISR & \arq{generateStaticParams} lista os caminhos \emph{publicados} (RN15); \texttt{revalidate} alto (por exemplo, 3600 s) como rede de segurança; novos caminhos são gerados sob demanda; revalidação sob demanda pelas mutações\\
\addlinespace
\arq{/login} & Estático ou dinâmico & Formulário; sem dados sensíveis\\
\addlinespace
\arq{/admin/**} & Dinâmico & Lê a sessão por \en{cookies}; nunca em cache\\
\addlinespace
\arq{/api/**} & Dinâmico & Route Handlers; respostas sem cache\\
\bottomrule
\end{tabularx}
\fonte{Elaborado pelos autores.}
\end{table}

\subsection{Regras de Revalidação}

A Tabela~\ref{tab:revalidacao} implementa as regras RN10 e RN16. A invalidação é feita em Route Handlers com \arq{revalidatePath}, função que só pode ser chamada no servidor, não em Componentes Cliente nem no Proxy. Em Route Handlers, a função marca o caminho para revalidação, que ocorre na próxima visita \cite{nextjsdocs}. A escolha por \emph{layout} raiz quando a estrutura muda é necessária porque a \en{sidebar} e a barra superior, presentes em todas as páginas, dependem da árvore inteira.

\begin{table}[htb]
\caption{Eventos de mutação e a revalidação correspondente}\label{tab:revalidacao}
\small
\begin{tabularx}{\linewidth}{@{}L{5.2cm}Y@{}}
\toprule
\textbf{Evento} & \textbf{Ação de revalidação}\\
\midrule
Criar ou excluir conteúdo; alterar título, \en{slug}, pai, ordem ou tipo; mover/reordenar &\texttt{revalidatePath("/", "layout")}: invalida todas as páginas, pois a árvore de navegação mudou\\
\addlinespace
Alterar o estado de publicação (\texttt{DRAFT} $\leftrightarrow$ \texttt{PUBLISHED}) de qualquer nó & \texttt{revalidatePath("/", "layout")}: a visibilidade do nó e de seus descendentes altera a árvore de navegação e o conjunto de páginas (RN16)\\
\addlinespace
Alterar somente o corpo de um documento (estado inalterado) & \texttt{revalidatePath(caminho do documento)}\\
\addlinespace
Alterar a configuração (nome, logotipo, cores, links sociais) & \texttt{revalidatePath("/", "layout")}: o rodapé aparece em todas as páginas\\
\bottomrule
\end{tabularx}
\fonte{Elaborado pelos autores, com base em RN10 e RN16.}
\end{table}

\subsection{Visibilidade por Estado de Publicação}\label{sec:arq-visibilidade}

A regra RN15 --- um conteúdo só é público se ele e todos os seus ancestrais estiverem \texttt{PUBLISHED} --- é aplicada em um único lugar, a camada de dados, e nunca nos componentes. Os pontos de aplicação são: (i) a árvore usada pela \en{sidebar}, pela barra superior e por \arq{generateStaticParams}, podada conforme o Código~\ref{lst:algoritmo-arvore}; (ii) a resolução do caminho nos Server Components, que responde 404 (\arq{notFound()}) para um nó ausente da árvore podada, de modo que um rascunho não é acessível nem por URL direta; e (iii) o mapa do site e os metadados (RF25). As rotas \arq{/api/contents} e o painel recebem a árvore completa, com o campo \texttt{status} de cada nó (RF26). Como o estado entra na regeneração do \en{layout} raiz (Tabela~\ref{tab:revalidacao}), a mudança de estado reflete-se em todas as páginas na visita seguinte. Como o padrão do estado é \texttt{PUBLISHED} (\decisao{D-15}), a proteção contra a publicação acidental depende de a interface exibir o estado de forma explícita e de a API aceitar a criação já como rascunho (risco R11).

\vitrinefigura[tipo=diagrama, arquivo=seq-leitura-isr, agente=Arquiteto, largura=0.95,
  origem={doc/diagrams/src/seq-leitura-isr.puml},
  descricao={Sequência de leitura de uma página com os três caminhos do cache: HIT, STALE e MISS, incluindo a regeneração em segundo plano e o cabeçalho x-nextjs-cache.}]
  {Sequência de leitura de uma página com ISR}{fig:seq-leitura-isr}

\vitrinefigura[tipo=diagrama, arquivo=seq-publicacao-conteudo, agente=Arquiteto, largura=0.95,
  origem={doc/diagrams/src/seq-publicacao-conteudo.puml},
  descricao={Sequência de edição e publicação: editor, PATCH na API, validação de sessão e de entrada, atualização no banco e revalidação do caminho ou do layout.}]
  {Sequência de edição e publicação de um documento}{fig:seq-publicacao-conteudo}

\subsection{Limitações e Observabilidade}

A documentação do Next.js 16 registra ressalvas que condicionam a implantação: a ISR exige o \en{runtime} Node.js e não funciona em exportação estática; com mais de uma instância, o cache padrão em sistema de arquivos é por instância e a invalidação só atinge a instância que recebeu a chamada; e o Proxy não é executado em requisições de ISR sob demanda \cite{nextjsdocs}. Por isso, a implantação usa \textbf{uma única instância}, isto é, um único contêiner (o que também convém ao SQLite). Em testes, o cabeçalho \arq{x-nextjs-cache} (\texttt{HIT}, \texttt{STALE}, \texttt{MISS}, \texttt{REVALIDATED}) permite verificar o comportamento do cache (Capítulo~\ref{cap:testes}).

\section{Pipeline de Conteúdo Markdown}\label{sec:arq-markdown}

A renderização de um documento usa o componente \arq{MarkdownRenderer} (\arq{src/components/markdown}), compartilhado pela página pública e pela pré-visualização do editor, para garantir que o autor veja o que o leitor verá. A cadeia de processamento é:

\begin{enumerate}
  \item \arq{react-markdown} converte o Markdown em elementos React, sem HTML bruto (RN11);
  \item \arq{remark-math} reconhece \texttt{\$...\$} e \texttt{\$\$...\$\$} (RF09);
  \item \arq{remark-gfm} acrescenta tabelas, tachado, listas de tarefas e links automáticos (RF11; pacote ainda não instalado, \decisao{D-08});
  \item \arq{rehype-katex} renderiza as fórmulas; a folha de estilos do KaTeX é importada uma vez no \en{layout};
  \item um \en{plugin} de realce de sintaxe colore os blocos de código (RF10; biblioteca a definir, \decisao{D-07}).
\end{enumerate}

Requisitos de segurança e de comportamento: manter a transformação padrão de URLs do \arq{react-markdown}, que restringe protocolos; não adicionar \arq{rehype-raw} sem sanitização; tratar expressões \LaTeX{} inválidas de forma tolerante (mensagem visível, sem derrubar a página); e mapear links relativos internos para as URLs por \en{slug}.

\section{Tematização Dinâmica}\label{sec:arq-tema}

A personalização visual (RF04, RF18, RF19) baseia-se em variáveis CSS alimentadas pelo banco. O mecanismo, ilustrado na Figura~\ref{fig:seq-personalizacao-tema}, tem três partes:

\begin{enumerate}
  \item \textbf{Tokens.} Em \arq{globals.css}, o bloco \texttt{@theme inline} do Tailwind mapeia os tokens de cor (por exemplo, \arq{--color-primary}) para variáveis (por exemplo, \arq{var(--vitrine-primary)}); os utilitários compilados referenciam a variável, não o valor;
  \item \textbf{Injeção.} O \en{layout} raiz, um Server Component, lê a configuração e define as variáveis no atributo \texttt{style} do elemento \texttt{<html>} --- as três cores e a cor do texto sobre a primária, \arq{--vitrine-on-primary}, calculada por \texttt{readableTextOn} (Seção~\ref{sec:arq-contraste}) ---, além do nome do site, do logotipo e do \en{favicon}; os valores são validados como \texttt{\#RRGGBB} antes de persistir, o que impede a injeção de CSS;
  \item \textbf{Invalidação.} Ao salvar a configuração, o Route Handler invalida o \en{layout} raiz; na próxima requisição, as páginas são regeneradas com as novas variáveis, sem \en{build}.
\end{enumerate}

\begin{lstlisting}[language=VCSS, caption={Tokens de tema em globals.css (resumo do arquivo integrado no merge M-005)}, label={lst:tema-css}]
@import "tailwindcss";

:root {
  /* Paleta configurável (contrato do M11): <html style> sobrescreve o padrão */
  --vitrine-primary: #3b82f6;
  --vitrine-secondary: #1e293b;
  --vitrine-accent: #f59e0b;
  /* Texto sobre a primária = readableTextOn(primária); definida no layout raiz */
  --vitrine-on-primary: #171717;
  --vitrine-on-secondary: #ffffff;

  /* Neutros fixos (D-17: sem modo escuro automático) e foco independente */
  --vitrine-background: #ffffff;  --vitrine-foreground: #171717;
  --vitrine-surface: #f5f5f5;     --vitrine-muted: #525252;
  --vitrine-border: #d4d4d4;      --vitrine-border-strong: #737373;
  --vitrine-danger: #b91c1c;
  --vitrine-focus: #171717;       --vitrine-focus-halo: #ffffff;
}

@theme inline {
  --color-primary: var(--vitrine-primary);
  --color-secondary: var(--vitrine-secondary);
  --color-accent: var(--vitrine-accent);
  --color-on-primary: var(--vitrine-on-primary);
  --color-on-secondary: var(--vitrine-on-secondary);
  /* neutros no mesmo padrão: --color-background, -foreground, -surface... */
}

/* D-17: dark: só vale sob a classe .dark, nunca pela preferência do sistema */
@custom-variant dark (&:where(.dark, .dark *));
\end{lstlisting}

O mapeamento semântico é: \textbf{primária} para ações e ênfase (botão principal e item ativo da navegação, com o fundo \texttt{bg-primary} e o texto \texttt{text-on-primary}); \textbf{secundária} para a estrutura (fundo da barra superior e cabeçalhos, com texto branco, \texttt{text-on-secondary}) e para o \en{hover} dos botões; e \textbf{destaque} para marcadores e realce não textual (marcador do item ativo e \en{hover} da navegação). O \textbf{indicador de foco} não usa a cor de destaque: é neutro e global --- contorno de 2~px em \texttt{\#171717} com halo branco (\arq{--vitrine-focus} e \arq{--vitrine-focus-halo}) ---, de modo que fica visível sobre qualquer fundo, inclusive sobre uma paleta que o administrador troque (RNF05; \adr{ADR-020}). Os demais tons da interface --- fundo, texto, superfícies, bordas e mensagens de erro --- são neutros fixos, não configuráveis. Nesta versão não há modo escuro automático: vale somente a paleta escolhida pelo administrador (\decisao{D-17}, \adr{ADR-017}), e a variante \texttt{dark:} do Tailwind só vale sob a classe \texttt{.dark}, nunca pela preferência do sistema; por isso o fundo da página é fixo e branco, o texto sobre a cor secundária é branco, o texto sobre a cor primária é escolhido por \texttt{readableTextOn} (Seção~\ref{sec:arq-contraste}) e só as três cores do tema são verificadas. Os componentes usam apenas os utilitários gerados pelos tokens (por exemplo, \texttt{bg-primary}, \texttt{text-on-primary}, \texttt{bg-surface} e \texttt{text-muted}), nunca as cores fixas do Tailwind (RF19).

\subsection{Avisos de Contraste}\label{sec:arq-contraste}

Ao salvar a configuração, o sistema calcula a razão de contraste da WCAG 2.1 (luminância relativa do sRGB, razão $(L_1 + 0{,}05)/(L_2 + 0{,}05)$) entre cada cor do tema e o branco, conforme a Tabela~\ref{tab:pares-contraste}, e emite um aviso, sem bloquear o salvamento, quando a razão é inferior ao mínimo do par (RN09, \req{RNF05}). A razão é simétrica, e o par é sempre contra o branco: o par ``primária e branco'' cobre a cor primária usada como texto sobre o fundo da página (por exemplo, em links). O texto \emph{sobre} um fundo da cor primária, ao contrário, não é branco fixo: é calculado, como descrito adiante.

\begin{table}[htb]
\caption{Pares de contraste verificados na configuração}\label{tab:pares-contraste}
\small
\begin{tabularx}{\linewidth}{@{}L{2.8cm}L{3.6cm}L{1.7cm}Y@{}}
\toprule
\textbf{Campo avisado} & \textbf{Par (cor e branco)} & \textbf{Mínimo} & \textbf{Uso que a razão protege}\\
\midrule
\texttt{primaryColor} & primária $\times$ \texttt{\#FFFFFF} & 4,5:1 & Texto na cor primária sobre a página (por exemplo, links); o texto sobre um fundo da cor primária é tratado por \texttt{readableTextOn}, e não por este par\\
\addlinespace
\texttt{secondaryColor} & secundária $\times$ \texttt{\#FFFFFF} & 4,5:1 & Texto branco sobre a barra superior e os cabeçalhos\\
\addlinespace
\texttt{accentColor} & destaque $\times$ \texttt{\#FFFFFF} & 3:1 & Marcadores e realce não textual (item ativo e \en{hover} da navegação): contraste de componentes de interface (WCAG 1.4.11); o foco não usa esta cor\\
\bottomrule
\end{tabularx}
\fonte{Elaborado pelos autores, a partir da resposta ao pedido P-002 (item 5) e da revisão do merge M-005.}
\end{table}

As funções de cálculo pertencem à UI e estão em \arq{src/lib/theme.js} (merge M-005), um módulo puro --- não importa o Next.js nem o Prisma e não usa APIs do navegador ---, que o servidor e o navegador compartilham:

\begin{itemize}
  \item \texttt{contrastRatio(a, b)} devolve a razão da WCAG 2.1 entre duas cores \texttt{\#RRGGBB}, sem arredondar e independentemente da ordem dos argumentos;
  \item \texttt{contrastWarnings(\{ primaryColor, secondaryColor, accentColor \})} devolve a lista de avisos \texttt{\{ code, field, ratio, minimum, message \}}, com o código \texttt{LOW\_CONTRAST}, vazia se nenhuma cor falha; a mensagem vem em português, por exemplo ``A cor primária tem contraste de 3,67:1 com o branco, abaixo do mínimo de 4,5:1 (WCAG 2.1 AA).'';
  \item \texttt{readableTextOn(cor)} devolve \texttt{\#FFFFFF} ou o texto escuro padrão, o de maior razão de contraste com a cor dada (no empate, o branco).
\end{itemize}

A comparação com o mínimo usa a razão \emph{sem arredondar}; o campo \texttt{ratio} e a mensagem trazem a razão \emph{truncada} em duas casas, para que 4,499 não apareça como 4,50 e pareça cumprir o mínimo (por isso a primária padrão aparece como 3,67, e não 3,68). O formulário de tema (M11) usa \texttt{contrastWarnings} na pré-visualização local, e o Route Handler de \texttt{PUT /api/configuration} (Core) a chama, depois da escrita, com as cores já gravadas --- inclusive quando o corpo da requisição não traz cores, pois o \texttt{PUT} é parcial --- e devolve a lista no campo \texttt{warnings} da resposta, ao lado de \texttt{data} (Seção~\ref{sec:arq-api}). O campo está sempre presente na resposta do \texttt{PUT} (lista vazia quando nenhuma cor falha), o aviso não bloqueia o salvamento (o status continua \texttt{200}) e o \texttt{GET} não o devolve. Cores fora do formato \texttt{\#RRGGBB} lançam \texttt{TypeError}: não chegam às funções, pois a validação de formato as recusa antes (RN09).

\textbf{Texto sobre a cor primária.} O branco fixo não atende ao nível AA com a paleta padrão: sobre a primária \texttt{\#3b82f6}, ele dá 3,68:1. Por isso o texto sobre um fundo da cor primária (o botão principal e o item ativo da navegação) usa a variável \arq{--vitrine-on-primary}, exposta como o utilitário \texttt{text-on-primary}, cujo valor é \texttt{readableTextOn(primaryColor)}; com a paleta padrão, \texttt{\#171717}, que dá 4,87:1. O \en{layout} raiz define essa variável junto com as três cores no atributo \texttt{style} de \texttt{<html>}, e quem sobrescrever \arq{--vitrine-primary} deve sobrescrever também \arq{--vitrine-on-primary}; do contrário, uma primária escura ficaria com texto escuro. O texto sobre a cor secundária continua branco (\arq{--vitrine-on-secondary}), protegido pelo par da secundária. A decisão, tomada com o Engenheiro Chefe, consta na \adr{ADR-020}.

\textbf{Limites conhecidos.} (i) Como o texto escuro é \texttt{\#171717}, e não o preto, existe uma faixa de primárias de luminância média (luminância relativa entre 0,183 e 0,214, cerca de 6\% das cores \texttt{\#RRGGBB}; por exemplo, \texttt{\#7b7b7b} e \texttt{\#676fe3}) em que nem o branco nem o escuro chegam a 4,5:1; o pior caso é 4,23:1. Nessa faixa o par da primária também reprova e o aviso da RN09 é emitido, de modo que o administrador é informado. Com o preto puro (\texttt{\#000000}), o pior caso passa a 4,58:1 e a faixa desaparece; essa troca foi decidida na \adr{ADR-020} e está pendente na UI. (ii) A primária padrão tem 3,68:1 também sobre a página, e a de destaque padrão, 2,15:1: com a paleta padrão da RN08, texto na cor primária (links) não atende ao nível AA, e marcadores na cor de destaque ficam abaixo de 3:1; o primeiro \texttt{PUT /api/configuration} com essa paleta devolve, portanto, dois avisos (primária e destaque). Por isso os links do conteúdo (M07 e M08) não podem depender apenas da cor primária para atingir 4,5:1 com uma paleta qualquer.

\vitrinependencia{UI}{Trocar o texto escuro de \texttt{readableTextOn}, e o valor padrão de \arq{--vitrine-on-primary}, de \texttt{\#171717} para \texttt{\#000000}, ajustando \arq{tests/ui/theme.test.js}; e propor, no pedido do M07 ou do M08, o mecanismo de cor dos links do conteúdo, que atenda a 4,5:1 com uma paleta qualquer.}

\vitrinefigura[tipo=diagrama, arquivo=seq-personalizacao-tema, agente=Arquiteto, largura=0.95,
  origem={doc/diagrams/src/seq-personalizacao-tema.puml},
  descricao={Sequência de personalização: formulário, PUT na API, validação de hexadecimais e contraste, atualização no banco, revalidação do layout e leitura da configuração pelo RootLayout na próxima visita.}]
  {Sequência de personalização visual}{fig:seq-personalizacao-tema}

\section{Interface de Programação HTTP}\label{sec:arq-api}

\subsection{Convenções}

As rotas JSON seguem um REST pragmático (Seção~\ref{sec:fund-renderizacao}). Toda rota de conteúdo e de configuração exige sessão válida, verificada no servidor antes de qualquer outra ação. Entradas e saídas são JSON em UTF-8; as rotas de escrita aceitam apenas \texttt{Content-Type: application/json}, e todas as respostas levam \texttt{Cache-Control: no-store}. Códigos de status: \texttt{200}/\texttt{201}/\texttt{204} para sucesso; \texttt{400} para JSON malformado ou \texttt{Content-Type} diferente de JSON (código \arq{VALIDATION_ERROR}); \texttt{401} sem sessão válida ou com mutação vinda de outra origem; \texttt{404} recurso inexistente; \texttt{409} conflito (por exemplo, \en{slug} duplicado); \texttt{422} falha de validação; \texttt{429} excesso de tentativas de login (resposta do better-auth); \texttt{500} erro interno. Essas convenções foram aprovadas na \adr{ADR-018}.

\textbf{Resposta de sucesso.} Toda resposta de sucesso com corpo usa o envelope \texttt{\{ "data": ... \}}, que deixa espaço para campos adicionais ao lado de \texttt{data}, como \texttt{warnings} (Seção~\ref{sec:arq-contraste}). O \texttt{POST} de criação responde \texttt{201} e o cabeçalho \texttt{Location} com o caminho do novo recurso; o \texttt{DELETE} responde \texttt{204} sem corpo; e a reordenação devolve em \texttt{data} a árvore atualizada.

\begin{lstlisting}[language=VJSON, caption={Formato padrão de sucesso}, label={lst:sucesso-json}]
{
  "data": { "id": "cm...", "title": "Introdução", "slug": "introducao", "status": "PUBLISHED" }
}
\end{lstlisting}

\textbf{Resposta de erro.} Os erros têm formato único; \texttt{details} é uma lista com um item por campo, no formato \texttt{\{ field, message \}} (validação) ou \texttt{\{ field, value \}} (valor recusado). Um \texttt{409} de \en{slug} traz em \texttt{details[0].suggestion} um \en{slug} livre entre os irmãos de destino (UC05, fluxo A1).

\begin{lstlisting}[language=VJSON, caption={Formato padrão de erro}, label={lst:erro-json}]
{
  "error": {
    "code": "SLUG_CONFLICT",
    "message": "Já existe um conteúdo com este slug na mesma coleção.",
    "details": [{ "field": "slug", "value": "introducao", "suggestion": "introducao-2" }]
  }
}
\end{lstlisting}

Códigos de erro previstos: \arq{VALIDATION_ERROR}, \arq{UNAUTHENTICATED}, \arq{NOT_FOUND}, \arq{SLUG_CONFLICT}, \arq{INVALID_PARENT}, \arq{CYCLE_DETECTED} e \arq{INTERNAL_ERROR}.

\textbf{Regras comuns de entrada.} Campos desconhecidos são descartados, inclusive \texttt{authorId}, que vem sempre da sessão (RN13). Um \texttt{PATCH} ou \texttt{PUT} sem nenhum campo conhecido responde \texttt{422}. O \en{slug} reservado também é conferido ao mover um conteúdo para a raiz (por \texttt{PATCH} ou em lote). A mutação cujo cabeçalho \texttt{Origin} indica outro site é recusada com \texttt{401} (\arq{UNAUTHENTICATED}), pois o cookie de sessão não é aceito em requisições de outra origem e o contrato não prevê \texttt{403} (Seção~\ref{sec:arq-seguranca}).

\subsubsection{Erros da Camada de Dados}\label{sec:arq-dataerror}

As funções de \arq{src/lib/db.js} lançam a classe \texttt{DataError} (com os campos \texttt{code}, \texttt{message} e \texttt{details}) quando uma regra que depende do estado do banco é violada: existência, unicidade, tipo do pai, ciclo e domínio de \texttt{type}, \texttt{status} e \texttt{order}. A validação de \emph{formato} (expressão regular e \en{slugs} reservados da RN03, URL \texttt{http}/\texttt{https} da RN17, cor \texttt{\#RRGGBB} da RN09) permanece nos validadores do Core (\arq{validators.js}). Cada Route Handler captura o \texttt{DataError} e o traduz conforme a Tabela~\ref{tab:dataerror}; qualquer outro erro vira \texttt{500} com o código \arq{INTERNAL_ERROR} e mensagem genérica, sem expor detalhes internos. As funções de escrita não chamam \arq{revalidatePath} (RN10): a revalidação é da camada de rota.

\begingroup\small
\begin{longtable}{@{}L{3.9cm}L{1.2cm}L{3.3cm}L{5.8cm}@{}}
\caption{Mapeamento dos códigos de \texttt{DataError} para a resposta HTTP}\label{tab:dataerror}\\
\toprule
\textbf{Código do \texttt{DataError}} & \textbf{HTTP} & \textbf{Código da API} & \textbf{Situação}\\
\midrule
\endfirsthead
\toprule
\textbf{Código do \texttt{DataError}} & \textbf{HTTP} & \textbf{Código da API} & \textbf{Situação}\\
\midrule
\endhead
\bottomrule
\endfoot
\arq{NOT_FOUND} & 404 & \arq{NOT_FOUND} & Conteúdo (ou, ao reordenar, algum item do lote) inexistente\\
\addlinespace
\arq{SLUG_CONFLICT} & 409 & \arq{SLUG_CONFLICT} & \en{Slug} repetido entre irmãos ou na raiz (RN03)\\
\addlinespace
\arq{PARENT_NOT_FOUND} & 422 & \arq{INVALID_PARENT} & O pai informado não existe\\
\addlinespace
\arq{PARENT_NOT_COLLECTION} & 422 & \arq{INVALID_PARENT} & O pai informado é um documento (RN02)\\
\addlinespace
\arq{CYCLE} & 422 & \arq{CYCLE_DETECTED} & O destino é o próprio nó ou um descendente (RN06)\\
\addlinespace
\arq{TYPE_HAS_CHILDREN} & 422 & \arq{VALIDATION_ERROR} & Coleção com filhos não pode virar documento (\texttt{details.field} = \texttt{type})\\
\addlinespace
\arq{AUTHOR_NOT_FOUND} & 422 & \arq{VALIDATION_ERROR} & Autor inexistente; como o autor vem da sessão, indica sessão inconsistente\\
\addlinespace
\arq{INVALID_INPUT} & 422 & \arq{VALIDATION_ERROR} & Campo obrigatório ausente ou fora do domínio (\texttt{type}, \texttt{status}, \texttt{order}, lote de reordenação)\\
\addlinespace
(outro erro) & 500 & \arq{INTERNAL_ERROR} & Falha não prevista; registrar no servidor, resposta genérica\\
\end{longtable}
\fonte{Elaborado pelos autores, a partir dos códigos de \arq{src/lib/db.js} e da sugestão do agente Dados (merge M-003).}
\endgroup

Duas observações de implementação. Em \texttt{reorderContents}, os itens do lote são aplicados em sequência e validados contra o estado resultante dos anteriores, de modo que a ordem dos itens importa quando há ciclo intermediário. E o redirecionamento de uma coleção ao primeiro documento publicado (RN14) não tem função própria: o Core usa \texttt{getContentTree} com \texttt{onlyPublished} e \texttt{resolvePath}.

\subsection{Contrato dos \emph{Endpoints}}

\begin{longtable}{@{}L{1.5cm}L{3.9cm}L{1.6cm}L{4.4cm}L{2.4cm}@{}}
\caption{Contrato das rotas HTTP (normativo)}\label{tab:api}\\
\toprule
\textbf{Método} & \textbf{Caminho} & \textbf{Sessão} & \textbf{Descrição} & \textbf{Sucesso}\\
\midrule
\endfirsthead
\toprule
\textbf{Método} & \textbf{Caminho} & \textbf{Sessão} & \textbf{Descrição} & \textbf{Sucesso}\\
\midrule
\endhead
\bottomrule
\endfoot
\texttt{GET} & \arq{/api/contents} & Sim & Árvore completa, inclusive rascunhos, sem o corpo dos documentos e com o \texttt{status} de cada nó & 200\\[1pt]
\texttt{POST} & \arq{/api/contents} & Sim & Cria conteúdo (\texttt{title}, \texttt{slug}, \texttt{type}, \texttt{parentId}, \texttt{body}, \texttt{order}, \texttt{status}; este, se omitido, assume \texttt{PUBLISHED}); o autor vem da sessão & 201 e \texttt{Location}\\[1pt]
\texttt{GET} & \arq{/api/contents/:id} & Sim & Detalhe, incluindo o corpo e o \texttt{status} & 200\\[1pt]
\texttt{PATCH} & \arq{/api/contents/:id} & Sim & Atualização parcial, inclusive de \texttt{status}; dispara revalidação (RN10, RN16) & 200\\[1pt]
\texttt{DELETE} & \arq{/api/contents/:id} & Sim & Exclui o nó e seus descendentes (RN07) & 204\\[1pt]
\texttt{PATCH} & \arq{/api/contents/reorder} & Sim & Lote \texttt{[\{id, parentId, order\}]}, em transação (RN05, RN06); devolve a árvore atualizada & 200\\[1pt]
\texttt{GET} & \arq{/api/configuration} & Sim & Configuração vigente & 200\\[1pt]
\texttt{PUT} & \arq{/api/configuration} & Sim & Atualização \emph{parcial} (campos ausentes não mudam) de identidade, cores e dos seis links sociais (RN17); logotipo e \en{favicon} aceitam URL \texttt{http(s)} ou caminho de \arq{public/} iniciado por uma única barra (D-10); devolve \texttt{warnings} de contraste & 200\\[1pt]
\texttt{POST} & \arq{/api/auth/sign-in/email} & Não & Login (rota do better-auth), com limite de 5 tentativas por minuto por IP & 200 (429 acima do limite)\\[1pt]
\texttt{POST} & \arq{/api/auth/sign-up/email} & Não & Recusado: cadastro público desabilitado (RF22) & 400 (\arq{EMAIL_PASSWORD_SIGN_UP_DISABLED})\\[1pt]
\texttt{POST} & \arq{/api/auth/sign-out} & Sim & Encerra a sessão (RF20); o painel usa, em seu lugar, a \en{Server Function} \texttt{signOut} (\arq{admin/_lib/actions.js}) & 200\\[1pt]
\texttt{GET} & \arq{/api/auth/get-session} & Não & Sessão corrente ou nula & 200\\[1pt]
\end{longtable}
\fonte{Elaborado pelos autores. As rotas \arq{/api/auth/*} são as padrão do better-auth, conferidas na versão 1.7.6 (merge M-004); as demais rotas padrão da biblioteca (por exemplo, troca de senha) também ficam expostas em \arq{/api/auth/*} e exigem sessão.}

A leitura pública da árvore, dos documentos e da configuração \textbf{não} usa essas rotas: ocorre em Server Components, por chamadas diretas às funções de \arq{db.js}, o que mantém as páginas estáticas e aplica o filtro de publicação (RN15). As rotas de conteúdo e de configuração servem ao painel administrativo.

Os campos de entrada têm as seguintes regras de validação no servidor (\texttt{422} em caso de falha): \texttt{status} deve ser \texttt{DRAFT} ou \texttt{PUBLISHED}; cada link social, quando preenchido, deve ser uma URL absoluta com protocolo \texttt{http} ou \texttt{https}, e a cadeia vazia equivale a \texttt{null} (RN17); as cores devem ser \texttt{\#RRGGBB} (RN09); o \en{slug} segue a RN03 (1 a 80 caracteres, sem \en{slugs} reservados na raiz). A resposta de \texttt{422} traz um item de \texttt{details} por campo inválido.

\section{Roteamento e Navegação}\label{sec:arq-rotas}

A Figura~\ref{fig:nav-mapa-do-site} resume as rotas. Segmentos estáticos (\texttt{admin}, \texttt{login}, \texttt{api}) têm precedência sobre o \en{catch-all} \arq{[...slug]}, razão pela qual esses nomes são reservados no nível raiz (RN03). A estrutura visual de cada página segue os protótipos da IC (Seção~\ref{sec:prototipagem}): barra superior com logotipo, nome e coleções raiz; \en{sidebar} da coleção corrente; área de conteúdo; e rodapé.

\vitrinefigura[tipo=diagrama, arquivo=nav-mapa-do-site, agente=Arquiteto, largura=0.9,
  origem={doc/diagrams/src/nav-mapa-do-site.puml},
  descricao={Mapa mental das rotas: área pública (início e leitura de documento), login, painel administrativo (visão geral, árvore, novo conteúdo, editor, configurações) e API.}]
  {Mapa de rotas e navegação}{fig:nav-mapa-do-site}

\begingroup\small
\begin{longtable}{@{}L{3.3cm}L{2.9cm}L{6.0cm}L{2.2cm}@{}}
\caption{Inventário de componentes de interface e situação após o merge M-005}\label{tab:componentes-ui}\\
\toprule
\textbf{Componente} & \textbf{Diretório} & \textbf{Responsabilidade} & \textbf{Situação}\\
\midrule
\endfirsthead
\toprule
\textbf{Componente} & \textbf{Diretório} & \textbf{Responsabilidade} & \textbf{Situação}\\
\midrule
\endhead
\bottomrule
\endfoot
Navbar & \arq{components/layout} & Logotipo, nome do site e coleções raiz (RF13) & Planejado\\
Sidebar & \arq{components/layout} & Árvore da coleção corrente, item ativo, recolher níveis (RF12) & Planejado\\
Footer, SocialLinks & \arq{components/layout} & Rodapé: direitos gerados a partir do nome do site e do ano, e ícones apenas das redes com link preenchido, com os atributos \texttt{noopener} e \texttt{noreferrer} (RF24, RN17) & Planejado\\
MarkdownRenderer & \arq{components/markdown} & Pipeline de Markdown, matemática e código (RF09 a RF11) & Planejado\\
Editor, Toolbar, Preview & \arq{components/editor} & Edição com painel dividido, rascunho local e publicação (RF08, RF15, RF16) & Planejado\\
AdminShell & \arq{components/admin} & Casca do painel: atalho para o conteúdo, barra superior, navegação e área principal; \en{Server Component} de apresentação, que não valida a sessão (M10) & Implementado\\
AdminHeader & \arq{components/admin} & Barra superior: nome do site, nome e e-mail do usuário e botão ``Sair'', que recebe a \en{Server Function} \texttt{signOut} por \en{prop} (RF20) & Implementado\\
AdminSidebar & \arq{components/admin} & Navegação do painel, com item ativo pela rota e recolhível em telas estreitas (RNF06); Componente Cliente; seções sem página aparecem como ``em breve'' & Implementado\\
AdminOverview & \arq{components/admin} & Visão geral provisória: identifica a sessão e anuncia as seções futuras & Implementado\\
ContentTree & \arq{components/admin} & Árvore do painel com criar, editar, mover e excluir, mostrando o estado de cada nó (RF03, RF06, RF07, RF26) & Planejado\\
ContentForm & \arq{components/admin} & Formulário de título, \en{slug}, pai e estado de publicação (RF05, RF26) & Planejado\\
StatusToggle & \arq{components/admin} & Indicador e alternância \texttt{DRAFT}/\texttt{PUBLISHED} na árvore e no editor, com destaque para rascunhos (RF26, RN16) & Planejado\\
ThemeForm & \arq{components/admin} & Identidade, cores e links sociais, com pré-visualização e aviso de contraste (RF18, RF19, RF24) & Planejado\\
LoginForm & \arq{components/admin} & Formulário de autenticação (RF01, UC08): mensagem de erro única, aviso de excesso de tentativas e falha de conexão; Componente Cliente & Implementado\\
Button, icons & \arq{components/ui} & Botão com as variantes \texttt{primary} e \texttt{inverse}, só com tokens de tema, e ícones em SVG embutido, decorativos (RF19) & Implementado\\
Demais primitivos & \arq{components/ui} & Diálogo de confirmação, \en{toast}, campo de cor & Planejado\\
\end{longtable}
\fonte{Elaborado pelos autores. O inventário nasceu como proposta para o agente UI; os componentes ``Implementado'' vieram do merge M-005.}
\endgroup

\vitrinefigura[tipo=diagrama, arquivo=seq-reordenacao-arvore, agente=Arquiteto, largura=0.95,
  origem={doc/diagrams/src/seq-reordenacao-arvore.puml},
  descricao={Sequência de reordenação e movimentação na árvore: atualização otimista, PATCH em lote, validação de pai, ciclos e slugs, transação e revalidação do layout.}]
  {Sequência de reordenação e movimentação de conteúdos}{fig:seq-reordenacao-arvore}

\section{Segurança}\label{sec:arq-seguranca}

A Tabela~\ref{tab:ameacas} relaciona as ameaças consideradas aos controles previstos, tomando como referência o OWASP Top 10 \cite{owasptop10} e as práticas de verificação do ASVS \cite{owaspasvs}.

\begin{longtable}{@{}L{3.3cm}L{4.4cm}L{5.4cm}@{}}
\caption{Ameaças e controles}\label{tab:ameacas}\\
\toprule
\textbf{Ameaça} & \textbf{Vetor} & \textbf{Controle}\\
\midrule
\endfirsthead
\toprule
\textbf{Ameaça} & \textbf{Vetor} & \textbf{Controle}\\
\midrule
\endhead
\bottomrule
\endfoot
XSS armazenado & Markdown com HTML ou URL \texttt{javascript:} & Sem HTML bruto (RN11); manter a transformação de URL do \arq{react-markdown}; revisar qualquer \en{plugin} novo\\[1pt]
CSRF & Requisição de outro site com o \en{cookie} do administrador & \en{Cookie} \texttt{SameSite=Lax}; verificação de origem nas mutações (um \texttt{Origin} diferente do próprio site, ou \texttt{null}, responde \texttt{401} \arq{UNAUTHENTICATED} e nada é gravado; \adr{ADR-018}); rotas de escrita aceitam apenas JSON\\[1pt]
Injeção de SQL & Entradas em consultas & Prisma parametriza; proibido \texttt{\$queryRawUnsafe}\\[1pt]
Força bruta e enumeração & Tentativas repetidas de login & No máximo cinco tentativas por minuto por IP (resposta 429); mensagem genérica de erro, igual para e-mail inexistente e senha errada (UC08, CT02); o \en{proxy} deve definir \texttt{x-forwarded-for} (Seção~\ref{sec:arq-auth})\\[1pt]
Acesso indevido à API & Chamada direta sem passar pelo painel & Sessão verificada em cada Route Handler (RF21), antes de ler o corpo; Proxy apenas otimista\\[1pt]
Rotas padrão do better-auth & Rotas de \arq{/api/auth/*} além do login (por exemplo, troca de senha e listagem de sessões) & Exigem sessão e passam pela mesma limitação de requisições; o cadastro público está desabilitado (RF22)\\[1pt]
Exposição de segredos & \arq{.env} no repositório, na imagem Docker ou variáveis no navegador & \arq{.env} ignorado pelo Git e excluído da imagem por \arq{.dockerignore}, fornecido ao contêiner em tempo de execução; nenhuma variável com prefixo \arq{NEXT_PUBLIC_} para segredos\\[1pt]
Vazamento de rascunho & Nó \texttt{DRAFT} acessível por URL, pela árvore ou por página em cache & Filtro de publicação na camada de dados (RN15); revalidação do \en{layout} ao alterar o estado (RN16); testes CT31 e CT32\\[1pt]
Link social malicioso & URL \texttt{javascript:} ou \texttt{data:} em um dos seis campos do rodapé & Validação de protocolo \texttt{http}/\texttt{https} no servidor (RN17); atributos \texttt{noopener} e \texttt{noreferrer}; teste CT33\\[1pt]
Injeção de CSS & Valor de cor malicioso & Validação \texttt{\#RRGGBB} no servidor (RN09)\\[1pt]
Dependências vulneráveis & Pacotes desatualizados, com vulnerabilidade conhecida ou maliciosos & Revisão de cada nova dependência pelo Arquiteto; \texttt{npm audit} a cada entrega; \en{lockfile} versionado; os quatro alertas atuais foram avaliados e aceitos (Seção~\ref{sec:arq-npm-audit}, \adr{ADR-019})\\[1pt]
\en{Clickjacking} & Página incorporada em \texttt{iframe} & Cabeçalho \arq{frame-ancestors}/\arq{X-Frame-Options}\\[1pt]
Exclusão acidental & Cascata indevida & Confirmação explícita; \en{backup} (RNF10)\\[1pt]
Perda do banco & Remoção do diretório de dados, falha do disco da VPS ou atualização que descarta o volume & Banco em \en{bind mount} fora da imagem; \en{backup} periódico com \texttt{VACUUM INTO} e cópia externa (Apêndice~\ref{ap:implantacao}; risco R12)\\[1pt]
\end{longtable}
\fonte{Elaborado pelos autores.}

Além dos controles acima, o \arq{next.config.mjs} deve definir cabeçalhos de segurança (\texttt{X-Content-Type-Options: nosniff}, \arq{Referrer-Policy}, política de enquadramento e \arq{Permissions-Policy} restritiva); HSTS e TLS ficam a cargo do \en{proxy} reverso. Uma política CSP com \en{nonce} é desejável e será avaliada pelo agente Core.

\subsection{Vulnerabilidades Conhecidas em Dependências}\label{sec:arq-npm-audit}

Em 01/10/2026, na revisão do merge M-004, o \texttt{npm audit} acusou \textbf{quatro vulnerabilidades de severidade alta}, todas em dependências transitivas do Prisma 7 (Tabela~\ref{tab:npm-audit}). O npm só as corrige com \texttt{npm audit fix -{}-force}, que instalaria o \texttt{prisma@6.19.3}: um rebaixamento de versão maior, incompatível com o que o projeto adota (\arq{prisma.config.ts}, gerador \texttt{prisma-client} e \en{driver adapter}; \adr{ADR-014}). A \texttt{7.10.0}, instalada, é a última versão 7 publicada, e o \texttt{npm audit} lista como afetada toda a faixa do \arq{prisma} até a \texttt{8.1.0-dev.6}; portanto, não há versão 7 corrigida. Decidiu-se \textbf{manter o Prisma 7 e não aplicar a correção automática}, aceitando os alertas, pelas razões da tabela (\adr{ADR-019}). Como \arq{prisma} agora está em \texttt{dependencies} (para aplicar migrações na imagem), esses pacotes também estão presentes na imagem de produção, embora não sejam exercidos.

\begingroup\small
\begin{longtable}{@{}L{3.3cm}L{5.4cm}L{5.5cm}@{}}
\caption{Alertas do \texttt{npm audit} aceitos em 01/10/2026}\label{tab:npm-audit}\\
\toprule
\textbf{Pacote} & \textbf{Aviso} & \textbf{Avaliação para o Vitrine}\\
\midrule
\endfirsthead
\toprule
\textbf{Pacote} & \textbf{Aviso} & \textbf{Avaliação para o Vitrine}\\
\midrule
\endhead
\bottomrule
\endfoot
\arq{deepmerge-ts} (antes da 8.0.0; usado por \arq{@prisma/config}) & Esgotamento de pilha ao mesclar grafos de objetos recursivos (GHSA-ggr8-5vv4-36mx) & Roda na CLI do Prisma, ao mesclar a configuração de \arq{prisma.config.ts}, escrita pela equipe; não recebe entrada de usuário nem é exercido pelo servidor\\
\addlinespace
\arq{@prisma/config} e \arq{prisma} & Alertas em cadeia, por dependerem do \arq{deepmerge-ts} e do \arq{mysql2} & O mesmo das linhas vizinhas; sem correção em nenhuma versão 7 disponível\\
\addlinespace
\arq{mysql2} (até a 3.23.0; dependência do \arq{prisma} e do \arq{better-auth}) & Rebaixamento do plugin de autenticação para \arq{mysql_clear_password}, com vazamento de credenciais (GHSA-3f6p-5ww8-9rcr), e descompressão (\en{inflate}) sem limite no protocolo comprimido, com negação de serviço (GHSA-rgwj-5xj2-c3m3) & Só é exercido ao conectar a um servidor MySQL; o Vitrine usa SQLite por \arq{better-sqlite3} e nunca abre conexão MySQL\\
\end{longtable}
\fonte{Elaborado pelos autores, a partir da saída do \texttt{npm audit} e de \texttt{npm ls} (merge M-004).}
\endgroup

A decisão será reavaliada a cada entrega (o \texttt{npm audit} é critério de saída, Seção~\ref{sec:testes-saida}) e sempre que o Prisma publicar uma versão 7 com as dependências corrigidas. Enquanto isso, os alertas aceitos não contam como vulnerabilidade crítica pendente.

\section{Implantação e Operação}\label{sec:arq-implantacao}

Por decisão do Engenheiro Chefe (\decisao{D-04}, \adr{ADR-011}), o alvo de implantação é uma VPS tradicional com Docker Engine e Docker Compose. A aplicação é empacotada em uma imagem construída em dois estágios (instalação e \en{build}; execução) que executa o servidor Next.js (\texttt{next start}) como usuário sem privilégios, em um único contêiner. O arquivo do SQLite fica em um \en{bind mount} do hospedeiro (\arq{./data}, montado em \arq{/data}), fora da imagem, e as variáveis de ambiente vêm de um arquivo \arq{.env} também fora da imagem. Um \en{proxy} reverso com TLS (Caddy ou Nginx) encaminha as requisições para a porta 3000 do contêiner, publicada apenas em \arq{127.0.0.1} (Figura~\ref{fig:impl-implantacao}). O passo a passo, com esboços do \arq{Dockerfile} e do \arq{docker-compose.yml}, está no Apêndice~\ref{ap:implantacao}.

\vitrinefigura[tipo=diagrama, arquivo=impl-implantacao, agente=Arquiteto, largura=0.95,
  origem={doc/diagrams/src/impl-implantacao.puml},
  descricao={Diagrama de implantação: VPS com Docker Engine; proxy reverso com TLS; contêiner vitrine com o build do Next.js, o cliente do Prisma e o adaptador better-sqlite3; volume ./data com o arquivo SQLite (e WAL) e cópias de segurança; arquivo .env fora da imagem.}]
  {Diagrama de implantação}{fig:impl-implantacao}

Os requisitos normativos sobre a imagem, que o módulo M13 deve atender, são:

\begin{itemize}
  \item \textbf{Módulo nativo.} O \arq{better-sqlite3} é um módulo nativo (\adr{ADR-014}); os estágios de \en{build} e de execução devem usar a mesma versão do Node.js e a mesma distribuição (\en{glibc}, por exemplo \texttt{node:22-bookworm-slim}), e o estágio de \en{build} precisa de compilador (\texttt{python3}, \texttt{make}, \texttt{g++}) caso não haja binário pré-compilado;
  \item \textbf{Cliente gerado.} \texttt{prisma generate} roda antes de \texttt{next build}, pois \arq{src/generated/prisma} não é versionado;
  \item \textbf{Migrações como passo explícito.} \texttt{prisma migrate deploy} é executado por um comando próprio antes de subir a nova versão (por exemplo, \texttt{docker compose run -{}-rm vitrine npx prisma migrate deploy}), e não pelo comando de partida da aplicação; isso exige que a CLI do Prisma e o \arq{dotenv} estejam na imagem, o que está garantido desde o merge M-002, que os moveu para \texttt{dependencies} (\adr{ADR-014});
  \item \textbf{Administrador inicial.} Depois das migrações e do \en{seed}, o administrador é criado por \texttt{docker compose run -{}-rm vitrine npm run admin:create}, que exige \arq{ADMIN_EMAIL} e \arq{ADMIN_PASSWORD} (e aceita \arq{ADMIN_NAME}) no ambiente da execução. O \en{seed} (\texttt{prisma db seed}) \emph{não} cria o administrador: ele só garante a configuração padrão (\adr{ADR-018});
  \item \textbf{\texttt{tsx} na imagem.} O \texttt{admin:create} e o \texttt{prisma db seed} rodam com o \texttt{tsx}, hoje em \texttt{devDependencies}; a imagem de produção deve manter as dependências de desenvolvimento (\texttt{npm ci} completo) ou mover o \texttt{tsx} para \texttt{dependencies}, caso contrário esses dois comandos deixam de funcionar;
  \item \textbf{\texttt{DATABASE\_URL} no \en{build}.} O \texttt{next build} importa as rotas, que importam \arq{db.js}, e este lança erro se \arq{DATABASE_URL} não está definida; o estágio de \en{build} do Dockerfile deve defini-la com qualquer valor bem formado (por exemplo, \texttt{ENV} \arq{DATABASE_URL=file:/data/vitrine.db}), pois nenhuma consulta ocorre no \en{build} e o arquivo não precisa existir. Segredos (\arq{BETTER_AUTH_SECRET}, \arq{ADMIN_PASSWORD}) \emph{não} entram no \en{build}: o aviso do segredo padrão registrado nessa etapa é esperado e inofensivo;
  \item \textbf{Segredo em execução.} \arq{BETTER_AUTH_SECRET} e \arq{BETTER_AUTH_URL} (em \texttt{https}, para o \en{cookie} \texttt{Secure}) são obrigatórios na execução; sem o segredo, as rotas de autenticação e de API respondem 500;
  \item \textbf{Instância única.} Um único contêiner, pois o cache de ISR é por instância (Seção~\ref{sec:arq-isr}), o SQLite tem um escritor por vez e a limitação de tentativas de login guarda os contadores em memória;
  \item \textbf{Segredos.} \arq{.env} e \arq{data/} constam do \arq{.dockerignore} e do \arq{.gitignore}.
\end{itemize}

O relatório de IC citava a plataforma Vercel como exemplo de implantação \cite{reis2026ic}. O sistema de arquivos de funções \emph{serverless} é efêmero, de modo que um arquivo SQLite local não persistiria entre invocações nesse tipo de ambiente; a decisão por uma VPS com volume elimina esse obstáculo, que constava como risco R01, para esta versão. Plataformas dessa natureza só seriam viáveis com a troca do \en{driver adapter} por um banco remoto compatível (por exemplo, libSQL/Turso), o que está fora do escopo (\adr{ADR-011}).

\section{Desempenho e Escalabilidade}\label{sec:arq-desempenho}

O perfil de carga é dominado por leituras. A ISR evita consultas ao banco na maioria das requisições públicas; quando há regeneração, a consulta da árvore é linear no número de nós e a resolução do caminho é proporcional à profundidade. O SQLite em modo WAL permite leitores simultâneos à escrita, que é serializada \cite{sqlitewal}; como as escritas ocorrem apenas por ação do administrador, a limitação de um escritor por vez não é relevante para o perfil esperado \cite{sqlitewhentouse}. Os limites práticos --- tamanho da árvore, número de documentos pré-renderizados no \en{build} e tempo de regeneração do \en{layout} raiz --- serão medidos no Capítulo~\ref{cap:testes}.

\section{Síntese do Capítulo}

O capítulo especificou o monólito modular, suas visões, regras de dependência, modelo de dados (com o estado de publicação e os links sociais), autenticação em camadas, estratégia de cache e de visibilidade, pipeline de conteúdo, tematização, contrato de API, segurança e implantação por imagem Docker. As observações sobre o esquema atual (Tabela~\ref{tab:obs-schema}) e os valores padrão da \adr{ADR-017} formam a lista de tarefas de refinamento dos agentes; o Capítulo~\ref{cap:implementacao} acompanha a execução.

%% ===================================================================================
%% Capítulo 6 -- Implementação
%% Este capítulo nasce como PLANO de implementação por módulo e é reescrito, módulo a
%% módulo, como RELATO de implementação à medida que os merges ocorrem.
%% Ciclo de vida de cada seção (responsável: Arquiteto):
%%   1. \statusmodulo{planejado}   -> a seção descreve tarefas e critérios de aceite;
%%   2. \statusmodulo{andamento}   -> existe ramificação em revisão no mural;
%%   3. \statusmodulo{implementado}-> após o merge: trocar as tarefas por relato real,
%%                                    incluir trechos (\vitrinecodigo) e capturas (\vitrinefigura).
%% Figuras: salvar os PNG em doc/tex/figuras/telas/<arquivo>.png -- aparecem sozinhas.
%% ===================================================================================
\chapter{Implementação}\label{cap:implementacao}

Este capítulo acompanha a construção do Vitrine módulo a módulo. Cada módulo tem um selo de status, um agente responsável, os requisitos que atende, as tarefas, os critérios de aceite e as evidências esperadas (capturas de tela e trechos de código). Os blocos tracejados amarelos indicam evidências ainda não produzidas; eles são substituídos automaticamente quando o arquivo correspondente é adicionado ao repositório (Apêndice~\ref{ap:pendencias}).

\section{Estado dos Módulos}\label{sec:modulos}

\begin{longtable}{@{}L{0.9cm}L{3.2cm}L{1.6cm}L{2.2cm}L{2.9cm}L{2.9cm}@{}}
\caption{Módulos, responsáveis, requisitos e estado}\label{tab:modulos}\\
\toprule
\textbf{ID} & \textbf{Módulo} & \textbf{Agente(s)} & \textbf{RF} & \textbf{Ramificação sugerida} & \textbf{Status}\\
\midrule
\endfirsthead
\toprule
\textbf{ID} & \textbf{Módulo} & \textbf{Agente(s)} & \textbf{RF} & \textbf{Ramificação sugerida} & \textbf{Status}\\
\midrule
\endhead
\bottomrule
\endfoot
M01 & Modelo de dados e migrações & Dados & RF02, RF03, RF22, RF23 & \arq{feat/dados-schema-inicial} & \statusmodulo{andamento}\\[1pt]
M02 & Acesso a dados e árvore & Dados & RF02, RF03, RF14 & \arq{feat/dados-db-e-arvore} & \statusmodulo{planejado}\\[1pt]
M03 & Autenticação e sessão & Core & RF01, RF20 a RF22 & \arq{feat/core-auth} & \statusmodulo{planejado}\\[1pt]
M04 & API de conteúdo & Core & RF05 a RF07, RF16 & \arq{feat/core-api-conteudo} & \statusmodulo{planejado}\\[1pt]
M05 & API de configuração & Core & RF18, RF19, RF21 & \arq{feat/core-api-configuracao} & \statusmodulo{planejado}\\[1pt]
M06 & Revalidação de cache (ISR) & Core & RF17 & \arq{feat/core-revalidacao} & \statusmodulo{planejado}\\[1pt]
M07 & Leitura pública e navegação & UI, Core & RF12 a RF14, RF24 & \arq{feat/ui-leitura-publica} & \statusmodulo{planejado}\\[1pt]
M08 & Renderização de Markdown & UI & RF09 a RF11 & \arq{feat/ui-markdown} & \statusmodulo{planejado}\\[1pt]
M09 & Editor de conteúdo & UI & RF08, RF15, RF16 & \arq{feat/ui-editor} & \statusmodulo{planejado}\\[1pt]
M10 & Painel de gestão da árvore & UI, Core & RF03, RF06, RF07 & \arq{feat/ui-painel-arvore} & \statusmodulo{planejado}\\[1pt]
M11 & Personalização visual & UI, Core, Dados & RF04, RF18, RF19 & \arq{feat/ui-personalizacao} & \statusmodulo{planejado}\\[1pt]
M12 & SEO e metadados & Core & RF25 & \arq{feat/core-seo} & \statusmodulo{planejado}\\[1pt]
M13 & Implantação e operação & Core, Dados & RF22; RNF01 & \arq{feat/core-implantacao} & \statusmodulo{planejado}\\[1pt]
M14 & Testes automatizados & Todos & RNF08 & \arq{feat/<dominio>-testes} & \statusmodulo{planejado}\\[1pt]
\end{longtable}
\fonte{Elaborado pelos autores. O estado de M01 é ``em andamento'' porque o esquema inicial já existe em \arq{prisma/schema.prisma}; faltam migrações, \en{seed} e conformidade com o Prisma 7.}

A ordem recomendada de execução respeita as dependências: M01 e M02 (dados) habilitam M03 a M06 (servidor), que habilitam M07 e M08 (leitura), seguidos de M09 a M11 (edição e personalização); M12 a M14 acompanham o restante.

\section{Ambiente de Desenvolvimento}\label{sec:impl-ambiente}

O Código~\ref{lst:comandos-dev} resume os comandos previstos. É um \emph{esboço}: a sequência exata depende da resolução de \decisao{D-09} (configuração do Prisma 7) e deve ser validada e, se necessário, corrigida pelos agentes Dados e Core, registrando a versão final no Apêndice~\ref{ap:implantacao}.

\begin{lstlisting}[language=VSh, caption={Esboço: comandos do ciclo de desenvolvimento}, label={lst:comandos-dev}]
# 1. Instalar dependências
npm install

# 2. Variáveis de ambiente (.env nunca é versionado; .env.example lista os nomes)
cp .env.example .env

# 3. Banco de dados (Dados) -- conferir as skills do Prisma em .claude/skills/
npx prisma generate
npx prisma migrate dev        # desenvolvimento
npx prisma db seed            # configuração padrão e primeiro administrador

# 4. Executar e verificar
npm run dev                   # http://localhost:3000
npm run lint
npm run build && npm run start   # verifica produção e o comportamento do ISR
\end{lstlisting}

O Next.js 16 fornece, em \arq{node_modules/next/dist/docs/}, a documentação da versão instalada; os agentes devem consultá-la antes de escrever código de roteamento, cache ou autenticação, pois a versão traz mudanças em relação ao conhecimento anterior (por exemplo, a renomeação de \en{Middleware} para \en{Proxy}).

\section{M01 e M02 --- Dados}\label{sec:impl-dados}

\noindent\statusmodulo{andamento}\quad Agente: \agente{Dados}. Requisitos: \req{RF02}, \req{RF03}, \req{RF14}, \req{RF22}, \req{RF23}.

\textbf{Situação.} O esquema inicial (Apêndice~\ref{ap:schema}) contém os seis modelos. Faltam migrações versionadas, \en{seed}, o módulo \arq{src/lib/db.js} e a adequação às práticas do Prisma 7. A Tabela~\ref{tab:prisma7} lista as divergências encontradas na configuração atual em relação às instruções do Prisma 7 incluídas nas \en{skills} do repositório (\arq{.claude/skills/prisma-orm-setup} e \arq{.claude/skills/prisma-upgrade-v7}).

\begin{table}[htb]
\caption{Conformidade da configuração do Prisma com a versão 7}\label{tab:prisma7}
\small
\begin{tabularx}{\linewidth}{@{}L{0.7cm}L{3.2cm}L{3.9cm}Y@{}}
\toprule
\textbf{\#} & \textbf{Item} & \textbf{Situação atual} & \textbf{Esperado / ação}\\
\midrule
1 & Arquivo de configuração & \arq{prisma7.config.ts} & Nome padrão \arq{prisma.config.ts}, ou uso explícito de \arq{--config}\\
2 & Gerador & \arq{prisma-client-js} (legado) & \arq{prisma-client} com \texttt{output}; alinhar com \arq{/src/generated/prisma} do \arq{.gitignore}\\
3 & URL do banco & \texttt{url = env("DATABASE\_URL")} no \texttt{datasource} & Mover a URL para o arquivo de configuração e manter só o \texttt{provider} no esquema\\
4 & \en{Driver adapter} & Ausente; \texttt{sqlite3} declarado & Instalar adapter do Prisma para SQLite (\arq{better-sqlite3} ou libSQL); remover \texttt{sqlite3}\\
5 & \arq{@prisma/client} & Em \arq{devDependencies} & Necessário em produção: mover para \texttt{dependencies}\\
6 & \arq{DATABASE_URL} relativa & Não verificada & Caminhos relativos podem resolver diferente na CLI e no \emph{runtime}; usar caminho explícito em produção\\
\bottomrule
\end{tabularx}
\fonte{Elaborado pelos autores, a partir do repositório e das \en{skills} do Prisma 7.}
\end{table}

\textbf{Tarefas.}
\begin{itemize}
  \item Resolver os itens da Tabela~\ref{tab:prisma7} e registrar a decisão final em \decisao{D-09};
  \item tratar as observações da Tabela~\ref{tab:obs-schema} (unicidade de \en{slug}, índices, cascata do autor, registro único de configuração);
  \item gerar a primeira migração; escrever o \en{seed} (configuração padrão e administrador inicial, RF22 e RF23);
  \item implementar \arq{src/lib/db.js} (instância única do cliente, funções de leitura da árvore, do documento por caminho e da configuração, e de escrita transacional);
  \item implementar \arq{src/lib/content-tree.js} conforme o Código~\ref{lst:algoritmo-arvore}.
\end{itemize}

\textbf{Critérios de aceite.} A migração aplica-se em um banco vazio; o \en{seed} é idempotente; a consulta da árvore retorna a estrutura ordenada sem carregar o corpo; a resolução de caminho devolve o nó correto ou nulo; nenhum arquivo fora de \arq{db.js} instancia o cliente.

\vitrinecodigo[linerange={1-60}]{Dados}{src/lib/db.js}{VJS}{Acesso a dados: instância única do cliente e consultas principais}

\vitrinecodigo{Dados}{src/lib/content-tree.js}{VJS}{Montagem da árvore e resolução de caminhos}

\vitrinefigura[tipo=teste, arquivo=teste-migracoes, agente=Dados, largura=0.85,
  descricao={Saída de terminal da aplicação da migração em banco vazio e da execução do seed (comandos e mensagens de sucesso).}]
  {Aplicação da migração e execução do \en{seed}}{fig:teste-migracoes}

\vitrinependencia{Dados}{Substituir esta seção por relato: decisões tomadas em D-05, D-06, D-09 e D-14, índices criados, estratégia de ordenação e exemplos de consulta.}

\section{M03 --- Autenticação e Sessão}\label{sec:impl-auth}

\noindent\statusmodulo{planejado}\quad Agente: \agente{Core}. Requisitos: \req{RF01}, \req{RF20}, \req{RF21}, \req{RF22}. Regras: RN01, RN12.

\textbf{Tarefas.} Configurar o better-auth com o adaptador do Prisma (Código~\ref{lst:auth-esboco}); expor o Route Handler; criar o cliente React; implementar \arq{src/proxy.js} (verificação otimista); criar a página \arq{/login} e o \en{logout}; validar a sessão nos layouts e nas rotas do painel; implementar o \en{bootstrap} do primeiro administrador com cadastro público desabilitado (\decisao{D-03}); habilitar a limitação de tentativas de login.

\textbf{Critérios de aceite.} Credenciais válidas redirecionam ao destino original; inválidas retornam mensagem genérica; \arq{/admin} sem \en{cookie} redireciona a \arq{/login}; qualquer rota de escrita sem sessão retorna 401 (mesmo com o Proxy desativado); \en{logout} invalida a sessão no banco; o \en{cookie} tem os atributos \texttt{HttpOnly} e \texttt{SameSite}, e \texttt{Secure} em produção.

\vitrinecodigo{Core}{src/lib/auth.js}{VJS}{Configuração do servidor de autenticação (better-auth)}

\vitrinecodigo{Core}{src/proxy.js}{VJS}{Proxy: verificação otimista de sessão}

\vitrinefigura[tipo=tela, arquivo=tela-login, agente=UI, largura=0.6,
  descricao={Tela de login: formulário de e-mail e senha, botão de entrar; estado de erro com mensagem genérica.},
  ref={Fig. 5 do relatório de IC (estilo visual do painel).}]
  {Tela de autenticação do painel administrativo}{fig:tela-login}

\vitrinependencia{Core}{Descrever a configuração final (opções de sessão, limitação de tentativas, variáveis de ambiente efetivas), o procedimento de bootstrap escolhido e eventuais diferenças entre a API do better-auth 1.7.6 e o esboço do Código~\ref{lst:auth-esboco}.}

\section{M04, M05 e M06 --- API, Configuração e Revalidação}\label{sec:impl-api}

\noindent\statusmodulo{planejado}\quad Agente: \agente{Core}. Requisitos: \req{RF05} a \req{RF07}, \req{RF16} a \req{RF19}, \req{RF21}.

\textbf{Tarefas.} Implementar as rotas da Tabela~\ref{tab:api} com validação de entrada (\decisao{D-11}), verificação de sessão em cada handler, formato de erro padronizado (Código~\ref{lst:erro-json}) e chamada às funções de \arq{db.js}; implementar \arq{src/lib/revalidate.js} com as regras da Tabela~\ref{tab:revalidacao}; proteger \arq{/api/configuration} e a atualização de tema.

\textbf{Critérios de aceite.} Cada rota responde com os códigos de status previstos; conflitos de \en{slug} retornam 409 com \arq{SLUG_CONFLICT}; ciclos retornam \arq{CYCLE_DETECTED}; uma alteração só de corpo revalida apenas a página e uma alteração estrutural revalida o \en{layout} raiz; nenhuma rota escreve sem sessão.

\vitrinecodigo{Core}{src/app/api/contents/route.js}{VJS}{Route Handler de coleção de conteúdos (GET e POST)}

\vitrinecodigo{Core}{src/lib/revalidate.js}{VJS}{Regras de revalidação sob demanda (RN10)}

\vitrinependencia{Core}{Relatar o contrato efetivamente implementado, divergências em relação à Tabela~\ref{tab:api} (que exigem aprovação prévia) e o resultado da verificação de revalidação com o cabeçalho x-nextjs-cache.}

\section{M07 e M08 --- Leitura Pública e Renderização}\label{sec:impl-leitura}

\noindent\statusmodulo{planejado}\quad Agentes: \agente{UI} e \agente{Core}. Requisitos: \req{RF09} a \req{RF14}, \req{RF24}.

\textbf{Tarefas.} Configurar o \en{layout} raiz (idioma \arq{pt-BR}, fontes, variáveis de tema, importação da folha de estilos do KaTeX); implementar barra superior, \en{sidebar} e rodapé; implementar a página de documento em \arq{src/app/[...slug]/page.js} com \arq{generateStaticParams}, \texttt{revalidate} e \texttt{await params}; tratar coleção (RN14) e 404; implementar o \arq{MarkdownRenderer} com a cadeia da Seção~\ref{sec:arq-markdown}.

\textbf{Critérios de aceite.} Uma página pública é servida de cache após a primeira requisição (\texttt{HIT}); a \en{sidebar} destaca o item ativo e é navegável por teclado; fórmulas, tabelas e código renderizam corretamente; HTML bruto no Markdown é exibido como texto, não interpretado; a página tem pontuação de acessibilidade e de desempenho conforme \req{RNF02} e \req{RNF05}.

\vitrinecodigo{UI}{src/components/markdown/MarkdownRenderer.js}{VJS}{Componente de renderização de Markdown com matemática e código}

\vitrinecodigo{UI}{src/components/layout/Sidebar.js}{VJS}{Componente da \en{sidebar} hierárquica}

\vitrinefigura[tipo=tela, arquivo=tela-leitura-documento, agente=UI, largura=0.9,
  descricao={Página de documento completa: barra superior com coleções raiz, sidebar com três níveis (coleção, subcoleção, documentos) e item ativo destacado, conteúdo renderizado e rodapé. A página inicial (D-16) pode ser registrada em captura adicional opcional, tela-home.},
  ref={Fig. 3 do relatório de IC (instância FromCaio).}]
  {Página de leitura de um documento}{fig:tela-leitura-documento}

\vitrinefigura[tipo=tela, arquivo=tela-leitura-matematica, agente=UI, largura=0.85,
  descricao={Trecho de documento com fórmula em linha e em bloco renderizadas pelo KaTeX (por exemplo, uma integral e um somatório), ao lado do texto LaTeX-fonte.}]
  {Fórmulas \LaTeX{} renderizadas no conteúdo}{fig:tela-matematica}

\vitrinefigura[tipo=tela, arquivo=tela-leitura-codigo, agente=UI, largura=0.85,
  descricao={Trecho de documento com bloco de código realçado (JavaScript e Python) e uma tabela GFM.}]
  {Blocos de código com realce e tabelas}{fig:tela-codigo}

\vitrinefigura[tipo=tela, arquivo=tela-responsiva-mobile, agente=UI, largura=0.45,
  descricao={Duas capturas lado a lado em largura de 375 px: página de leitura com sidebar recolhida e com sidebar aberta.}]
  {Layout responsivo em tela pequena}{fig:tela-mobile}

\vitrinependencia{UI}{Descrever a composição final do layout, o tratamento de coleções e da página inicial (D-16), a biblioteca de realce escolhida (D-07), a adoção do remark-gfm (D-08) e as decisões de acessibilidade.}

\section{M09 --- Editor de Conteúdo}\label{sec:impl-editor}

\noindent\statusmodulo{planejado}\quad Agente: \agente{UI}. Requisitos: \req{RF08}, \req{RF09}, \req{RF15}, \req{RF16}.

\textbf{Tarefas.} Implementar o editor em painel dividido --- edição e pré-visualização com o mesmo \arq{MarkdownRenderer} da página pública --- com barra de ferramentas (negrito, itálico, tachado, títulos, links, imagens, código, citação, listas, tabela e fórmula); rascunho local por conteúdo (``Salvar localmente''), com restauração ao reabrir; e envio da versão (``Enviar nova versão'') à API.

\textbf{Critérios de aceite.} A pré-visualização é idêntica à renderização pública; o rascunho sobrevive ao recarregamento da página; o envio exibe confirmação e trata erros da API (401, 409, 422); os botões da barra de ferramentas têm rótulos acessíveis.

\vitrinecodigo{UI}{src/components/editor/Editor.js}{VJS}{Componente do editor com pré-visualização}

\vitrinefigura[tipo=tela, arquivo=tela-editor-markdown, agente=UI, largura=0.9,
  descricao={Editor em painel dividido: Markdown com fórmulas à esquerda, pré-visualização à direita, barra de ferramentas e botões Salvar localmente e Enviar nova versão.},
  ref={Fig. 6 do relatório de IC.}]
  {Editor de Markdown com pré-visualização}{fig:tela-editor}

\vitrinependencia{UI}{Relatar a implementação do editor, o formato do rascunho local e o tratamento de conflitos entre rascunho e versão publicada.}

\section{M10 --- Painel de Gestão da Árvore}\label{sec:impl-painel}

\noindent\statusmodulo{planejado}\quad Agentes: \agente{UI} e \agente{Core}. Requisitos: \req{RF03}, \req{RF06}, \req{RF07}.

\textbf{Tarefas.} Implementar a visão geral do painel, a árvore de conteúdos com ações por nó (criar documento ou coleção, editar, mover, reordenar, excluir com confirmação) e os formulários de criação e edição de metadados, integrados às rotas da Tabela~\ref{tab:api}.

\textbf{Critérios de aceite.} Criação, edição e exclusão refletem-se na \en{sidebar} pública; movimentos inválidos (ciclo, pai que não é coleção) são rejeitados com mensagem clara; a exclusão de coleção com descendentes exige confirmação que informe a quantidade de itens afetados.

\vitrinefigura[tipo=tela, arquivo=tela-admin-arvore-conteudo, agente=UI, largura=0.9,
  descricao={Painel com a árvore de conteúdos: itens com ações de criar, editar, mover e excluir; botão de nova publicação.},
  ref={Fig. 4 do relatório de IC.}]
  {Painel de gestão da árvore de conteúdos}{fig:tela-arvore}

\vitrinefigura[tipo=tela, arquivo=tela-admin-novo-conteudo, agente=UI, largura=0.8,
  descricao={Formulário de criação de conteúdo: título, slug sugerido, tipo (documento ou coleção) e seleção da coleção pai.}]
  {Formulário de criação de conteúdo}{fig:tela-novo}

\vitrinependencia{UI}{Relatar a implementação do painel, a solução de arrastar e soltar (ou ações de subir e descer) e os estados de erro.}

\section{M11 --- Personalização Visual}\label{sec:impl-personalizacao}

\noindent\statusmodulo{planejado}\quad Agentes: \agente{UI}, \agente{Core} e \agente{Dados}. Requisitos: \req{RF04}, \req{RF18}, \req{RF19}, \req{RF24}.

\textbf{Tarefas.} Implementar o formulário de identidade (nome, descrição, logotipo, \en{favicon}) e de paleta (três cores), com pré-visualização local e aviso de contraste; a rota \texttt{PUT /api/configuration}; a leitura da configuração no \en{layout} raiz; os tokens do Código~\ref{lst:tema-css}; e a decisão sobre armazenamento de logotipo/\en{favicon} (\decisao{D-10}) e dos links do rodapé (\decisao{D-19}).

\textbf{Critérios de aceite.} Alterar uma cor no painel muda a aparência de todas as páginas na próxima requisição, sem novo \en{build}; valores inválidos são recusados; o aviso de contraste aparece para combinações abaixo de 4,5:1; nome e logotipo aparecem na barra superior e nos metadados.

\vitrinecodigo{UI}{src/lib/theme.js}{VJS}{Geração das variáveis CSS e cálculo de contraste}

\vitrinefigura[tipo=tela, arquivo=tela-admin-configuracoes, agente=UI, largura=0.85,
  descricao={Tela de configurações: identidade (nome, descrição, logotipo, favicon) e paleta (seletores de cor e campos hexadecimais) com pré-visualização e aviso de contraste.},
  ref={Fig. 5 do relatório de IC.}]
  {Configuração da identidade e da paleta de cores}{fig:tela-configuracoes}

\vitrinefigura[tipo=tela, arquivo=tela-tema-antes-depois, agente=UI, largura=0.95,
  descricao={Duas capturas lado a lado da mesma página antes e depois de alterar a paleta, comprovando o efeito sem novo build.}]
  {Efeito da personalização sobre uma mesma página}{fig:tela-antes-depois}

\vitrinependencia{UI}{Relatar a implementação da personalização, o mapeamento final das cores (Seção~\ref{sec:arq-tema}) e o tratamento do modo escuro (D-17).}

\section{M12, M13 e M14 --- SEO, Implantação e Testes}\label{sec:impl-demais}

\noindent\statusmodulo{planejado}\quad

\textbf{M12 --- SEO e metadados (\agente{Core}).} Gerar \texttt{title} e \texttt{description} por página a partir do conteúdo e da configuração, \arq{robots.txt} e mapa do site (\req{RF25}). Critério de aceite: metadados presentes no HTML pré-renderizado de cada documento.

\textbf{M13 --- Implantação e operação (\agente{Core} e \agente{Dados}).} Criar o \arq{.env.example} (e a exceção correspondente no \arq{.gitignore}, pois o padrão \arq{.env*} o ignora); documentar o roteiro de instalação, atualização e \en{backup}; definir os cabeçalhos de segurança em \arq{next.config.mjs}. Critério de aceite: uma instalação limpa, seguindo apenas o Apêndice~\ref{ap:implantacao}, produz um sistema funcional.

\textbf{M14 --- Testes automatizados (todos).} Definir as ferramentas (\decisao{D-12}) e implementar os casos de teste do Capítulo~\ref{cap:testes} no domínio de cada agente.

\vitrinecodigo{Core}{.env.example}{VSh}{Variáveis de ambiente necessárias (sem valores)}

\section{Dificuldades e Lições Aprendidas}\label{sec:impl-licoes}

\vitrinependencia{Arquiteto}{Ao final da implementação, consolidar aqui as dificuldades técnicas encontradas, as decisões que mudaram em relação ao projeto do Capítulo 5 (com a referência à ADR correspondente), as lições do processo de DDD com agentes e os desvios entre o planejado e o realizado.}

\section{Síntese do Capítulo}

O capítulo organiza a implementação em catorze módulos com responsáveis, critérios de aceite e evidências, e registra o estado inicial: esquema de dados presente, demais módulos planejados. À medida que as integrações ocorrem, cada seção é reescrita como relato e as evidências substituem os marcadores.

%% ===================================================================================
%% Capítulo 7 -- Verificação, Validação e Testes
%% Define a estratégia, os critérios e os casos de teste (CT) rastreados aos requisitos.
%% Os RESULTADOS de execução entram no Capítulo 8. Dono: Arquiteto/Redator; cada agente
%% implementa e executa os casos do seu domínio e entrega as evidências em
%% doc/tex/figuras/testes/ (ver DESENVOLVEDORES.md).
%% ===================================================================================
\chapter{Verificação, Validação e Testes}\label{cap:testes}

A \emph{verificação} responde se o sistema foi construído corretamente, isto é, conforme a especificação dos Capítulos~\ref{cap:requisitos} e~\ref{cap:arquitetura}; a \emph{validação} responde se foi construído o sistema correto, isto é, se atende às necessidades do usuário \cite{sommerville2011,pressman2021}. Este capítulo define como ambas serão realizadas e quais evidências serão coletadas.

\section{Objetivos e Escopo}\label{sec:testes-objetivos}

Os objetivos da verificação e validação são: (i) demonstrar que cada requisito essencial e importante é atendido (\req{RF01} a \req{RF26}); (ii) comprovar as regras de negócio invariantes (RN01 a RN17) no servidor; (iii) medir as qualidades declaradas como metas (\req{RNF02}, \req{RNF04}, \req{RNF05}); e (iv) verificar o comportamento do cache ISR em modo de produção. Ficam fora do escopo testes de carga em grande escala, pois o perfil de uso esperado é de um único administrador e tráfego de leitura moderado.

\section{Estratégia de Testes}\label{sec:testes-estrategia}

A estratégia segue a proporção clássica entre tipos de teste: muitos testes de unidade, menos testes de integração e poucos testes de ponta a ponta (\emph{E2E}).

\begin{table}[htb]
\caption{Tipos de teste, alvo e ferramentas adotadas}\label{tab:tipos-teste}
\small
\begin{tabularx}{\linewidth}{@{}L{2.6cm}YL{4.2cm}@{}}
\toprule
\textbf{Tipo} & \textbf{Alvo} & \textbf{Ferramentas adotadas (D-12)}\\
\midrule
Unidade & Funções puras: montagem da árvore, resolução de caminho, validadores, cálculo de contraste, \en{slug} & Vitest \cite{vitest}\\
\addlinespace
Integração & Funções de \arq{db.js} contra um SQLite temporário e isolado; Route Handlers com sessão simulada & Vitest com banco de teste (Seção~\ref{sec:testes-banco})\\
\addlinespace
Ponta a ponta & Fluxos completos no navegador sobre o build de produção (login, publicação, leitura, personalização) & Playwright \cite{playwright}\\
\addlinespace
Não funcionais & Desempenho, acessibilidade, segurança e ISR & Lighthouse \cite{lighthouse}, verificação de cabeçalhos, \texttt{npm audit}\\
\bottomrule
\end{tabularx}
\fonte{Elaborado pelos autores.}
\end{table}

\textbf{Ambiente de teste.} Cada suíte de integração usa um arquivo SQLite temporário, criado a partir das migrações, com \en{seed} controlado (Seção~\ref{sec:testes-banco}). Os testes de ISR e de desempenho usam \texttt{next build} seguido de \texttt{next start}, porque o comportamento de cache em \texttt{next dev} não representa a produção \cite{nextjsdocs}. Nenhum teste usa o banco de desenvolvimento nem segredos reais. O caso CT29 exige uma máquina com Docker e Compose, de preferência uma VPS limpa, para validar a instalação por imagem (\adr{ADR-011}).

\textbf{Ferramenta oficial.} O Vitest \cite{vitest} (versão 5, dependência de desenvolvimento) é a ferramenta oficial de testes de unidade e de integração, aprovada pelo Engenheiro Chefe (\decisao{D-12}, \adr{ADR-017}). O comando é \texttt{npm test} (\texttt{vitest run}) e a configuração está em \arq{vitest.config.mjs}, na raiz, com o mesmo alias \arq{@/} do \arq{jsconfig.json}. Os testes ficam em \arq{tests/}, por domínio: \arq{tests/setup/} (infraestrutura compartilhada), \arq{tests/dados/}, \arq{tests/core/} e \arq{tests/ui/}, que reutilizam \arq{tests/setup/}. Os testes de ponta a ponta usam o Playwright \cite{playwright}. O \arq{vitest.config.mjs} inclui um \emph{plugin} que transforma o JSX dos arquivos \texttt{.js} de \arq{src/} e de \arq{tests/}: a \adr{ADR-008} manda escrever JSX em \texttt{.js}, e o Vite só o transforma em \texttt{.jsx}; sem o \emph{plugin}, um teste não consegue importar um arquivo de rota nem um componente. Ele usa \texttt{transformWithOxc}, do próprio Vite 8 (dependência transitiva do Vitest, sem pacote novo), e não afeta o \texttt{.ts} do cliente gerado do Prisma; numa troca de versão do Vitest ou do Vite, esse \emph{import} precisa ser reavaliado.

\textbf{Responsabilidades.} Cada agente implementa os testes do seu domínio (Tabela~\ref{tab:casos-teste}); o Arquiteto verifica a cobertura dos requisitos e consolida os resultados no Capítulo~\ref{cap:resultados}.

\section{Testes de Integração com Banco Isolado}\label{sec:testes-banco}

Os testes que exercitam o banco seguem uma estratégia de isolamento, implementada pelo agente Dados em \arq{tests/setup/} (merge M-003), que todo teste de integração deve reutilizar:

\begin{enumerate}
  \item \textbf{Banco-modelo migrado.} Uma vez por rodada, o \texttt{globalSetup} cria um diretório temporário do sistema e aplica nele as migrações reais com \texttt{prisma migrate deploy}, o mesmo comando da produção. Assim, o esquema testado é exatamente o das migrações, inclusive o índice parcial da raiz (Seção~\ref{sec:arq-indice-parcial}), que o Prisma não modela.
  \item \textbf{Cópia por arquivo.} Antes de cada arquivo de teste, o \texttt{setupFiles} copia o banco-modelo para um arquivo próprio e define \texttt{DATABASE\_URL} com o caminho absoluto dessa cópia, antes de qualquer \texttt{import} de \arq{db.js} (que lê a variável ao ser carregado). Os arquivos de teste não interferem entre si.
  \item \textbf{Nunca o banco de desenvolvimento.} O \arq{dev.db} e o \arq{.env} não entram em jogo: a variável é sempre definida pelo \en{setup}, e uma salvaguarda aborta o teste se o banco estiver fora do diretório temporário. Ao final da rodada o diretório é removido.
  \item \textbf{Velocidade sem perder as constraints.} Nos bancos de teste, e só neles, usa-se \texttt{PRAGMA synchronous = OFF}; isso dispensa o \texttt{fsync} e não altera o comportamento das restrições verificadas.
  \item \textbf{Ausência de \en{drift}.} Os testes de migração conferem que \texttt{migrate status} e \texttt{migrate diff} não acusam divergência por causa do índice parcial escrito à mão.
\end{enumerate}

A primeira rodada, no merge M-003, tem 72 testes em cinco arquivos: restrições físicas do banco (8), \arq{content-tree} (14), conteúdo em \arq{db.js} (40), configuração (6) e migração com \en{seed} (4). Cobrem os casos CT11, CT12, CT13 e CT28 e três restrições da migração inicial: \en{slug} único na raiz, \en{slug} único entre irmãos e bloqueio da exclusão de usuário que possui conteúdo. Para validar a eficácia da suíte, o agente Dados introduziu mutações temporárias (remover o índice parcial, remover o índice composto, desativar a checagem de ciclo, desativar o filtro \texttt{PUBLISHED}, executar a reordenação sem transação) e confirmou que, em cada caso, os testes esperados falharam. Testes de Route Handlers devem definir \texttt{DATABASE\_URL} antes do primeiro \texttt{import} de \arq{db.js}, como faz \arq{tests/setup/per-file.js}.

O Core reutiliza essa infraestrutura em \arq{tests/core/} (merge M-004), com 269 testes em nove arquivos: validadores (93), utilitários HTTP e mapeamento de erros (33), revalidação (18), autenticação (16), comando \texttt{admin:create} (6), guarda de sessão (37), API de conteúdo (48), API de configuração (13) e Proxy (5). A diferença é que esses testes usam o \textbf{better-auth real}: o administrador é criado no banco de teste, o login passa por \texttt{auth.handler} e os Route Handlers recebem o \en{cookie} verdadeiro, de modo que a sessão é validada no banco, como em produção. O comando \texttt{admin:create} é executado de verdade com o \texttt{tsx}, com o \arq{.env} desligado. Apenas o \texttt{revalidatePath} é substituído por um duplo (\texttt{vi.mock("next/cache")}), pois exige o servidor do Next.js; o arquivo \arq{tests/core/env.js} define um segredo fictício. As recusas de escrita sem sessão (CT04) chamam os \en{handlers} diretamente, sem o Proxy, o que comprova que a barreira real está na API. Com as 72 do Dados, a suíte tinha 341 testes nesse merge. O Core também aplicou mutações temporárias (22, das quais 21 fizeram testes falharem; a restante era redundante), e o Arquiteto repetiu, na revisão, o \texttt{lint}, o \texttt{build} e a suíte completa.

A UI acrescenta, em \arq{tests/ui/} (merge M-005), 36 casos em um arquivo, \arq{theme.test.js}, sobre funções puras e sobre o \arq{globals.css}: \texttt{contrastRatio} (preto e branco = 21, cores iguais = 1, simetria, maiúsculas e minúsculas, valores conhecidos e recusa de entradas fora de \texttt{\#RRGGBB}), \texttt{contrastWarnings} (lista vazia; paleta padrão com avisos da primária e do destaque; limiar de 4,5:1 entre \texttt{\#767676} e \texttt{\#777777}; limiar de 3:1 do destaque; truncamento da razão) e \texttt{readableTextOn}, além da consistência dos tokens (a paleta padrão da RN08, o mapeamento \texttt{--color-*}, o valor de \arq{--vitrine-on-primary} igual ao de \texttt{readableTextOn(primária)}, o contraste dos neutros e a ausência de \texttt{prefers-color-scheme}). Cobrem a parte de unidade do CT24; a parte de API foi coberta pelo Core no merge M-006 (adiante). Com os 72 do Dados e os 269 do Core, a suíte tinha 377 testes em 15 arquivos nesse merge, reproduzidos pelo Arquiteto (\texttt{lint}, \texttt{build} e suíte completa). Não há testes automatizados de componentes, pois o ambiente de DOM não foi adotado: o comportamento da interface foi verificado no navegador pela UI (Seção~\ref{sec:impl-auth}), e o Playwright segue previsto para os casos de ponta a ponta.

O Core ampliou, no merge M-006, a API de configuração em \arq{tests/core/api-configuration.test.js}, de 13 para 23 testes (CT24, parte de API), de modo que \arq{tests/core/} passa a ter \textbf{279 testes em nove arquivos}. Os casos novos afirmam o campo, a razão truncada e o mínimo de cada aviso, e não só a presença da chave \texttt{warnings}: o destaque padrão (\texttt{\#f59e0b}) já gera aviso em qualquer \texttt{PUT} com a paleta padrão, e um teste genérico passaria pelo motivo errado. Cobrem a paleta padrão gravada (primária 3,67 e destaque 2,14, nessa ordem); o formato \texttt{\{ code, field, ratio, minimum, message \}}; a paleta sem falhas (\texttt{warnings} presente e vazio); a correção só da primária, que mantém o aviso do destaque; a secundária logo abaixo do limiar (4,47); a persistência dos avisos em um \texttt{PUT} posterior sem cores; o cálculo sobre as três cores \emph{lidas do banco}, e não sobre o corpo; a gravação e a revalidação normais com três avisos; e a falha no cálculo, que responde \texttt{500} depois de a configuração já estar gravada e o \en{layout} revalidado (resiliência; Seção~\ref{sec:arq-contraste}). O \texttt{GET} e a resposta \texttt{422} não trazem \texttt{warnings}.

Para observar e fazer falhar a função da UI sem alterá-la, o arquivo usa um \emph{espião transparente}: \texttt{vi.mock("@/lib/theme")} substitui \texttt{contrastWarnings} por um \texttt{vi.fn} que repassa a chamada à implementação real, o que permite conferir os argumentos e simular a falha com \texttt{mockImplementationOnce}, sem mudar o comportamento nos demais testes. O Core provou a eficácia dos testes com oito mutações temporárias da rota (calcular sobre o corpo; usar as cores padrão quando o corpo não traz cor; omitir \texttt{warnings} quando vazio; devolver sempre lista vazia; devolver \texttt{warnings} no \texttt{GET}; calcular antes de revalidar; responder 422 quando há aviso; avaliar só a primária), todas detectadas; o Arquiteto reproduziu três delas (calcular antes de revalidar, omitir \texttt{warnings} quando vazio e calcular sobre o corpo), também detectadas. Com os 72 do Dados, os 279 do Core e os 36 da UI, a suíte tinha 387 testes em 15 arquivos nesse merge, reproduzidos pelo Arquiteto (\texttt{lint}, \texttt{build} e suíte completa).

No merge M-007, o Core acrescentou 48 testes em dois arquivos novos de \arq{tests/core/}, de modo que o Core passa a ter \textbf{327 testes em 11 arquivos}. O \arq{site.test.js} (30) cobre o módulo \arq{src/app/_lib/site.js}: \texttt{buildThemeStyle} (as quatro variáveis; o texto sobre a primária igual a \texttt{readableTextOn(primária)}, comparado com a função e não com um literal, pois o texto escuro vai mudar para \texttt{\#000000}; doze entradas inválidas omitidas sem lançar; \texttt{null} dá objeto vazio), o nome e a descrição com os padrões e \texttt{getSiteConfiguration} (sem consulta quando \arq{NEXT_PHASE} vale \texttt{phase-production-build}; \texttt{null} quando a consulta lança). O \arq{layouts.test.js} (18) renderiza de verdade, com \texttt{renderToStaticMarkup} de \texttt{react-dom/server} e sobre o SQLite temporário, o \en{layout} raiz (\texttt{lang}; as quatro variáveis no \texttt{style} do \texttt{<html>} e nada no \texttt{<body>}; a paleta nova na renderização seguinte; ausência de \texttt{style} sem banco ou no \en{build}; cor inválida gravada direto no banco que não chega ao HTML), o \texttt{generateMetadata}, o \en{layout} do painel (a \en{prop} \texttt{siteName} e só ela; o nome no cabeçalho; a sessão validada \emph{antes} de ler a configuração) e a página 404 (texto em português, link para o início, nenhuma cor fixa do Tailwind, nenhuma menção a rascunho). Só se substituem os módulos que dependem do servidor do Next.js (\texttt{next/font/google}, a sessão do painel, a Server Function de sair e \texttt{usePathname}); o banco é o real, com um espião transparente em \texttt{getConfiguration}, como no M-006. O Core provou a eficácia dos testes com 21 mutações temporárias, todas detectadas; o Arquiteto reproduziu quatro (consultar no \en{build}, aceitar cor sem validação, ler a configuração antes da sessão e remover o \emph{plugin} de JSX), também detectadas.

Além dos testes, o CT23 foi verificado sem navegador por HTTP, sobre o \en{build} de produção com \texttt{next start}, pelo Core e, de forma independente, pelo Arquiteto: \arq{/login} e o painel trazem a paleta e o nome do banco; a 404 responde em português, com \texttt{noindex}; e, depois de um \texttt{PUT /api/configuration}, as páginas passam à nova paleta sem novo \en{build} (Seção~\ref{sec:arq-tema-estatico}). O \en{build} também passou com \texttt{DATABASE\_URL=file:/data/vitrine.db}, de diretório inexistente, e com um banco existente sem tabelas, sem criar pasta nem arquivo (o arquivo seguiu vazio). Com os 72 do Dados, os 327 do Core e os 36 da UI, a suíte tem \textbf{435 testes em 17 arquivos}, reproduzidos pelo Arquiteto (\texttt{lint}, \texttt{build} e suíte completa) na revisão do merge M-007.

\section{Casos de Teste}\label{sec:casos-teste}

A Tabela~\ref{tab:casos-teste} relaciona cada caso de teste aos requisitos e regras que ele verifica. Um caso é considerado \emph{aprovado} quando o resultado observado coincide com o esperado; um requisito essencial só é considerado atendido quando todos os seus casos de teste aprovam.

\begin{longtable}{@{}L{1.1cm}L{7.0cm}L{2.2cm}L{1.6cm}L{1.6cm}@{}}
\caption{Casos de teste e rastreabilidade}\label{tab:casos-teste}\\
\toprule
\textbf{ID} & \textbf{Caso de teste (ação e resultado esperado)} & \textbf{Verifica} & \textbf{Tipo} & \textbf{Agente}\\
\midrule
\endfirsthead
\toprule
\textbf{ID} & \textbf{Caso de teste (ação e resultado esperado)} & \textbf{Verifica} & \textbf{Tipo} & \textbf{Agente}\\
\midrule
\endhead
\bottomrule
\endfoot
CT01 & Login com credenciais válidas redireciona ao painel e cria a sessão & RF01 & E2E & Core\\[2pt]
CT02 & Login inválido retorna mensagem genérica, sem indicar o campo incorreto & RF01; RNF03 & API & Core\\[2pt]
CT03 & Acesso a \arq{/admin} sem sessão redireciona a \arq{/login} & RF01; RF21 & E2E & Core\\[2pt]
CT04 & Qualquer rota de escrita sem sessão retorna 401 & RF21; RN01 & API & Core\\[2pt]
CT05 & Encerrar sessão invalida a sessão no banco e remove o \en{cookie} & RF20 & API & Core\\[2pt]
CT06 & O cadastro público está desabilitado após o \en{bootstrap} & RF22 & API & Core\\[2pt]
CT07 & Criar documento válido retorna 201 e o item aparece na árvore & RF05; RF02 & API & Core\\[2pt]
CT08 & \en{Slug} duplicado entre irmãos retorna 409 (\arq{SLUG_CONFLICT}) & RN03 & API & Core\\[2pt]
CT09 & Pai que não é coleção retorna 422 (\arq{INVALID_PARENT}) & RN02 & API & Core\\[2pt]
CT10 & Mover um nó para dentro de um descendente é rejeitado (\arq{CYCLE_DETECTED}) & RN06 & API & Core\\[2pt]
CT11 & Reordenar irmãos persiste a nova ordem e a listagem a respeita & RF03; RN05 & Integração & Dados\\[2pt]
CT12 & Excluir uma coleção remove todos os descendentes & RF07; RN07 & Integração & Dados\\[2pt]
CT13 & Resolver caminho existente e publicado retorna o nó; inexistente, ou não publicado (RN15), retorna nulo (404) & RF14; RN04; RN15 & Unidade & Dados\\[2pt]
CT14 & Abrir uma coleção redireciona ao primeiro documento descendente & RN14 & E2E & UI\\[2pt]
CT15 & Fórmulas em linha e em bloco são renderizadas pelo KaTeX & RF09 & Unidade & UI\\[2pt]
CT16 & HTML bruto no Markdown é exibido como texto e não é executado & RN11; RNF03 & Unidade & UI\\[2pt]
CT17 & Tabelas, tachado e listas de tarefas (GFM) são renderizados & RF11 & Unidade & UI\\[2pt]
CT18 & Blocos de código recebem realce conforme a linguagem & RF10 & E2E & UI\\[2pt]
CT19 & A pré-visualização do editor é idêntica à página pública & RF08 & E2E & UI\\[2pt]
CT20 & O rascunho local sobrevive ao recarregamento da página & RF15 & E2E & UI\\[2pt]
CT21 & Publicar nova versão atualiza a página pública (revalidação da página) & RF16; RF17; RN10 & E2E & Core\\[2pt]
CT22 & Mudança estrutural (título, \en{slug}, pai, ordem) revalida o \en{layout} raiz & RF17; RN10 & E2E & Core\\[2pt]
CT23 & Alterar cores reflete-se nas páginas sem novo \en{build} & RF19 & E2E & UI\\[2pt]
CT24 & Cor fora do formato é recusada; contraste insuficiente gera aviso & RN09 & API/Unid. & Core, UI\\[2pt]
CT25 & Metadados de SEO presentes no HTML pré-renderizado & RF25; RNF07 & E2E & Core\\[2pt]
CT26 & A \en{sidebar} é navegável por teclado e destaca o item ativo & RF12; RNF05 & E2E & UI\\[2pt]
CT27 & O layout permanece utilizável em largura de 320 px & RNF06 & E2E & UI\\[2pt]
CT28 & A migração aplica-se em banco vazio e o \en{seed} é idempotente & RF23 & Integração & Dados\\[2pt]
CT29 & Uma instalação limpa em VPS com Docker, seguindo o Apêndice~\ref{ap:implantacao}, produz sistema funcional, e o banco persiste após recriar o contêiner & RNF01; D-04 & Manual & Core\\[2pt]
CT30 & O \en{backup} do banco (\texttt{VACUUM INTO}) pode ser restaurado com os dados íntegros & RNF10 & Manual & Dados\\[2pt]
CT31 & Um conteúdo \texttt{DRAFT}, e toda a sua subárvore, responde 404 por URL direta e não aparece na \en{sidebar}, na barra superior nem no mapa do site; o painel continua listando todos os conteúdos & RF26; RN15; RF25 & E2E & Core\\[2pt]
CT32 & A criação sem estado explícito resulta em \texttt{PUBLISHED}; alternar o estado revalida o \en{layout} raiz e o conteúdo passa a aparecer ou a desaparecer na visita seguinte, sem novo \en{build} & RN16; RF17; RF26 & E2E & Core\\[2pt]
CT33 & Link social com protocolo diferente de \texttt{http} e \texttt{https} é recusado (422); campo vazio é aceito e o ícone não é exibido; links válidos recebem \texttt{noopener} e \texttt{noreferrer} & RF24; RN17; RNF03 & API/Unid. & Core, UI\\[2pt]
\end{longtable}
\fonte{Elaborado pelos autores.}

\section{Testes Não Funcionais}\label{sec:testes-nf}

\subsection{Desempenho e ISR}

O desempenho será medido com o Lighthouse \cite{lighthouse} sobre o \en{build} de produção, em perfil móvel simulado e em desktop, registrando as métricas centrais de experiência \cite{webvitals}, em especial o \en{Largest Contentful Paint} (LCP). A meta é LCP $\leq$ 2,5 s no percentil 75 e pontuação de desempenho $\geq$ 90 (\req{RNF02}). O comportamento do cache será verificado por requisições repetidas à mesma página, observando o cabeçalho \arq{x-nextjs-cache}: a primeira requisição deve resultar em \texttt{MISS} ou \texttt{HIT} conforme a pré-renderização, as seguintes em \texttt{HIT}, e a primeira após uma mutação em \texttt{MISS} ou \texttt{REVALIDATED} (casos CT21 e CT22).

\vitrinefigura[tipo=teste, arquivo=teste-lighthouse-mobile, agente=UI, largura=0.85,
  descricao={Relatório do Lighthouse em perfil móvel para uma página de documento no build de produção, com as quatro pontuações e as métricas LCP, CLS e TBT.}]
  {Relatório do Lighthouse (perfil móvel)}{fig:teste-lh-mobile}

\vitrinefigura[tipo=teste, arquivo=teste-isr-cabecalhos, agente=Core, largura=0.85,
  descricao={Saída de terminal com requisições repetidas e o cabeçalho x-nextjs-cache (MISS, HIT, HIT, e após mutação REVALIDATED ou MISS).}]
  {Evidência do comportamento do cache ISR}{fig:teste-isr}

\subsection{Segurança}

A verificação de segurança consiste em: executar os casos CT02, CT04, CT06, CT16, CT31 e CT33; conferir os cabeçalhos de segurança da resposta; executar \texttt{npm audit} e registrar as vulnerabilidades encontradas (os quatro alertas altos já avaliados e aceitos constam na Seção~\ref{sec:arq-npm-audit} e na \adr{ADR-019}; qualquer alerta novo exige nova avaliação); e percorrer a lista da Tabela~\ref{tab:ameacas}, confirmando cada controle. Como referência de itens de verificação utiliza-se o ASVS do OWASP \cite{owaspasvs}.

\subsection{Acessibilidade}

A acessibilidade será avaliada segundo o WCAG 2.1 nível AA \cite{wcag21}, combinando verificação automatizada (contraste, rótulos, estrutura de títulos) e verificação manual (navegação apenas por teclado, foco visível, leitura por leitor de tela em um fluxo representativo).

\vitrinefigura[tipo=teste, arquivo=teste-acessibilidade, agente=UI, largura=0.85,
  descricao={Resultado da verificação automatizada de acessibilidade (por exemplo, categoria Acessibilidade do Lighthouse ou ferramenta equivalente) para as páginas pública e do painel.}]
  {Resultado da verificação automatizada de acessibilidade}{fig:teste-a11y}

\subsection{Usabilidade}

A usabilidade será validada com usuários do perfil do Administrador (docentes e pesquisadores), em sessões individuais com tarefas roteirizadas e o questionário SUS \cite{brooke1996sus}. Estudos de usabilidade clássicos indicam que amostras pequenas já revelam a maior parte dos problemas de uma interface \cite{nielsen1993usability}; planeja-se um grupo de cinco a oito participantes. As tarefas previstas são:

\begin{enumerate}[label=\textbf{T\arabic*.}, leftmargin=1.2cm]
  \item entrar no painel e criar uma coleção chamada ``Cálculo'';
  \item criar, dentro dela, um documento com um título, um parágrafo e uma fórmula;
  \item conferir que o documento está publicado e encontrá-lo na página pública; colocá-lo em rascunho e confirmar que deixa de aparecer;
  \item mover o documento para outra coleção e reordená-lo;
  \item alterar a cor primária do site e confirmar o efeito na página pública;
  \item preencher o link do GitHub nas configurações e localizar o ícone correspondente no rodapé.
\end{enumerate}

A meta é pontuação SUS média igual ou superior a 68 (\req{RNF04}) e taxa de conclusão de pelo menos 80\% por tarefa.


\section{Critérios de Saída}\label{sec:testes-saida}

A versão é considerada pronta para a entrega (\en{release}) quando: (i) todos os casos de teste dos requisitos \emph{Essenciais} aprovam; (ii) os casos dos requisitos \emph{Importantes} aprovam ou têm justificativa de exceção registrada; (iii) \texttt{npm run lint}, \texttt{npm run build} e \texttt{npm test} passam sem erros; (iv) não há vulnerabilidades críticas pendentes no relatório de auditoria de dependências, e as vulnerabilidades altas se limitam às aceitas e documentadas na \adr{ADR-019}; (v) as metas de \req{RNF02}, \req{RNF04} e \req{RNF05} foram medidas e registradas, ainda que não atingidas, com análise no Capítulo~\ref{cap:resultados}.

\vitrinefigura[tipo=teste, arquivo=teste-suite-automatizada, agente=Core, largura=0.85,
  descricao={Saída do executor de testes (unidade, integração e E2E) mostrando o total de casos aprovados, falhos e ignorados e o tempo de execução.}]
  {Execução da suíte de testes automatizados}{fig:teste-suite}

\section{Síntese do Capítulo}

O capítulo definiu a estratégia de verificação e validação, trinta e três casos de teste rastreados a requisitos e regras, os procedimentos não funcionais (desempenho e ISR, segurança, acessibilidade e usabilidade) e os critérios de saída. A execução e a análise dos resultados são apresentadas no capítulo seguinte.

%% ===================================================================================
%% Capítulo 8 -- Resultados e Discussão
%% Parte 1 (IC) é factual, extraída do relatório de IC. As demais partes são TABELAS E
%% ROTEIROS a preencher pelo Arquiteto após a execução dos testes do Capítulo 7.
%% Regra: nunca inventar resultados. Célula sem medição permanece com "---".
%% ===================================================================================
\chapter{Resultados e Discussão}\label{cap:resultados}

Este capítulo apresenta os resultados obtidos na Iniciação Científica (Seção~\ref{sec:res-ic}), os resultados da consolidação do monólito (Seções~\ref{sec:res-requisitos} e~\ref{sec:res-metricas}), o atendimento aos objetivos (Seção~\ref{sec:res-objetivos}) e a discussão, incluindo as ameaças à validade (Seções~\ref{sec:res-discussao} e~\ref{sec:res-ameacas}).

\section{Resultados da Iniciação Científica}\label{sec:res-ic}

O relatório de IC \cite{reis2026ic} registra que os objetivos propostos foram atingidos por meio de uma abordagem de desenvolvimento incremental, e destaca as seguintes conquistas técnicas:

\begin{itemize}
  \item a implementação de um editor Markdown capaz de renderizar com precisão elementos essenciais da escrita científica, como fórmulas em \LaTeX{} e blocos de código com realce de sintaxe;
  \item o desenvolvimento de uma interface de personalização visual dinâmica, que permite gerenciar a estética do site por variáveis CSS persistidas em banco de dados, com reflexo instantâneo por ISR;
  \item a garantia da perenidade do conteúdo digital, uma vez que a arquitetura baseada em arquivos portáteis assegura que o pesquisador detenha a posse integral dos seus dados.
\end{itemize}

O mesmo relatório apresenta como limitações o prazo de execução --- que restringiu a modelagem da interface a protótipos de baixa fidelidade --- e o sistema de autenticação, que, embora seguro e funcional para a gestão básica, ainda poderia ser expandido para múltiplos níveis de permissão em cenários colaborativos.

\section{Cobertura dos Requisitos}\label{sec:res-requisitos}

A Tabela~\ref{tab:res-rf} registra a situação final de cada requisito funcional e a evidência que a sustenta (caso de teste aprovado, captura de tela ou trecho de código). A situação assume um dos valores: \emph{atendido}, \emph{parcial}, \emph{não atendido} ou \emph{adiado}. Até a conclusão da implementação, todas as células permanecem sem preenchimento.

\begin{longtable}{@{}L{1.4cm}L{4.4cm}L{2.4cm}L{5.4cm}@{}}
\caption{Situação dos requisitos funcionais}\label{tab:res-rf}\\
\toprule
\textbf{ID} & \textbf{Requisito} & \textbf{Situação} & \textbf{Evidência}\\
\midrule
\endfirsthead
\toprule
\textbf{ID} & \textbf{Requisito} & \textbf{Situação} & \textbf{Evidência}\\
\midrule
\endhead
\bottomrule
\endfoot
RF01 & Autenticação & --- & ---\\
RF02 & Estrutura hierárquica & --- & ---\\
RF03 & Ordenação & --- & ---\\
RF04 & Personalização & --- & ---\\
RF05 & Criar documento & --- & ---\\
RF06 & Criar coleção & --- & ---\\
RF07 & Editar e excluir & --- & ---\\
RF08 & Editor Markdown & --- & ---\\
RF09 & Fórmulas matemáticas & --- & ---\\
RF10 & Realce de sintaxe & --- & ---\\
RF11 & Markdown estendido & --- & ---\\
RF12 & \en{Sidebar} hierárquica & --- & ---\\
RF13 & Barra superior & --- & ---\\
RF14 & Página de documento & --- & ---\\
RF15 & Rascunho local & --- & ---\\
RF16 & Publicar versão & --- & ---\\
RF17 & Invalidação de cache & --- & ---\\
RF18 & Identidade do site & --- & ---\\
RF19 & Paleta de cores & --- & ---\\
RF20 & Encerrar sessão & --- & ---\\
RF21 & Autorização de escrita & --- & ---\\
RF22 & \en{Bootstrap} do administrador & --- & ---\\
RF23 & Configuração inicial & --- & ---\\
RF24 & Rodapé configurável & --- & ---\\
RF25 & Metadados de SEO & --- & ---\\
RF26 & Estado de publicação & --- & ---\\
\end{longtable}
\fonte{Elaborado pelos autores.}

\vitrinependencia{Arquiteto}{Preencher a Tabela de cobertura de requisitos ao final da fase F6, usando somente evidências existentes (casos de teste aprovados, capturas, trechos de código). Nenhum valor deve ser estimado.}

\section{Métricas de Qualidade}\label{sec:res-metricas}

A Tabela~\ref{tab:res-metricas} confronta as metas dos requisitos não funcionais com as medições obtidas no Capítulo~\ref{cap:testes}.

\begin{table}[htb]
\caption{Metas e medições dos requisitos não funcionais}\label{tab:res-metricas}
\small
\begin{tabularx}{\linewidth}{@{}L{1.4cm}YL{2.5cm}L{2.5cm}L{1.9cm}@{}}
\toprule
\textbf{ID} & \textbf{Métrica} & \textbf{Meta} & \textbf{Medido} & \textbf{Atende?}\\
\midrule
RNF02 & LCP (percentil 75, perfil móvel simulado) & $\leq$ 2,5 s & --- & ---\\
RNF02 & Pontuação de desempenho no Lighthouse & $\geq$ 90 & --- & ---\\
RNF05 & Pontuação de acessibilidade & $\geq$ 90 e WCAG AA & --- & ---\\
RNF04 & Pontuação SUS média & $\geq$ 68 & --- & ---\\
RNF04 & Taxa de conclusão das tarefas T1 a T6 & $\geq$ 80\% & --- & ---\\
RNF03 & Vulnerabilidades críticas em dependências & 0 & --- & ---\\
RNF01 & Comandos da instalação limpa (Apêndice~\ref{ap:implantacao}) & Informativo & --- & ---\\
RNF08 & Erros de \en{lint} / Falhas de \en{build} & 0 / 0 & --- & ---\\
--- & Casos de teste aprovados / total & Essenciais: 100\% & --- & ---\\
\bottomrule
\end{tabularx}
\fonte{Elaborado pelos autores. Medições a preencher ao término da verificação.}
\end{table}

\section{Atendimento aos Objetivos}\label{sec:res-objetivos}

\begin{table}[htb]
\caption{Atendimento aos objetivos específicos}\label{tab:res-objetivos}
\small
\begin{tabularx}{\linewidth}{@{}L{1.3cm}YL{3.6cm}L{2.0cm}@{}}
\toprule
\textbf{Objetivo} & \textbf{Resumo} & \textbf{Evidência} & \textbf{Situação}\\
\midrule
OE1 & Publicação independente por clonagem e implantação simples (imagem Docker em VPS) & Apêndice~\ref{ap:implantacao}; CT29 & ---\\
OE2 & Editor Markdown com fórmulas e realce de sintaxe & RF08 a RF11; CT15 a CT19 & ---\\
OE3 & Personalização visual sem novo \en{build} & RF19; CT23, CT24 & ---\\
OE4 & Arquitetura de dados portátil (Prisma e SQLite) & Seção~\ref{sec:arq-dados}; CT28, CT30 & ---\\
OE5 & Autenticação segura (better-auth) & RF01, RF20 a RF22; CT01 a CT06 & ---\\
OE6 & Monólito modular documentado e rastreável & Capítulos~\ref{cap:requisitos} e~\ref{cap:arquitetura}; Apêndice~\ref{ap:adr} & ---\\
OE7 & Processo DDD com agentes e rastreabilidade & Seção~\ref{sec:processo-ddd}; histórico do Git & ---\\
OE8 & Verificação e validação do sistema & Capítulo~\ref{cap:testes} & ---\\
\bottomrule
\end{tabularx}
\fonte{Elaborado pelos autores.}
\end{table}

\section{Discussão}\label{sec:res-discussao}

\subsection{Comparação com os Trabalhos Relacionados}

A discussão deve retomar a Tabela~\ref{tab:trabalhos-relacionados} e examinar, com base nas medições, em que critérios (C1 a C5) o Vitrine de fato se diferencia e em quais se aproxima ou fica aquém das categorias comparadas.

\subsection{O Processo de Desenvolvimento com Agentes}

Esta subseção avalia o método \en{Documentation-Driven Development} com agentes. Devem ser respondidas, com base no histórico do Git e no mural: quantos pedidos de integração foram aprovados de primeira e quantos foram vetados, e por quais motivos; quantas vezes o código divergiu da documentação e como a divergência foi detectada; quantas atualizações de documentação foram solicitadas pelos desenvolvedores; e qual o esforço de revisão do Arquiteto.

\vitrinependencia{Arquiteto}{Coletar, do histórico do Git e do mural, os indicadores do processo (pedidos aprovados e vetados, motivos de veto, divergências detectadas) e analisá-los nesta subseção; os indicadores só podem ser calculados após as primeiras integrações.}

\section{Ameaças à Validade}\label{sec:res-ameacas}

\begin{itemize}
  \item \textbf{Validade de construção:} as métricas de desempenho e de usabilidade são aproximações das qualidades desejadas; o LCP medido em ambiente simulado pode diferir do experimentado por usuários reais;
  \item \textbf{Validade interna:} o autor do sistema também conduz a avaliação, o que pode enviesar a interpretação; mitiga-se com casos de teste automatizados, critérios objetivos e evidências registradas;
  \item \textbf{Validade externa:} o estudo de usabilidade envolve poucos participantes e um contexto institucional específico; os resultados não devem ser generalizados para outros perfis ou áreas;
  \item \textbf{Validade de conclusão:} o volume de dados é pequeno para testes estatísticos; os resultados são descritos de modo qualitativo e com estatística descritiva;
  \item \textbf{Variabilidade dos agentes:} o comportamento de agentes de IA varia entre execuções e versões de modelos; o processo depende de revisão humana e de critérios de aceite objetivos, e os resultados do processo podem não se repetir integralmente.
\end{itemize}

\section{Limitações}

São limitações conhecidas desta versão: administrador único, sem papéis (\decisao{D-13}); ausência de histórico de versões do conteúdo, que se limita ao estado de rascunho ou publicado (\decisao{D-15}); ausência de busca textual e de envio de mídia (\decisao{D-10}); implantação restrita a ambientes com disco persistente, como uma VPS com Docker, sem suporte a plataformas \emph{serverless} (\decisao{D-04}); e o alcance da validação com usuários.

\section{Síntese do Capítulo}

O capítulo consolida os resultados factuais da IC e estrutura a apresentação dos resultados do monólito, com tabelas de cobertura de requisitos, de metas de qualidade e de atendimento aos objetivos. Os valores serão preenchidos após a verificação, sem estimativas.

%% ===================================================================================
%% Capítulo 9 -- Considerações Finais
%% Base: Conclusão do relatório de IC (resultados, limitações, trabalhos futuros),
%% ampliada para o TCC. As conclusões sobre o monólito são escritas pelo Arquiteto
%% somente após o Capítulo 8 estar preenchido.
%% ===================================================================================
\chapter{Considerações Finais}\label{cap:conclusao}

\section{Conclusões}\label{sec:conclusoes}

Este trabalho tratou do desenvolvimento do Vitrine, uma plataforma CMS de código aberto voltada à publicação acadêmica independente. Partiu de um problema concreto --- a dependência de plataformas proprietárias ou a complexidade de CMS tradicionais --- e de uma base já concebida e prototipada na Iniciação Científica \cite{reis2026ic}, que aponta a soberania do autor sobre a sua produção intelectual como princípio orientador.

A contribuição deste trabalho é dupla. No plano do produto, a consolidação do sistema em um monólito modular sobre Next.js, SQLite e better-auth, descrito por requisitos rastreáveis, modelo de dados, contrato de API, estratégia de cache e decisões de projeto (Capítulos~\ref{cap:requisitos} e~\ref{cap:arquitetura}). No plano do processo, a proposta e a aplicação de um método de \en{Documentation-Driven Development} em que uma documentação viva, versionada junto ao código e com um único redator, coordena agentes especializados (Capítulo~\ref{cap:metodologia}).

\vitrinependencia{Arquiteto}{Após o preenchimento do Capítulo 8, reescrever estas conclusões retomando o objetivo geral e cada objetivo específico (OE1 a OE8) com base nas evidências, e responder às perguntas de pesquisa P1 e P2 da Seção 1.3.}

\section{Contribuições}\label{sec:contribuicoes-finais}

\begin{enumerate}
  \item Um artefato de \en{software} de código aberto para publicação científica autônoma, com edição Markdown e fórmulas \LaTeX{}, personalização visual dinâmica e entrega por ISR;
  \item uma especificação documentada e rastreável do sistema, reutilizável por terceiros;
  \item um método de desenvolvimento por documentação e agentes, com regras de governança, fluxo de Git e mural de comunicação, descrito de forma replicável;
  \item um relato de verificação e validação, com as limitações e as ameaças à validade.
\end{enumerate}

\section{Limitações}\label{sec:limitacoes-finais}

Além das limitações registradas na Seção~\ref{sec:res-ameacas}, o relatório de IC já apontava que a modelagem da interface se restringira a protótipos de baixa fidelidade, em prol da funcionalidade e da usabilidade, e que a autenticação poderia ser expandida para múltiplos níveis de permissão em cenários colaborativos \cite{reis2026ic}. Ambos os pontos permanecem como oportunidades de evolução.

\section{Trabalhos Futuros}\label{sec:trabalhos-futuros}

Do relatório de IC \cite{reis2026ic} permanecem como propostas de continuidade:

\begin{itemize}
  \item produzir documentação técnica detalhada e tutoriais em vídeo que demonstrem o processo de clonagem do repositório e de implantação --- considerando a ressalva sobre plataformas \emph{serverless} e a persistência do SQLite (\decisao{D-04});
  \item lançar o site oficial do projeto, denominado Vitrine, que servirá como repositório central para tutoriais, documentação e instâncias de demonstração;
  \item explorar a integração com repositórios de dados abertos para a citação automática de publicações;
  \item expandir a biblioteca de componentes visuais, permitindo maior customização da interface;
  \item implementar a renderização nativa de gráficos de funções de uma variável (gráficos 2D).
\end{itemize}

Da consolidação do monólito surgem novas linhas de evolução, já mapeadas como decisões em aberto (Apêndice~\ref{ap:decisoes}): papéis e permissões com o \en{plugin} \texttt{admin} do better-auth (\decisao{D-13}); estados de publicação e histórico de versões (\decisao{D-15}); envio e gerenciamento de mídia (\decisao{D-10}); busca textual no conteúdo; internacionalização (\req{RNF12}); exportação e importação do conteúdo em Markdown, para reforçar a soberania dos dados; e temas pré-definidos.

\section{Considerações Finais}

A experiência de desenvolver o Vitrine reafirma a premissa de que sistemas de informação devem atuar como facilitadores da disseminação de conhecimento, e não como obstáculos técnicos \cite{rainer2019intro}. O projeto entrega uma base sólida e moderna que elimina o trabalho repetitivo de configuração estrutural, permitindo que o foco permaneça na produção de conteúdo qualificado. O Vitrine busca ser não apenas uma ferramenta técnica, mas um instrumento em prol da autonomia digital e da democratização do saber científico na rede.


% ----------------------------------------------------------
% ELEMENTOS PÓS-TEXTUAIS
% ----------------------------------------------------------
\postextual
\bookmarksetup{startatroot}

\bibliographystyle{abntex2-num}
\bibliography{references}

\begin{apendicesenv}
\partapendices
%% ===================================================================================
%% Apêndice A -- Dicionário de Dados
%% Espelha prisma/schema.prisma. Quando o agente Dados alterar o esquema e o merge for
%% aprovado, o Arquiteto atualiza esta tabela no mesmo dia (ver ARQUITETO.md).
%% Tipos: Prisma -> SQLite (String=TEXT, Int=INTEGER, Boolean=BOOLEAN, DateTime=DATETIME).
%% ===================================================================================
\chapter{Dicionário de Dados}\label{ap:dicionario}

Este apêndice descreve as tabelas do banco de dados conforme o esquema do Prisma em vigor (Apêndice~\ref{ap:schema}). Os nomes de tabela vêm da diretiva \texttt{@@map}; os nomes de coluna são os campos dos modelos. \emph{PK} indica chave primária; \emph{FK}, chave estrangeira; ``?'' após o tipo indica valor opcional (nulo permitido). As chaves estrangeiras têm exclusão em cascata, exceto \texttt{contents.authorId}, que usa \texttt{Restrict}. \texttt{ContentStatus} é um enumerado do Prisma (\texttt{DRAFT} ou \texttt{PUBLISHED}) que o SQLite não impõe no banco; a validação ocorre no Prisma e na aplicação.

\begin{longtable}{@{}L{3.9cm}L{1.8cm}L{3.3cm}L{5.0cm}@{}}
\caption{Dicionário de dados}\label{tab:dicionario}\\
\toprule
\textbf{Coluna} & \textbf{Tipo} & \textbf{Restrição / padrão} & \textbf{Descrição}\\
\midrule
\endfirsthead
\toprule
\textbf{Coluna} & \textbf{Tipo} & \textbf{Restrição / padrão} & \textbf{Descrição}\\
\midrule
\endhead
\bottomrule
\endfoot
\multicolumn{4}{@{}l}{\textbf{Tabela \texttt{user} (modelo \texttt{User}) --- gerenciada pelo better-auth}}\\[2pt]
\texttt{id} & String & PK; \texttt{cuid()} & Identificador do usuário\\
\texttt{name} & String & Obrigatório & Nome de exibição\\
\texttt{email} & String & Único; obrigatório & E-mail, usado como login\\
\texttt{emailVerified} & Boolean & Obrigatório & Indica e-mail verificado\\
\texttt{image} & String? & Opcional & URL da foto de perfil\\
\texttt{createdAt} & DateTime & \texttt{now()} & Data de criação\\
\texttt{updatedAt} & DateTime & \texttt{@updatedAt} & Data da última alteração\\
\midrule
\multicolumn{4}{@{}l}{\textbf{Tabela \texttt{session} (modelo \texttt{Session}) --- gerenciada pelo better-auth}}\\[2pt]
\texttt{id} & String & PK; \texttt{cuid()} & Identificador da sessão\\
\texttt{token} & String & Único; obrigatório & Identificador opaco guardado no \en{cookie}\\
\texttt{expiresAt} & DateTime & Obrigatório & Instante de expiração (RN12)\\
\texttt{ipAddress} & String? & Opcional & Endereço IP de criação\\
\texttt{userAgent} & String? & Opcional & Navegador de criação\\
\texttt{userId} & String & FK $\rightarrow$ \arq{user.id} & Usuário dono da sessão\\
\texttt{createdAt}, \texttt{updatedAt} & DateTime & \texttt{now()}; \texttt{@updatedAt} & Carimbos de tempo\\
\midrule
\multicolumn{4}{@{}l}{\textbf{Tabela \texttt{account} (modelo \texttt{Account}) --- gerenciada pelo better-auth}}\\[2pt]
\texttt{id} & String & PK; \texttt{cuid()} & Identificador da conta\\
\texttt{accountId} & String & Obrigatório & Identificador da conta no provedor\\
\texttt{providerId} & String & Obrigatório & Identificador do provedor de login\\
\texttt{userId} & String & FK $\rightarrow$ \arq{user.id} & Usuário dono da conta\\
\texttt{password} & String? & Opcional & \en{Hash} da senha (login por e-mail e senha)\\
\texttt{accessToken}, \texttt{refreshToken}, \texttt{idToken} & String? & Opcionais & \en{Tokens} de provedores externos\\
\arq{accessTokenExpiresAt}, \arq{refreshTokenExpiresAt} & DateTime? & Opcionais & Validade dos \en{tokens}\\
\texttt{scope} & String? & Opcional & Escopos autorizados\\
\texttt{createdAt}, \texttt{updatedAt} & DateTime & \texttt{now()}; \texttt{@updatedAt} & Carimbos de tempo\\
\midrule
\multicolumn{4}{@{}l}{\textbf{Tabela \texttt{verification} (modelo \texttt{Verification}) --- gerenciada pelo better-auth}}\\[2pt]
\texttt{id} & String & PK; \texttt{cuid()} & Identificador\\
\texttt{identifier} & String & Obrigatório & Alvo da verificação (por exemplo, e-mail)\\
\texttt{value} & String & Obrigatório & Valor a conferir\\
\texttt{expiresAt} & DateTime & Obrigatório & Validade\\
\texttt{createdAt}, \texttt{updatedAt} & DateTime? & \texttt{now()}; \texttt{@updatedAt} & Carimbos de tempo\\
\midrule
\multicolumn{4}{@{}l}{\textbf{Tabela \texttt{contents} (modelo \texttt{Content}) --- conteúdo hierárquico}}\\[2pt]
\texttt{id} & String & PK; \texttt{cuid()} & Identificador do conteúdo\\
\texttt{title} & String & Obrigatório & Título exibido no menu e na página\\
\texttt{slug} & String & Obrigatório; único entre irmãos & Trecho amigável da URL (RN03; índices na Tabela~\ref{tab:indices})\\
\texttt{body} & String? & Opcional & Corpo em Markdown (vazio para coleções)\\
\texttt{type} & String & Obrigatório & \texttt{DOCUMENT} ou \texttt{COLLECTION} (RN02)\\
\texttt{status} & ContentStatus & \texttt{PUBLISHED} & Estado de publicação: \texttt{DRAFT} (rascunho) ou \texttt{PUBLISHED} (RN15, RN16; D-15)\\
\texttt{order} & Int & \texttt{0} & Ordem entre irmãos (RN05)\\
\texttt{parentId} & String? & FK $\rightarrow$ \arq{contents.id} & Conteúdo pai; nulo indica raiz\\
\texttt{authorId} & String & FK $\rightarrow$ \arq{user.id}; \texttt{Restrict} & Autor (RN13); excluir um usuário que tenha conteúdo é bloqueado (RNF10)\\
\texttt{createdAt}, \texttt{updatedAt} & DateTime & \texttt{now()}; \texttt{@updatedAt} & Carimbos de tempo\\
\midrule
\multicolumn{4}{@{}l}{\textbf{Tabela \arq{configurations} (modelo \texttt{Configuration}) --- registro único}}\\[2pt]
\texttt{id} & String & PK; padrão \texttt{default} & Identificador fixo do registro único (RN08); acesso sempre por \emph{upsert}\\
\texttt{siteName} & String & \texttt{"Vitrine"} & Nome do site\\
\arq{siteDescription} & String? & Opcional & Descrição do site\\
\texttt{primaryColor} & String & \texttt{"\#3b82f6"} & Cor primária (RN09)\\
\arq{secondaryColor} & String & \texttt{"\#1e293b"} & Cor secundária\\
\texttt{accentColor} & String & \texttt{"\#f59e0b"} & Cor de destaque\\
\texttt{logoUrl}, \texttt{faviconUrl} & String? & Opcionais & Logotipo e \en{favicon} (D-10)\\
\arq{discordUrl}, \arq{githubUrl}, \arq{youtubeUrl}, \arq{instagramUrl}, \arq{linkedinUrl}, \arq{twitchUrl} & String? & Opcionais; somente http(s) (RN17) & Links sociais do rodapé (RF24, D-19); vazio = ícone omitido\\
\texttt{updatedAt} & DateTime & \texttt{@updatedAt} & Data da última alteração\\
\end{longtable}
\fonte{Elaborado pelos autores, a partir de \arq{prisma/schema.prisma}.}

A Tabela~\ref{tab:indices} lista os índices de \texttt{contents}, criados pela migração \arq{20261001034824_init}. O índice parcial da raiz não é modelado pelo Prisma: foi escrito à mão no SQL da migração (Seção~\ref{sec:arq-indice-parcial}).

\begin{table}[htb]
\caption{Índices da tabela \texttt{contents}}\label{tab:indices}
\small
\begin{tabularx}{\linewidth}{@{}L{4.6cm}L{4.4cm}Y@{}}
\toprule
\textbf{Índice} & \textbf{Colunas} & \textbf{Finalidade}\\
\midrule
\arq{contents_parentId_slug_key} & único: (\texttt{parentId}, \texttt{slug}) & RN03: \en{slug} único entre irmãos de um mesmo pai\\
\addlinespace
\arq{contents_root_slug_key} & único parcial: \texttt{slug} \texttt{WHERE parentId IS NULL} & RN03 na raiz, onde o SQLite trata \texttt{NULL} como distinto (D-05)\\
\addlinespace
\arq{contents_parentId_order_idx} & (\texttt{parentId}, \texttt{order}) & RNF02, RN05: consulta da árvore e ordenação dos irmãos\\
\bottomrule
\end{tabularx}
\fonte{Elaborado pelos autores, a partir da migração inicial.}
\end{table}

%% ===================================================================================
%% Apêndice B -- Esquema do Prisma (incluído diretamente do repositório)
%% Os listings abaixo leem os arquivos reais: nunca ficam desatualizados em relação ao
%% que está na main. Não copie o conteúdo para este arquivo.
%% ===================================================================================
\chapter{Esquema e Configuração do Prisma}\label{ap:schema}

Este apêndice reproduz, diretamente do repositório, o esquema do Prisma e o arquivo de configuração da CLI em vigor na data de compilação deste documento. O conteúdo é incluído por referência aos arquivos, de modo que sempre reflete a \texttt{main}. O estado de conformidade com as práticas do Prisma 7 está na Tabela~\ref{tab:prisma7} e as observações sobre o modelo, na Tabela~\ref{tab:obs-schema}. Até o módulo M15 ser integrado, o esquema reproduzido é o \emph{anterior} ao escopo multiusuário; o modelo-alvo (usuário com papel e \en{handle}, \texttt{appearances}, \texttt{notifications}, dono único por árvore e índices por dono) está no Apêndice~\ref{ap:dicionario} (Tabelas~\ref{tab:dicionario-alvo} e~\ref{tab:indices-alvo}) e na Seção~\ref{sec:arq-dados}, e substitui estes arquivos no merge. A migração do M15 terá SQL escrito à mão (índice parcial por dono, \texttt{UPDATE} condicional das cores), que o Prisma não enxerga.

\lstinputlisting[language=VPrisma, caption={Conteúdo de prisma/schema.prisma}, label={lst:schema-prisma}]{../../prisma/schema.prisma}

\lstinputlisting[language=VJS, caption={Conteúdo de prisma.config.ts}, label={lst:prisma-config}]{../../prisma.config.ts}

A migração inicial, reproduzida a seguir, contém o SQL gerado pelo Prisma e, ao final, o índice parcial escrito à mão (Seção~\ref{sec:arq-indice-parcial}).

\lstinputlisting[caption={Conteúdo de prisma/migrations/20261001034824\_init/migration.sql}, label={lst:migracao-init}]{../../prisma/migrations/20261001034824_init/migration.sql}

%% ===================================================================================
%% Apêndice C -- Fontes dos Diagramas
%% O catálogo completo e as instruções de renderização estão em doc/diagrams/README.md.
%% Aqui entram apenas dois fontes representativos, incluídos diretamente do repositório.
%% ===================================================================================
\chapter{Fontes dos Diagramas}\label{ap:diagramas}

Os diagramas deste documento são mantidos como código: PlantUML para diagramas UML e de arquitetura, e DBML para o modelo de dados. Os fontes vivem em \arq{doc/diagrams/src/} e são a fonte da verdade visual; as imagens renderizadas ficam em \arq{doc/tex/figuras/diagramas/}. O catálogo completo --- um arquivo por figura, com o nome do arquivo igual ao identificador da figura --- e as instruções de renderização estão em \arq{doc/diagrams/README.md}. A Tabela~\ref{tab:catalogo-diagramas} resume o catálogo.

\begin{table}[htb]
\caption{Catálogo dos diagramas-fonte}\label{tab:catalogo-diagramas}
\small
\begin{tabularx}{\linewidth}{@{}L{4.6cm}L{2.6cm}Y@{}}
\toprule
\textbf{Arquivo em \arq{doc/diagrams/src/}} & \textbf{Formato} & \textbf{Figura}\\
\midrule
\arq{uc-visao-geral.puml} & PlantUML & Casos de uso (Figura~\ref{fig:uc-visao-geral})\\
\arq{er-vitrine.dbml} & DBML & Entidade-relacionamento (Figura~\ref{fig:er-vitrine})\\
\arq{classes-dominio.puml} & PlantUML & Modelo de domínio (Figura~\ref{fig:classes-dominio})\\
\arq{arq-contexto.puml}, \arq{arq-conteineres.puml} & PlantUML & Contexto e contêineres (Figuras~\ref{fig:arq-contexto} e~\ref{fig:arq-conteineres})\\
\arq{arq-componentes.puml}, \arq{arq-camadas.puml} & PlantUML & Componentes e camadas (Figuras~\ref{fig:arq-componentes} e~\ref{fig:arq-camadas})\\
\arq{seq-autenticacao.puml} & PlantUML & Autenticação (Figura~\ref{fig:seq-autenticacao})\\
\arq{seq-leitura-isr.puml} & PlantUML & Leitura com ISR (Figura~\ref{fig:seq-leitura-isr})\\
\arq{seq-publicacao-conteudo.puml} & PlantUML & Publicação (Figura~\ref{fig:seq-publicacao-conteudo})\\
\arq{seq-personalizacao-tema.puml} & PlantUML & Personalização (Figura~\ref{fig:seq-personalizacao-tema})\\
\arq{seq-reordenacao-arvore.puml} & PlantUML & Reordenação (Figura~\ref{fig:seq-reordenacao-arvore})\\
\arq{nav-mapa-do-site.puml} & PlantUML & Mapa de rotas (Figura~\ref{fig:nav-mapa-do-site})\\
\arq{impl-implantacao.puml} & PlantUML & Implantação (Figura~\ref{fig:impl-implantacao})\\
\arq{fluxo-ddd-agentes.puml}, \arq{fluxo-git-mural.puml} & PlantUML & Processo (Figuras~\ref{fig:fluxo-ddd-agentes} e~\ref{fig:fluxo-git-mural})\\
\bottomrule
\end{tabularx}
\fonte{Elaborado pelos autores.}
\end{table}

O conteúdo integral de cada fonte pode ser lido no repositório; para evitar duplicação e divergência, este apêndice não o reproduz. A regra é: \textbf{o arquivo em \arq{doc/diagrams/src/} prevalece}, e a figura do Capítulo correspondente deve ser sempre uma renderização do fonte vigente.

%% ===================================================================================
%% Apêndice D -- Guia de Implantação e Operação
%% ESBOÇO normativo escrito antes do código. Os agentes Core e Dados devem validar cada
%% passo numa instalação limpa (CT29) e pedir ao Arquiteto, pelo mural, as correções.
%% ===================================================================================
\chapter{Guia de Implantação e Operação}\label{ap:implantacao}

Este apêndice descreve como instalar, executar, atualizar e proteger uma instância do Vitrine. Destina-se ao Implantador (Seção~\ref{sec:atores}). Os comandos foram redigidos a partir da especificação e da documentação das ferramentas; devem ser validados em uma instalação limpa (caso de teste CT29) e corrigidos conforme a implementação real.

\vitrinependencia{Core}{Validar este guia em uma instalação limpa, corrigindo comandos e nomes de variáveis, e confirmar com o Arquiteto pelo mural o texto final (caso CT29).}

\section{Pré-requisitos}

\begin{itemize}
  \item Servidor Linux, máquina virtual ou contêiner com \textbf{Node.js 20.9 ou superior} e \texttt{npm};
  \item \textbf{disco persistente} para o arquivo do banco (um volume dedicado, em contêineres);
  \item \textbf{HTTPS} em produção, tipicamente por um \en{proxy} reverso com TLS;
  \item acesso ao repositório Git do projeto.
\end{itemize}

\section{Variáveis de Ambiente}

\begin{table}[htb]
\caption{Variáveis de ambiente (nomes; nunca versionar valores)}\label{tab:env}
\small
\begin{tabularx}{\linewidth}{@{}L{4.1cm}L{2.2cm}Y@{}}
\toprule
\textbf{Variável} & \textbf{Obrigatória} & \textbf{Descrição}\\
\midrule
\arq{DATABASE_URL} & Sim & Localização do arquivo SQLite, no formato \texttt{file:CAMINHO}. Em produção, usar caminho absoluto em volume persistente (por exemplo, \arq{file:/data/vitrine.db})\\
\addlinespace
\arq{BETTER_AUTH_SECRET} & Sim (nome a confirmar) & Segredo usado pelo better-auth para assinar/criptografar dados de sessão; longo e aleatório\\
\addlinespace
\arq{BETTER_AUTH_URL} & Sim (nome a confirmar) & URL pública base do site (por exemplo, \arq{https://meusite.exemplo})\\
\addlinespace
\arq{NODE_ENV} & Sim & Deve ser \texttt{production} na execução em produção\\
\addlinespace
\arq{PORT} & Não & Porta do servidor Next.js; o padrão é 3000\\
\bottomrule
\end{tabularx}
\fonte{Elaborado pelos autores. Os nomes das variáveis do better-auth devem ser confirmados pelo agente Core na documentação da versão instalada.}
\end{table}

\section{Instalação}

\begin{lstlisting}[language=VSh, caption={Esboço: instalação em servidor}, label={lst:instalacao}]
# 1. Obter o código e instalar as dependências
git clone <url-do-repositorio> vitrine && cd vitrine
npm ci

# 2. Configurar o ambiente (.env nunca é versionado)
cp .env.example .env
# editar .env: DATABASE_URL, segredo e URL do better-auth

# 3. Aplicar as migrações e criar os dados iniciais
npx prisma migrate deploy
npx prisma db seed        # configuração padrão + primeiro administrador (RF22, RF23)

# 4. Compilar e iniciar
npm run build
npm run start             # http://localhost:3000
\end{lstlisting}

Em produção, o processo deve ser supervisionado por um gerenciador (por exemplo, \texttt{systemd} ou um gerenciador de processos Node.js) ou executado em contêiner com reinício automático. Deve haver \textbf{uma única instância} (Seção~\ref{sec:arq-isr}).

\section{Proxy Reverso e HTTPS}

Um \en{proxy} reverso com TLS recebe as conexões públicas e as encaminha à porta interna do Node.js. Exemplo mínimo com o servidor Caddy, que obtém e renova certificados automaticamente:

\begin{lstlisting}[caption={Esboço: Caddyfile para o Vitrine}, label={lst:caddy}]
meusite.exemplo {
    reverse_proxy 127.0.0.1:3000
}
\end{lstlisting}

\section{Cópia de Segurança e Restauração}\label{sec:backup}

O banco é um único arquivo; porém, com o modo WAL, a cópia de arquivos com o sistema em execução pode ser inconsistente. A forma recomendada é gerar uma cópia consistente com o próprio SQLite, por exemplo:

\begin{lstlisting}[language=VSh, caption={Esboço: cópia de segurança consistente do SQLite}, label={lst:backup}]
# Requer a CLI sqlite3 (SQLite 3.27 ou superior para VACUUM INTO)
sqlite3 /data/vitrine.db "VACUUM INTO '/backup/vitrine-$(date +%F).db'"
\end{lstlisting}

Para restaurar: parar o servidor, substituir o arquivo do banco pela cópia e iniciar novamente. A rotina deve ser agendada (por exemplo, diariamente) e as cópias guardadas fora do servidor (RNF10).

\section{Atualização}

Antes de atualizar, faça uma cópia de segurança. Em seguida: obter o novo código (\texttt{git pull}), reinstalar dependências (\texttt{npm ci}), aplicar migrações (\texttt{npx prisma migrate deploy}), recompilar (\texttt{npm run build}) e reiniciar o serviço. A recompilação gera um novo conjunto de páginas pré-renderizadas, o que renova as páginas servidas pelo cache da aplicação.

\section{Solução de Problemas}

\begin{longtable}{@{}L{4.2cm}L{4.5cm}L{5.3cm}@{}}
\caption{Problemas comuns de implantação}\label{tab:problemas}\\
\toprule
\textbf{Sintoma} & \textbf{Causa provável} & \textbf{Ação}\\
\midrule
\endfirsthead
\toprule
\textbf{Sintoma} & \textbf{Causa provável} & \textbf{Ação}\\
\midrule
\endhead
\bottomrule
\endfoot
Erro de banco não encontrado & Caminho relativo em \arq{DATABASE_URL} resolvido de forma diferente pela CLI e pelo servidor & Usar caminho absoluto em produção\\[1pt]
Login não persiste & URL base incorreta ou \en{cookie} \texttt{Secure} sem HTTPS & Conferir a URL do better-auth e o TLS do \en{proxy}\\[1pt]
Alteração não aparece no site público & Mais de uma instância (cache por instância) ou falha na revalidação & Manter uma única instância; verificar o cabeçalho \arq{x-nextjs-cache}\\[1pt]
Falha de escrita no banco & Permissão de escrita ausente na pasta do arquivo & Ajustar o dono e as permissões do volume\\[1pt]
Dados perdidos após reinício em plataforma \en{serverless} & Sistema de arquivos efêmero & Usar servidor com disco persistente ou banco remoto (D-04)\\[1pt]
\end{longtable}
\fonte{Elaborado pelos autores.}

%% ===================================================================================
%% Apêndice E -- Registros de Decisão Arquitetural (ADR)
%% Formato: Nygard (2011). Uma decisão aceita só muda por uma NOVA ADR que a substitui
%% (status "Substituída por ADR-0xx"); nunca se apaga uma ADR. Questões ainda abertas
%% ficam no Apêndice F. Dono: Arquiteto/Redator.
%% ===================================================================================
\chapter{Registros de Decisão Arquitetural}\label{ap:adr}
\begingroup\small

Este apêndice reúne as decisões arquiteturais aceitas, no formato de \en{Architecture Decision Records} \cite{nygard2011adr}: contexto, decisão, consequências e alternativas. Uma decisão só é alterada por uma nova ADR que a substitua; as questões ainda não decididas estão no Apêndice~\ref{ap:decisoes}.

\adrentry{ADR-001}{Monólito modular com Next.js (App Router)}{Aceita.}
{O Vitrine é um sistema de um único \en{tenant}, com tráfego moderado, implantado pelo próprio docente; a equipe é pequena e há agentes de desenvolvimento em paralelo. O requisito de implantação simples (\req{RNF01}) e o baixo custo operacional são prioritários.}
{Adotar uma única aplicação Next.js (App Router), organizada internamente em módulos por domínio --- Dados, Core e UI --- com dependências unidirecionais (Figuras~\ref{fig:arq-componentes} e~\ref{fig:arq-camadas}).}
{Positivas: implantação e depuração simples, transações locais, uma única unidade de versionamento. Negativas: escalabilidade horizontal limitada (cache ISR por instância, SQLite em arquivo); a interface e a API são implantadas juntas.}
{\en{Front-end} separado com API independente; microsserviços; CMS \en{headless} de terceiros. Descartadas pelo custo de operação e de coordenação desproporcional ao escopo \cite{newman2019monolith,fowler2015monolithfirst}.}

\adrentry{ADR-002}{SQLite como banco de dados embutido}{Aceita.}
{O objetivo OE4 pede uma arquitetura de dados portátil, sem servidor de banco de dados, com baixo custo para o usuário final. O perfil de uso é de leitura predominante, com poucas escritas realizadas por um administrador.}
{Usar o SQLite, com o banco em um único arquivo em volume persistente e, preferencialmente, em modo WAL.}
{Positivas: sem servidor a administrar, cópia de segurança simples, custo zero. Negativas: escrita serializada; necessidade de disco persistente (incompatível com sistemas de arquivos efêmeros, risco R01, tratado pela \adr{ADR-011}); ausência de acesso remoto nativo; valores de enumeração não impostos pelo banco \cite{sqliteabout,sqlitewal}.}
{PostgreSQL ou MySQL (exigem servidor); libSQL/Turso (banco remoto compatível com SQLite, via possível para plataformas \emph{serverless}, descartada nesta versão pela \adr{ADR-011}).}

\adrentry{ADR-003}{Prisma ORM 7 como camada de acesso a dados}{Aceita.}
{É necessário um modelo de dados declarativo, versionável e com migrações, e um cliente consistente para o SQLite. O projeto usa o Prisma 7, cujas práticas diferem das versões anteriores (configuração em arquivo próprio, novo gerador, \en{driver adapter}).}
{Adotar o Prisma ORM 7. Somente \arq{src/lib/db.js} instancia o cliente; o restante do código usa as funções exportadas por ele.}
{Positivas: esquema único como fonte do modelo, migrações e cliente tipado. Negativas: a adequação à versão 7 gerou trabalho inicial, resolvido na \adr{ADR-014}; o Prisma acrescenta uma camada de abstração entre a aplicação e o SQL.}
{Drizzle ORM; SQL direto com um \en{driver} de SQLite (mais controle, menos produtividade); uso apenas do \en{query builder} interno do better-auth.}

\adrentry{ADR-004}{better-auth para autenticação, em substituição ao Auth.js}{Aceita.}
{O relatório de IC previu o Auth.js \cite{reis2026ic}. Na consolidação do monólito, o repositório adotou o better-auth: a dependência está declarada, o esquema de dados contém as tabelas \texttt{user}, \texttt{session}, \texttt{account} e \texttt{verification} no formato do better-auth, e o pacote instalado oferece adaptador para Prisma e integração com Next.js.}
{Adotar o better-auth 1.7.x com adaptador do Prisma, autenticação por e-mail e senha, sessões persistidas em banco e cadastro público desabilitado após o \en{bootstrap}.}
{Positivas: esquema alinhado ao que já existe; integração com Next.js; \en{plugins} para evolução (papéis, autenticação de dois fatores). Negativas: biblioteca mais recente, exigindo acompanhar mudanças de API; todas as menções ao Auth.js do relatório de IC passam a ser históricas.}
{Auth.js (prevista na IC); implementação própria, descartada pelo risco de segurança. \textbf{Motivação original a confirmar com o autor}: esta ADR registra o estado do repositório e justificativas técnicas verificáveis; a motivação da troca deve ser completada pelo autor.}

\adrentry{ADR-005}{ISR com revalidação sob demanda}{Aceita.}
{As páginas públicas precisam de desempenho de página estática (\req{RNF02}) sem exigir novo \en{build} a cada alteração de conteúdo ou de tema (OE3).}
{Pré-renderizar as páginas públicas com \arq{generateStaticParams}; definir \texttt{revalidate} alto como rede de segurança; invalidar sob demanda com \arq{revalidatePath} nas mutações (regra RN10); não habilitar Cache Components.}
{Positivas: desempenho e atualização imediata na próxima visita. Negativas: suporta apenas o \en{runtime} Node.js; o cache é por instância (exige instância única); a invalidação do \en{layout} raiz regenera todas as páginas \cite{nextjsdocs}.}
{SSR a cada requisição; SSG puro com novo \en{build}; \texttt{revalidateTag} com granularidade por etiqueta (reavaliar se o volume de conteúdo crescer); Cache Components.}

\adrentry{ADR-006}{Markdown com \texttt{unified}, KaTeX e sem HTML bruto}{Aceita.}
{O objetivo OE2 exige escrita científica: fórmulas \LaTeX{}, código com realce e tabelas. A segurança exige que o conteúdo do administrador nunca execute \en{scripts} no navegador do leitor.}
{Renderizar com \arq{react-markdown}, \arq{remark-math} e \arq{rehype-katex}, acrescentando GFM e realce de sintaxe; proibir HTML bruto.}
{Positivas: texto legível e portátil; segurança por padrão; matemática nativa. Negativas: sem HTML arbitrário; dependência de \en{plugins} do ecossistema \texttt{unified} \cite{unifiedjs,reactmarkdown,katex}.}
{MDX (permite código executável, risco de segurança); editor visual que gera HTML; renderização apenas no cliente.}

\adrentry{ADR-007}{Tematização por variáveis CSS e Tailwind CSS 4}{Aceita.}
{A personalização visual deve refletir-se sem novo \en{build} (OE3).}
{Mapear tokens do Tailwind para variáveis CSS (\texttt{@theme inline}) e definir os valores no atributo \texttt{style} do elemento raiz, a partir da configuração do banco; validar cores como \texttt{\#RRGGBB}.}
{Positivas: mudança de tema sem recompilar CSS; baixa complexidade. Negativas: exige disciplina para usar sempre os tokens; convivência com o modo escuro é uma questão em aberto (\decisao{D-17}) \cite{tailwinddocs,mdncssvars}.}
{CSS-in-JS em tempo de execução; gerar CSS no \en{build} (exigiria recompilar); temas fixos predefinidos.}

\adrentry{ADR-008}{JavaScript com JSDoc, sem TypeScript}{Aceita.}
{O projeto foi iniciado em JavaScript (presença de \arq{jsconfig.json}) e precisa minimizar a fricção entre agentes.}
{Manter JavaScript (módulos ES), documentando as funções públicas de \arq{src/lib} com JSDoc. Reavaliar TypeScript em versões futuras.}
{Positivas: menor fricção e configuração. Negativas: menos garantias estáticas; erros de tipo aparecem em testes ou em execução, o que aumenta a importância dos testes de unidade e de integração \cite{bierman2014typescript,andreasen2017survey}.}
{TypeScript desde já.}

\adrentry{ADR-009}{Documentation-Driven Development com \LaTeX{} e diagramas como código}{Aceita.}
{O desenvolvimento por agentes em paralelo exige um contrato único e versionado; o trabalho é um TCC e a documentação também é o produto acadêmico.}
{Manter a documentação oficial em \arq{doc/tex} (\LaTeX{}, base abntex2) e os diagramas em \arq{doc/diagrams/src} (PlantUML e DBML), com um único redator (Arquiteto), figuras com \en{placeholder} automático e mapa de pendências.}
{Positivas: documentação revisável como código e pronta para a entrega; coerência entre agentes. Negativas: custo de redação e de manutenção pelo Arquiteto; dependência de ferramentas de composição e de renderização \cite{parnas1986rational,writethedocs}.}
{Documentação em \en{wiki} ou Markdown solto; documentos de texto proprietários; documentar só ao final.}

\adrentry{ADR-010}{Fluxo Git por domínios com mural na \texttt{main}}{Aceita.}
{Vários agentes trabalham simultaneamente e precisam de um canal de coordenação e de um ponto único de integração. O repositório não possui, atualmente, um remoto configurado.}
{Desenvolver apenas em ramificações por domínio; proibir código na \texttt{main}; usar \arq{MURAL_DE_AGENTES.md}, editado somente na \texttt{main}, como canal de pedidos de integração e de comunicação; o Arquiteto revisa, integra ou veta.}
{Positivas: rastreabilidade no histórico do Git; conflitos de comunicação evitados; integração controlada. Negativas: o Arquiteto é ponto único de revisão; trocar de ramificação exige árvore de trabalho limpa \cite{chacon2014progit,conventionalcommits}.}
{Pedidos de integração em plataforma Git hospedada (sem remoto hoje); trabalho direto na \texttt{main}.}

\adrentry{ADR-011}{Implantação por imagem Docker em VPS tradicional, com volume persistente}{Aceita (decisão do Engenheiro Chefe em 30/09/2026, pedido P-001; resolve \decisao{D-04}).}
{O SQLite exige disco persistente (\adr{ADR-002}), e o sistema de arquivos de plataformas \emph{serverless} é efêmero (risco R01); o relatório de IC citava a Vercel como exemplo de implantação \cite{reis2026ic}. O Implantador é, em geral, o próprio docente com apoio técnico, e a ISR exige \en{runtime} Node.js e instância única (\adr{ADR-005}).}
{Distribuir o Vitrine como \textbf{imagem Docker} executada em uma \textbf{VPS tradicional} (Linux com Docker Engine). O arquivo SQLite reside em um \textbf{volume persistente} montado no contêiner (por exemplo, em \arq{/data}); executa-se uma única instância; um \en{proxy} reverso fornece TLS; os segredos entram por variáveis de ambiente, fora da imagem; as cópias de segurança usam \texttt{VACUUM INTO}. Plataformas \emph{serverless} não são alvo desta versão.}
{Positivas: persistência garantida do banco; ambiente reprodutível e documentável por \texttt{Dockerfile} e \texttt{docker-compose.yml}; coerência com a ISR e com o modelo de instância única \cite{dockerdocs}. Negativas: exige uma VPS e noções básicas de Docker pelo Implantador (o requisito \req{RNF01} é preservado por poucos comandos de \en{compose}); o módulo nativo \arq{better-sqlite3} precisa ser instalado no \en{build} da imagem; o volume e as cópias de segurança são responsabilidade do operador; a CLI do Prisma (hoje em \texttt{devDependencies}) é necessária para aplicar migrações, o que deve ser resolvido no módulo M13.}
{Plataformas \emph{serverless} (Vercel) com banco remoto compatível com SQLite, como libSQL/Turso (descartada nesta versão: troca do \en{driver}, serviço externo e maior complexidade); PaaS com disco persistente; instalação direta com Node.js sem contêiner (continua possível, mas não é o alvo documentado).}

\adrentry{ADR-012}{Estado de publicação do conteúdo: DRAFT e PUBLISHED, com padrão PUBLISHED}{Aceita (decisão do Engenheiro Chefe em 30/09/2026; resolve \decisao{D-15}; integrada pelo merge M-001).}
{Os protótipos da IC trazem o comando ``Enviar nova versão'' e a possibilidade de preparar conteúdo antes de exibi-lo. A recomendação inicial do Arquiteto era a publicação imediata, com rascunho apenas local no navegador (\req{RF15}).}
{Adicionar o campo \texttt{Content.status}, do tipo enumerado \texttt{ContentStatus} (\texttt{DRAFT} ou \texttt{PUBLISHED}), com padrão \texttt{PUBLISHED}. Um conteúdo é visível ao público somente se ele e todos os seus ancestrais estiverem publicados (regra RN15). A API aceita e devolve o estado; alterá-lo é uma mudança estrutural para efeito de revalidação (RN10 e RN16). O histórico de versões continua fora do escopo.}
{Positivas: permite preparar e revisar conteúdo no servidor sem expô-lo. Negativas: toda leitura pública deve filtrar por estado (risco de vazamento de rascunhos, coberto pelo caso de teste CT31); a geração estática e o mapa do site devem excluir rascunhos; como o padrão é \texttt{PUBLISHED}, um conteúdo criado sem escolha explícita fica público, e a interface deve tornar o estado visível; o SQLite não impõe o enumerado no banco, que é validado pelo Prisma e pela aplicação.}
{Publicação imediata sem campo de estado (recomendação inicial); estado mais tabela de versões (descartada por exceder o escopo).}

\adrentry{ADR-013}{Links sociais do rodapé como campos individuais opcionais}{Aceita (decisão do Engenheiro Chefe em 30/09/2026; resolve \decisao{D-19}; integrada pelo merge M-001).}
{O rodapé dos protótipos exibe ícones de redes sociais, mas o esquema original não previa onde guardar os endereços. A recomendação inicial era um único campo de texto em Markdown.}
{Acrescentar a \texttt{Configuration} seis campos opcionais e independentes: \arq{discordUrl}, \arq{githubUrl}, \arq{youtubeUrl}, \arq{instagramUrl}, \arq{linkedinUrl} e \arq{twitchUrl}. Apenas URLs absolutas com protocolo \texttt{http} ou \texttt{https} são aceitas (RN17); campo vazio significa ícone omitido. O texto de direitos do rodapé é gerado a partir do nome do site e do ano corrente, sem campo próprio nesta versão.}
{Positivas: formulário simples e validável por campo; ícones padronizados; ausência de Markdown arbitrário no rodapé, o que reduz a superfície de ataque. Negativas: a lista de redes é fixa, e incluir outra rede exige migração e ajuste da interface.}
{Campo único em Markdown (recomendação inicial); campo JSON com a lista de links; tabela de links extensível (descartada por excesso de estrutura para a versão inicial).}

\adrentry{ADR-014}{Driver adapter better-sqlite3 e conformidade com o Prisma 7}{Aceita (proposta pelo agente Dados no pedido M-001 e aprovada pelo Arquiteto em 30/09/2026; resolve \decisao{D-09}, itens 1 a 5).}
{O Prisma 7 exige configuração da CLI em \arq{prisma.config.ts}, o gerador \texttt{prisma-client} e um \en{driver adapter} para o SQLite; a configuração inicial do repositório seguia práticas do Prisma 6 e declarava o pacote \arq{sqlite3}, que não é adaptador do Prisma. Com a implantação em VPS (\adr{ADR-011}), não há necessidade de banco remoto.}
{Adotar o adaptador \arq{@prisma/adapter-better-sqlite3} com o \arq{better-sqlite3} (série 12, em \texttt{dependencies}); renomear a configuração para \arq{prisma.config.ts}; usar o gerador \texttt{prisma-client} com saída em \arq{src/generated/prisma} (não versionada); retirar a URL do \texttt{datasource} do esquema; declarar o \arq{dotenv} em \texttt{devDependencies} para a CLI; mover \arq{@prisma/client} para \texttt{dependencies}; remover o \arq{sqlite3}.}
{Positivas: acesso direto ao arquivo SQLite, sem rede, coerente com a implantação em contêiner. Negativas: \arq{better-sqlite3} é módulo nativo (o \en{script} de instalação precisa ser permitido em \texttt{allowScripts}, e o \en{build} da imagem precisa de ferramentas de compilação ou binários pré-compilados); o cliente gerado é escrito em TypeScript dentro de um projeto JavaScript, e sua compilação pelo Next.js deve ser validada ao implementar o \arq{db.js}; \arq{prisma} e \arq{dotenv} permanecem em \texttt{devDependencies}, o que afeta a aplicação de migrações na imagem de produção (módulo M13).}
{\arq{@prisma/adapter-libsql} com \arq{@libsql/client} (banco remoto ou \en{edge}; mantido como via futura, não adotada pela \adr{ADR-011}); o pacote \arq{sqlite3} (não é adaptador do Prisma).}

\endgroup


%% ===================================================================================
%% Apêndice F -- Decisões em Aberto
%% Registro vivo das perguntas que ainda não têm resposta. Quando uma decisão for tomada,
%% o Arquiteto cria a ADR correspondente (Apêndice E) e marca aqui "Decidida -> ADR-0xx".
%% Identificadores D-xx são referenciados com \decisao{D-xx} nos demais capítulos.
%% "Responsável" é quem deve propor a resposta no mural; o Arquiteto decide.
%% ===================================================================================
\chapter{Decisões em Aberto}\label{ap:decisoes}

Este apêndice registra as decisões de projeto ainda não tomadas, com as opções conhecidas, a recomendação inicial do Arquiteto e o responsável pela proposta. Decisões que dependem do autor ou do orientador (marcadas ``Autor'') não podem ser tomadas pelos agentes. Quando uma decisão é tomada, ela passa a constar como ADR no Apêndice~\ref{ap:adr}.

\begin{longtable}{@{}L{1.1cm}L{6.3cm}L{4.5cm}L{1.9cm}@{}}
\caption{Registro de decisões em aberto}\label{tab:decisoes}\\
\toprule
\textbf{ID} & \textbf{Questão e opções} & \textbf{Recomendação inicial} & \textbf{Responsável}\\
\midrule
\endfirsthead
\toprule
\textbf{ID} & \textbf{Questão e opções} & \textbf{Recomendação inicial} & \textbf{Responsável}\\
\midrule
\endhead
\bottomrule
\endfoot
\decisaotarget{D-01} & \textbf{Nome oficial.} Os comentários de \arq{doc/ref/er.dbml} e de \arq{prisma/schema.prisma} usam ``Vitreo CMS''; o repositório, a governança e o relatório de IC usam ``Vitrine''. & Adotar ``Vitrine'' e corrigir os comentários na próxima alteração do esquema. & Autor\\[1pt]
\decisaotarget{D-02} & \textbf{Licença de código aberto.} Opções: MIT, Apache-2.0, GPL-3.0, AGPL-3.0, entre outras. & Preferir licença permissiva para facilitar a adoção; decidir ouvindo orientador e a política da instituição. & Autor\\[1pt]
\decisaotarget{D-03} & \textbf{Política de cadastro e \en{bootstrap} do administrador.} Opções: \en{script} de \en{seed} com credenciais informadas na instalação; comando de linha de comando interativo; página de ``primeiro acesso'' ativa enquanto não há usuários. & Comando ou \en{script} explícito na instalação e cadastro público desabilitado; evitar página aberta de primeiro acesso (risco de sequestro da instalação). & Core\\[1pt]
\decisaotarget{D-04} & \textbf{Alvo de implantação e \en{driver} de banco.} Servidor ou contêiner com SQLite em volume persistente; plataforma \emph{serverless} (por exemplo, Vercel) com banco remoto compatível (libSQL/Turso). & Documentar o primeiro como suportado e o segundo como experimental, dependente do \en{driver adapter}. & Arquiteto e Autor\\[1pt]
\decisaotarget{D-05} & \textbf{Unicidade de \en{slug} e \en{slugs} reservados.} Índice único composto (\texttt{parentId}, \texttt{slug}) mais índice parcial para a raiz; ou apenas validação na aplicação. & Ambos: validação na aplicação e índice no banco; reservar \texttt{admin}, \texttt{api}, \texttt{login} e \arq{_next}. & Dados\\[1pt]
\decisaotarget{D-06} & \textbf{Resolução de caminhos.} Calculada em memória a partir da árvore; ou caminho materializado em cada nó. & Calculada (Código~\ref{lst:algoritmo-arvore}); reavaliar se o volume crescer. & Dados\\[1pt]
\decisaotarget{D-07} & \textbf{Realce de sintaxe.} Opções: \arq{rehype-highlight}, \arq{rehype-prism}, Shiki (por exemplo, via \arq{rehype-pretty-code}). & Preferir solução que realce no servidor, sem JavaScript adicional no cliente; avaliar temas e tamanho. & UI\\[1pt]
\decisaotarget{D-08} & \textbf{Markdown estendido.} Adicionar \arq{remark-gfm} (pacote ainda não instalado) para tabelas, tachado e listas de tarefas. & Sim (\req{RF11}). & UI\\[1pt]
\decisaotarget{D-09} & \textbf{Conformidade com o Prisma 7.} Seis divergências listadas na Tabela~\ref{tab:prisma7} (arquivo de configuração, gerador, URL do banco, \en{driver adapter}, classificação de \arq{@prisma/client}, caminho de \arq{DATABASE_URL}). & Resolver na primeira ramificação de Dados, antes da primeira migração. & Dados\\[1pt]
\decisaotarget{D-10} & \textbf{Logotipo e \en{favicon}.} URL externa; arquivo em \arq{public/}; envio (\en{upload}) com armazenamento local ou em objeto. & Nesta versão, URL externa ou arquivo em \arq{public/}, sem \en{upload}. & Core e UI\\[1pt]
\decisaotarget{D-11} & \textbf{Validação de entrada.} Zod (presente como dependência transitiva do better-auth) ou validação manual. & Zod como dependência direta. & Core\\[1pt]
\decisaotarget{D-12} & \textbf{Ferramentas de teste.} Vitest, Playwright e Testing Library; ou Jest. & Vitest e Playwright (Tabela~\ref{tab:tipos-teste}). & Arquiteto\\[1pt]
\decisaotarget{D-13} & \textbf{Papéis e permissões.} Administrador único; ou papéis com o \en{plugin} \texttt{admin} do better-auth. & Administrador único nesta versão; manter o esquema compatível com evolução futura. & Core\\[1pt]
\decisaotarget{D-14} & \textbf{Estratégia de ordenação.} Inteiros contíguos com reindexação em transação; inteiros espaçados; chaves fracionárias. & Inteiros contíguos com reindexação (árvore pequena). & Dados\\[1pt]
\decisaotarget{D-15} & \textbf{Estados de publicação e versionamento.} Publicação imediata; campo de estado (rascunho e publicado); tabela de versões. & Publicação imediata e rascunho apenas local (\req{RF15}); versionamento fora do escopo. & Arquiteto e Autor\\[1pt]
\decisaotarget{D-16} & \textbf{Página inicial e coleções.} Página inicial com lista das coleções raiz; redirecionamento ao primeiro documento. & Início com a lista de coleções e documentos da raiz; coleção redireciona ao primeiro documento (RN14). & UI e Core\\[1pt]
\decisaotarget{D-17} & \textbf{Modo escuro \emph{versus} paleta personalizada.} O \en{scaffold} define cores para \texttt{prefers-color-scheme: dark}. & Sem modo escuro automático nesta versão: respeitar somente a paleta escolhida. & UI\\[1pt]
\decisaotarget{D-18} & \textbf{Dados do TCC.} Título definitivo, orientador, coorientador, banca e datas. & Confirmar com o orientador. & Autor\\[1pt]
\decisaotarget{D-19} & \textbf{Links do rodapé.} O esquema não prevê onde guardá-los. Opções: campo JSON de links; campo de texto em Markdown para o rodapé; arquivo estático. & Campo de texto Markdown para o rodapé, em \texttt{Configuration}, sem nova tabela. & Dados e UI\\[1pt]
\end{longtable}
\fonte{Elaborado pelos autores.}

%% ===================================================================================
%% Apêndice G -- Mapa de Pendências (GERADO AUTOMATICAMENTE)
%% A tabela é montada a partir do .aux: cada \vitrinefigura sem imagem, cada \vitrinecodigo
%% sem arquivo e cada \vitrinependencia aparecem aqui. Exige DUAS compilações (latexmk faz).
%% Nada neste arquivo deve ser editado à mão, exceto o texto introdutório.
%% ===================================================================================
\chapter{Mapa de Pendências da Documentação}\label{ap:pendencias}

Este apêndice é gerado automaticamente a cada compilação em modo rascunho. Lista (i) as \textbf{figuras} cuja imagem ainda não existe em \arq{doc/tex/figuras/}, com o caminho exato onde o arquivo deve ser salvo; (ii) os \textbf{trechos de código} do repositório ainda inexistentes; e (iii) as \textbf{pendências de texto} endereçadas a um agente. Quando o arquivo de imagem ou de código é adicionado, o item desaparece na compilação seguinte; quando a pendência de texto é resolvida, o Arquiteto remove o aviso do capítulo.

\textbf{Como localizar as pendências de um agente:} no terminal, dentro de \arq{doc/tex}, execute \texttt{make pendentes AGENTE=UI} (ou \texttt{Core}, \texttt{Dados}, \texttt{Arquiteto}). Na versão final, a opção \texttt{final} do pacote \texttt{vitrine} oculta os avisos e transforma figuras e códigos ausentes em erro de compilação.

\vitrineimprimirpendencias

%% ===================================================================================
%% Apêndice H -- Glossário
%% Termos do projeto e da metodologia. Ordem alfabética. Dono: Arquiteto/Redator.
%% ===================================================================================
\chapter{Glossário}\label{ap:glossario}

\begin{description}[style=nextline, leftmargin=0.6cm, itemsep=4pt]
  \item[Administrador] Ator humano que se autentica para gerenciar o conteúdo e a identidade visual do site. Na versão atual, é o único perfil autenticado.
  \item[Agente] Participante especializado do desenvolvimento (Arquiteto, Dados, Core ou UI), em geral uma ferramenta de inteligência artificial operada sob supervisão humana.
  \item[App Router] Sistema de roteamento do Next.js baseado na estrutura do diretório \arq{src/app}.
  \item[Arquiteto] Agente que governa a documentação, revisa pedidos de integração e executa os \en{merges} na \texttt{main}; também chamado Redator Técnico e \en{Merge Master}.
  \item[ADR] \en{Architecture Decision Record}: registro curto de uma decisão arquitetural com contexto, decisão e consequências.
  \item[better-auth] Biblioteca de autenticação usada pelo Vitrine para login por e-mail e senha e gestão de sessões.
  \item[Coleção] Conteúdo do tipo \texttt{COLLECTION}, que funciona como pasta e pode ter documentos e outras coleções como filhos.
  \item[Conteúdo] Nó da árvore do site: documento ou coleção.
  \item[Core] Agente responsável pela lógica de servidor: API, autenticação, roteamento e revalidação.
  \item[Dados] Agente responsável pelo esquema, pelas migrações e pelo acesso ao banco.
  \item[DBML] \en{Database Markup Language}, usada para descrever o modelo de dados como texto.
  \item[Documento] Conteúdo do tipo \texttt{DOCUMENT}: uma página com corpo em Markdown; não tem filhos.
  \item[Documentation-Driven Development] Processo em que a documentação é escrita antes do código e é a fonte da verdade que o código deve atender.
  \item[Fonte da verdade] Artefato que prevalece em caso de conflito sobre determinada informação (Tabela~\ref{tab:fontes-verdade}).
  \item[ISR] \en{Incremental Static Regeneration}: estratégia do Next.js que serve páginas estáticas e as regenera de forma incremental.
  \item[KaTeX] Biblioteca de composição de fórmulas matemáticas para a web.
  \item[Merge] Integração de uma ramificação na \texttt{main}.
  \item[Monólito modular] Aplicação implantada como uma unidade, com módulos internos bem delimitados.
  \item[Mural] Arquivo \arq{MURAL_DE_AGENTES.md}, canal de comunicação e de pedidos de integração, editado somente na \texttt{main}.
  \item[Placeholder] Marcador de figura ou de código ausente, que indica o arquivo a ser criado e é substituído automaticamente quando o arquivo existe.
  \item[Proxy] Arquivo do Next.js 16 (antigo \en{Middleware}) que executa lógica antes da conclusão de uma requisição.
  \item[Revalidação] Invalidação de páginas em cache para que sejam regeneradas na próxima requisição.
  \item[Route Handler] Arquivo \arq{route.js} do App Router que implementa um \en{endpoint} HTTP.
  \item[Server Component] Componente React executado apenas no servidor.
  \item[Slug] Trecho legível da URL que identifica um conteúdo entre seus irmãos.
  \item[SQLite] Banco de dados SQL embutido, em arquivo único.
  \item[Tema] Conjunto de identidade (nome, logotipo) e cores do site, aplicado por variáveis CSS.
  \item[UI] Agente responsável pela interface: componentes, estilos e editor.
  \item[WAL] \en{Write-Ahead Logging}: modo do SQLite que permite leituras simultâneas a uma escrita.
\end{description}

\end{apendicesenv}

\end{document}
